\documentclass[11pt,a4paper]{ctexart}

% ── Packages ──
\usepackage[top=2.5cm, bottom=2.5cm, left=2.8cm, right=2.8cm]{geometry}
\usepackage{amsmath,amssymb,amsthm}
\usepackage{bm}
\usepackage{graphicx}
\usepackage{booktabs}
\usepackage{longtable}
\usepackage[colorlinks=true,linkcolor=blue,citecolor=blue,urlcolor=blue]{hyperref}
\usepackage{enumitem}
\usepackage{fancyhdr}
\usepackage{titlesec}
\usepackage{xcolor}
\usepackage{float}
\usepackage{tikz}

\pagestyle{fancy}
\fancyhf{}
\fancyhead[L]{\small HMAE 数学与物理逻辑基础}
\fancyhead[R]{\small MAPFM Medical AI}
\fancyfoot[C]{\thepage}
\renewcommand{\headrulewidth}{0.4pt}

% ── Theorems ──
\newtheorem{definition}{定义}[section]
\newtheorem{theorem}{定理}[section]
\newtheorem{lemma}{引理}[section]
\newtheorem{axiom}{公理}[section]
\theoremstyle{remark}
\newtheorem{remark}{注}[section]

% ── Math operators ──
\DeclareMathOperator{\softmax}{softmax}
\DeclareMathOperator{\sigmoid}{sigmoid}
\DeclareMathOperator{\clip}{clip}
\DeclareMathOperator{\simhash}{SimHash}
\DeclareMathOperator{\Ltwo}{L2}
\DeclareMathOperator{\TFIDF}{TF-IDF}
\DeclareMathOperator{\E}{\mathbb{E}}
\DeclareMathOperator{\Var}{Var}
\DeclareMathOperator{\Ham}{Hamming}
\DeclareMathOperator{\KL}{KL}

\begin{document}

% ============================================================
\title{\textbf{HMAE 异构多智能体医疗AI生态系统\\
底层数学与物理逻辑推导\\
—— 附公理体系推演}}
\author{MAPFM Medical AI Research Framework\\[4pt]
\small Heterogeneous Multi-Agent Ecosystem v2.0\\[4pt]
\small 丰富讲解版（含公理化推导）}
\date{\today}
\maketitle

\begin{center}
\small
\textbf{摘要} 本文从第一性原理出发，系统推导 HMAE 异构多智能体医疗 AI 生态系统的完整数学与物理基础。
全文分为五个部分：(I) 十大智能体的核心数学公式与物理类比；(II) 支撑基础设施的数学推导；
(III) 跨模块耦合与系统动力学；(IV) 输出评价体系；(V) 从数学结构中提炼的公理体系。
第五部分是本文的核心创新——我们从工程实现中逆向归纳出8条公理，构成本系统的公理化基础。
\end{center}

\tableofcontents
\newpage

% ============================================================
% PART I: 智能体部分
% ============================================================
\part{智能体部分 —— Agent Mathematical Foundations}

\vspace{6pt}
\begin{center}
\small
\textit{Part I 涵盖 HMAE 生态系统的十大核心智能体。每个智能体的数学推导都对应一个物理系统的类比，}\\
\textit{这一"数学-物理"双重对应关系是理解整个系统行为的关键。}
\end{center}
\vspace{12pt}

% ============================================================
\section{感知智能体 (Perception Agent)}

\subsection{设计动机与物理直觉}

感知智能体是多模态输入的"统一入口"。它的核心任务是将异质的临床数据（文本、图像）映射到
一个共同的低维语义流形 $\mathcal{M} \subset \mathbb{S}^{d-1}$（单位超球面）上。
这一映射在物理上等价于\textbf{降维投影测量（dimensionality-reducing projective measurement）}：
我们从一个极高维的原始信号空间（文本 token 空间 $\sim 10^4$ 维 + 图像像素空间 $\sim 10^5$ 维）
投影到 $d=128$ 维的语义流形，保留下游任务最需要的信息分量。

\subsection{文本编码}

输入临床文本 $X_{\text{txt}}$ 首先通过 MiniLM (all-MiniLM-L6-v2) 编码器映射为语义向量：

\begin{equation}
\mathbf{v}_{\text{txt}} = \text{MiniLM}(X_{\text{txt}}) \in \mathbb{R}^{d_{\text{MiniLM}}}
\end{equation}

其中 $d_{\text{MiniLM}} = 384$（MiniLM 原生维度），随后通过分块均值降维至项目维度 $d = 128$：

\begin{equation}
\mathbf{v}_{\text{txt}}^{(d)} = \text{Resize}(\mathbf{v}_{\text{txt}}, d) = \Ltwo\left(\Big[\operatorname{mean}(\mathbf{v}_{\text{chunk}_k})\Big]_{k=1}^{d}\right)
\end{equation}

\textbf{物理注解：} $\Ltwo$ 归一化将所有向量约束在单位球面 $S^{d-1}$ 上。这一约束的深层含义是：
在语义空间中，\textbf{方向比幅值更重要}。两个临床文本的相似度仅由它们之间的夹角决定，
文本长度（信息量大小）不影响语义比较。这在物理上对应量子态的描述——
态的物理内容由归一化后的态矢量方向完全确定，模长不承载物理信息。

\subsection{图像编码 (ViT Patching)}

对输入图像 $X_{\text{img}}$，采用 ViT 风格分块编码：

\begin{equation}
\mathbf{P} = \left[\mathbf{p}_1, \mathbf{p}_2, \mathbf{p}_3, \mathbf{p}_4\right]^{\top} \in \mathbb{R}^{4 \times d}
\end{equation}

其中每个 $\mathbf{p}_i$ 来自图像块的文本代理编码：$\mathbf{p}_i = \text{MiniLM}\left(\text{TextProxy}(\text{flatten}(X_{\text{img}})_{[i\cdot s:(i+1)\cdot s]})\right)$。

图像上下文向量：
\begin{equation}
\mathbf{v}_{\text{img}} = \Ltwo\left(\frac{1}{4}\sum_{i=1}^{4} \mathbf{p}_i\right)
\end{equation}

\textbf{设计原理：} 使用 TextProxy 将图像像素映射为文本描述再通过 MiniLM 编码，
这一"迂回"策略的深层原因是——MiniLM 是在自然语言语料上预训练的，直接输入像素值会产生
分布外 (out-of-distribution) 激活。TextProxy 将像素统计量翻译为自然语言描述（如
"region 1: high intensity, irregular texture"），使图像信息进入 MiniLM 的分布内区域，
从而获得有意义的语义编码。这在物理上等价于\textbf{量子测量中的经典仪器读数}：
量子系统（图像）的信息必须通过经典仪器（TextProxy）转换为经典语言才能被经典处理器（MiniLM）理解。

\subsection{交叉注意力融合 (Cross-Attention Fusion)}

融合层通过文本向量 $\mathbf{q} = \mathbf{v}_{\text{txt}}$（作为查询）与图像块 $\mathbf{P}$ 进行注意力计算：

\begin{equation}
\mathbf{a} = \softmax\left(\frac{\mathbf{P}\,\mathbf{q}^{\top}}{\sqrt{d}}\right) \in \mathbb{R}^{4}
\end{equation}

视觉上下文：
\begin{equation}
\mathbf{c}_{\text{vis}} = \sum_{i=1}^{4} a_i \cdot \mathbf{p}_i
\end{equation}

最终融合上下文向量通过串联 + 线性投影：

\begin{equation}
\mathbf{v}_{\text{fused}} = \Ltwo\left(\begin{bmatrix} \mathbf{v}_{\text{txt}} & \mathbf{c}_{\text{vis}} \end{bmatrix} \cdot \mathbf{W}_{\text{fusion}}\right)
\label{eq:fusion}
\end{equation}

其中 $\mathbf{W}_{\text{fusion}} \in \mathbb{R}^{2d \times d}$，初始化 $w_{ij} \sim \mathcal{N}(0, 1/\sqrt{d})$。

\begin{remark}[交叉注意力的物理意义]
式 (\ref{eq:fusion}) 的物理意义极为深刻：文本语义与视觉信号在高维流形上的交集投影。
\textbf{当图像与文本描述一致时}（如文本描述"肺部X光显示右下叶浸润"，图像确实显示该特征），
注意力权重 $a_i$ 会高度集中在包含该特征的图像块上，使得 $\mathbf{c}_{\text{vis}}$ 与 $\mathbf{v}_{\text{txt}}$ 方向对齐。
此时 $\mathbf{v}_{\text{fused}}$ 的范数接近1，表示"高置信度融合"。

\textbf{当图文不一致时}（如文本描述肺炎但图像显示正常肺），注意力权重分散，
$\mathbf{c}_{\text{vis}}$ 与 $\mathbf{v}_{\text{txt}}$ 方向偏离。L2 约束使融合向量仍位于单位球面上，
但其方向偏离了文本与图像的共同方向，系统通过向量夹角的增大自然地表达了"不确定性"。

这一机制在数学上等价于\textbf{量子纠缠态的投影测量}：
两个子系统（文本、图像）的联合态被投影到一个共同的本征基上，
融合结果取决于两个子系统是否"相容"（即是否处于同一本征态）。
\end{remark}

% ============================================================
\section{检索增强智能体 (RAA — Retrieval-Augmented Agent)}

\subsection{从感知到检索：为什么需要 RAA}

感知智能体（§1）将异构的临床输入映射到统一的语义流形 $S^{d-1}$，但这只是完成了"理解"的第一步。临床诊断还需要\textbf{外部知识}——医学文献、临床指南、病例报告、药物相互作用数据库——这些信息不可能被压缩进一个 128 维的向量中。RAA 的角色是\textbf{以感知智能体的输出为查询，在外部知识库中检索最相关的医学知识，弥合"当前病例"与"已建立的医学知识"之间的信息鸿沟}。

\subsection{双通道检索的互补性原理}

RAA 同时运行稀疏检索（TF-IDF）和稠密检索（FAISS）两个通道。
这并非工程上的冗余设计，而是基于深刻的\textbf{信号处理对偶性}：

\begin{itemize}[nosep]
\item \textbf{稀疏检索（词频空间）}：在词项的离散基上操作，擅长精确词匹配，等价于信号处理中的\textbf{时域分析}——定位到具体的关键词。
\item \textbf{稠密检索（语义空间）}：在连续嵌入基上操作，擅长语义近似匹配，等价于信号处理中的\textbf{频域分析}——捕获整体的语义"频谱"。
\end{itemize}

二者的融合可以用\textbf{Gabor 变换的离散类比}来理解：TF-IDF 提供了"哪个词在哪里出现"的局部信息，
FAISS 提供了"整体语义是什么"的全局信息。同时使用两者，类似于在时频平面上同时具有最优分辨率
（虽然受制于不确定性原理的 trade-off，见第五部分公理2）。

\subsection{稀疏检索：TF-IDF 余弦相似度}

给定查询 $q$ 和文档集 $\mathcal{D}$，构建 TF-IDF 矩阵 $\mathbf{M} \in \mathbb{R}^{|\mathcal{D}| \times k_{\max}}$（$k_{\max}=2400$）：

\begin{equation}
w_{ij} = \text{tf}(t_j, d_i) \cdot \log\frac{|\mathcal{D}|}{|\{d \in \mathcal{D}: t_j \in d\}|}
\end{equation}

采用次线性 TF 缩放 $\text{tf}_{\text{sub}}(t, d) = 1 + \log(1 + \log(\text{tf}(t, d)))$ 和 $(1,3)$-gram 范围。

\textbf{次线性缩放的物理原因：} 一个词在文档中出现100次并不比出现10次传递10倍的信息量。
次线性缩放符合\textbf{信息论中的递减边际信息原理}——每增加一次出现的边际信息增益随出现次数递减。
双重对数 $1 + \log(1 + \log(\cdot))$ 确保了即使对于极高频词，权重增长也被严格压制。

查询向量 $\mathbf{q}_{\text{tfidf}} = \text{Vectorizer}(q)$，相似度：

\begin{equation}
s_{\text{sparse}}(d_i) = \cos(\mathbf{q}_{\text{tfidf}}, \mathbf{m}_i) = \frac{\mathbf{q}_{\text{tfidf}} \cdot \mathbf{m}_i}{\|\mathbf{q}_{\text{tfidf}}\|_2 \cdot \|\mathbf{m}_i\|_2}
\end{equation}

加权相似度：$s'_{\text{sparse}}(d_i) = s_{\text{sparse}}(d_i) \cdot w_{\text{doc}}(d_i)$，其中 $w_{\text{doc}}(d_i)$ 为文档动态权重。

\subsection{稠密检索：FAISS 内积搜索}

语义向量通过 MiniLM 编码后经 FAISS 索引（IndexFlatIP）进行最大内积搜索（MIPS）：

\begin{equation}
\mathbf{e}_q = \text{MiniLM}(q) \in \mathbb{R}^{d}, \quad
\mathbf{M}_{\text{dense}} \in \mathbb{R}^{|\mathcal{D}| \times d}
\end{equation}

\begin{equation}
s_{\text{dense}}(d_i) = \frac{\langle \mathbf{e}_q, \mathbf{m}_i^{\text{dense}} \rangle + 1}{2} \in [0, 1]
\end{equation}

\begin{remark}[内积搜索的几何本质]
在 $\Ltwo$ 归一化后（即所有向量在单位球面上），内积 $\langle \mathbf{a}, \mathbf{b} \rangle$ 严格等于 $\cos(\theta_{\mathbf{a},\mathbf{b}})$。
因此 FAISS MIPS 实际上是在球面上寻找与查询向量夹角最小的文档——等价于球面上的最近邻搜索。
归一化至 $[0,1]$ 区间的变换 $(\cos\theta + 1)/2$ 将角度 $[0,\pi]$ 线性映射到 $[0,1]$，
保持了顺序关系（单调递减映射的逆）。
\end{remark}

\subsection{查询不确定性量化}

RAA 通过三组分解合不确定性估计查询难度：

\begin{align}
H(P) &= -\sum_{i=1}^{5} p_i \log p_i / \log(5) \quad \text{(归一化熵)} \\
G(P) &= 1 - \clip_{[0,1]}(p_1 - p_2) \quad \text{(Top-2 间隙)} \\
O(q) &= \mathbb{I}[|\text{tokens}(q)| \leq 2] \quad \text{(OOV 代理)}
\end{align}

\begin{equation}
U(q) = \alpha_H \cdot H(P) + \alpha_G \cdot G(P) + \alpha_O \cdot O(q)
\label{eq:uncertainty}
\end{equation}

其中 $\alpha_H = 0.58$, $\alpha_G = 0.32$, $\alpha_O = 0.10$，概率 $P = \softmax(\mathbf{s}, T=0.35)$ 来自前5个检索结果的相似度得分。

当 $U(q) > \tau_U = 0.58$ 时触发深层检索（扩大 $k$ 至10 + 多样化策略）。

\begin{remark}[不确定性量化的三重分解]
$H(P)$（归一化熵）度量检索得分分布的"平坦度"——如果前5个文档得分几乎相同（高熵），
说明检索系统无法区分它们，查询本身缺乏区分性信息。这对应\textbf{热力学中的最大熵状态}。

$G(P)$（Top-2 间隙）度量最佳文档与次佳文档之间的分离度——如果间隙很小，
系统在"两个都很像"的歧义状态中。这对应\textbf{量子力学中的近简并能级}。

$O(q)$（OOV 代理）标记极短查询（$\leq 2$ tokens），这种查询几乎一定是
OOV (out-of-vocabulary) 或信息不足的。这对应\textbf{信号处理中的欠采样}。
\end{remark}

\subsection{重排序 (Rerank-RAG)}

\begin{equation}
s_{\text{rerank}}(d) = \omega_s \cdot s(d) + \omega_o \cdot \frac{|\mathcal{T}_q \cap \mathcal{T}_d|}{|\mathcal{T}_q|} + \omega_r \cdot e^{-\Delta t / 540} + b_{\text{auth}} + b_{\text{med}}
\label{eq:rerank}
\end{equation}

其中 $\omega_s=0.62$, $\omega_o=0.22$, $\omega_r=0.10$, $b_{\text{auth}}=0.10$, $b_{\text{med}} \in [0, 0.12]$。
$\Delta t$ 以天为单位，540天对应约1.5年，时效性半衰期约为 $540 \cdot \ln 2 \approx 374$ 天。

\begin{remark}[信息检索的哈密顿量类比]
式 (\ref{eq:rerank}) 可视为信息检索中的\textbf{哈密顿量}（总能量函数）：
\begin{itemize}[nosep]
\item $\omega_s$ 项为\textbf{势能}（语义匹配——系统的"位置"在语义空间中的稳定程度）
\item $\omega_o$ 项为\textbf{动能}（词项交叠——系统的"运动"，即查询与文档之间的信息流动）
\item $\omega_r$ 项为\textbf{耗散}（信息时效性衰减——热力学第二定律在信息空间中的表达：信息随时间熵增）
\item $b_{\text{auth}}$ 为\textbf{外场偏置}（外部权威源施加的"磁场"，使检索偏好可信来源）
\end{itemize}
这一类比不仅仅是比喻：如果我们定义检索空间上的作用量 $S = \int (T - V) dt$，
最小作用量原理自然导致最优检索排序。
\end{remark}

% ============================================================
\section{决策智能体 (DMA — Decision-Making Agent)}

\subsection{从检索到决策：信息汇聚的自然终点}

感知智能体（§1）提供了对当前病例的"感知"（$\mathbf{v}_{\text{fused}} \in S^{d-1}$），RAA（§2）提供了与当前病例相关的外部医学知识（检索到的文档集 $\mathcal{K}_s$）。DMA 是这两股信息流的\textbf{自然汇聚点}——它将内部感知（病例本身是什么样）和外部知识（类似病例的医学文献怎么说）融合为最终的诊断决策。从信息处理的角度，DMA 实现了一个\textbf{贝叶斯决策规则}：以内部感知为似然函数，以外部知识为先验，输出后验概率分布 $P(c \mid q, \mathcal{K}_s)$。

\subsection{决策理论的几何解释}

DMA 是整个生态系统的核心。从决策理论的视角，DMA 实现了一个\textbf{贝叶斯决策规则}在
多源证据下的具体实例化。每个证据源（质心相似度、检索得分、词汇匹配、上下文信号、时序风险）
构成一个似然函数 $P(\text{evidence}_k \mid c)$，DMA 的融合得分 $s_c^{\text{raw}}$ 恰好是
这些对数似然的加权和（在 softmax 归一化之前）。

\subsection{基于质心的 TF-IDF 分类}

训练阶段，对每类 $c \in \mathcal{C}$ 计算质心：

\begin{equation}
\mathbf{c}_c = \Ltwo\left(\frac{1}{|\mathcal{D}_c|}\sum_{d \in \mathcal{D}_c} \text{TF-IDF}(d)\right) \in \mathbb{R}^{k_{\max}}
\end{equation}

推理阶段，查询 $q$ 对各质心的相似度：
\begin{equation}
s_c^{\text{centroid}}(q) = \langle \Ltwo(\text{TF-IDF}(q)), \mathbf{c}_c \rangle
\end{equation}

\textbf{几何解释：} 每个疾病类在 TF-IDF 空间中由一个单位向量（质心）代表。
分类过程等价于在 $k_{\max}$ 维单位球面上寻找与查询向量夹角最小的质心。
这构成一个\textbf{Voronoi 划分}——决策边界由相邻质心的中垂超平面确定。

\subsection{多源得分融合}

DMA 的原始得分由五部分组成：

\begin{equation}
s_c^{\text{raw}} = \underbrace{\lambda_c s_c^{\text{centroid}}}_{\text{质心}} + \underbrace{\lambda_r s_c^{\text{retrieval}}}_{\text{检索}} + \underbrace{\lambda_l s_c^{\text{lexical}}}_{\text{词汇}} + \underbrace{\lambda_k \cdot (1 - H(\mathbf{a}))}_{\text{上下文}} + \underbrace{\lambda_t \cdot r_{\text{TFA}}}_{\text{时序风险}} + b_c
\label{eq:dma_raw}
\end{equation}

其中权重系数：$\lambda_c=0.50$, $\lambda_r=0.34$, $\lambda_l=0.12$, $\lambda_k=0.02$, $\lambda_t=0.02$。
$H(\mathbf{a}) = -\sum_i a_i \log a_i$ 为交叉注意力权重向量 $\mathbf{a}$（式5）的熵，该项反映图文融合质量——
当图文一致时注意力集中（$H$ 低，$1-H$ 高），不一致时注意力分散（$H$ 高，$1-H$ 低）。

\begin{remark}[权重分配的信息论解释]
权重 $\boldsymbol{\lambda} = [0.50, 0.34, 0.12, 0.02, 0.02]$ 并非随意设定。
它们近似满足\textbf{信息衰减的幂律}：质心（最可靠的结构化知识）占主导权重 0.50；
检索（外部证据）次之，占 0.34；词汇（表面匹配）再次之，占 0.12。
这种分配在数学上近似于以 $r \approx 0.68$ 为公比的几何级数衰减，
反映了从"结构化先验"到"弱信号"的信息可靠性递减链。
\end{remark}

检索得分归一化：$s_c^{\text{retrieval}} = \sum_{d: \text{area}(d)=c} s(d) \cdot w_{\text{doc}}(d) / \sum_d s(d) \cdot w_{\text{doc}}(d)$。

词汇先验：$s_c^{\text{lexical}} = |\mathcal{T}_q \cap \mathcal{T}_c| / \max(|\mathcal{T}_c|, 1)$。

\subsection{易混淆疾病消歧}

对于查询标记与疾病标签的直接匹配、同义词匹配和部分重叠，应用离散奖励：

\begin{equation}
s_c^{\text{raw}} \leftarrow s_c^{\text{raw}} +
\begin{cases}
4.0 & \text{if } \mathcal{T}_c \subseteq \mathcal{T}_q \text{ or } |\text{Syn}(\mathcal{T}_c) \cap \mathcal{T}_q| \geq 2 \\
2.0 & \text{if } |\text{Syn}(\mathcal{T}_c) \cap \mathcal{T}_q| \geq 1 \\
2.0 & \text{if } |\mathcal{T}_c \cap \mathcal{T}_q| \geq 0.5 \cdot |\mathcal{T}_c|
\end{cases}
\end{equation}

\textbf{离散奖励的博弈论解释：} 这不是连续得分，而是一种\textbf{信号机制（signaling）}——
当查询中的词汇与某疾病的术语集中的词精确匹配（特别是同义词匹配）时，
这传递了一个几乎无噪声的"信号"，远远强于统计相似度。
+4.0 的加成相当于统计得分在因子 $\sim 1.5\times$ 级别上的偏移，
确保精确术语匹配在决策中产生决定性影响。

\subsection{Softmax 概率输出}

\begin{equation}
P(c|q) = \frac{\exp(s_c^{\text{raw}} / T)}{\sum_{c'}\exp(s_{c'}^{\text{raw}} / T)}, \quad T = 0.18
\end{equation}

\begin{remark}[低温极限的物理意义]
低温 $T=0.18$ 使 softmax 趋近于 argmax（低熵分布）。在统计力学中，
温度 $T \to 0$ 时系统冻结至基态——所有概率质量集中在最低能量状态。
在 DMA 中，低温度确保系统输出的概率分布高度尖锐（低熵），
\\ \textbf{当 $T \to 0$ 时：} $\softmax(\mathbf{s}/T) \to \text{one\_hot}(\arg\max \mathbf{s})$，
系统退化为确定性分类器。
\\ \textbf{当 $T \to \infty$ 时：} $\softmax(\mathbf{s}/T) \to \text{uniform}$，
系统退化为随机猜测。
\\ 物理上，$T=0.18$ 在"尖锐决策"和"保留不确定性度量"之间取得平衡——
系统有足够的置信度进行决策，同时保留了概率排序信息用于后续校准。
\end{remark}

\subsection{Platt 置信度校准}

原始置信度 $\bar{p} = \max_c P(c|q)$ 经 Platt scaling 校准：

\begin{equation}
p_{\text{cal}} = \frac{1}{1 + \exp(-(A \cdot \bar{p} + B))}
\end{equation}

其中 $(A, B)$ 通过 SGD 在滑动窗口（最近200对 $(\bar{p}, \text{is\_correct})$）上拟合逻辑回归得到：

\begin{equation}
a \leftarrow a - \eta \cdot \mathbb{E}[(p_{\text{cal}} - y) \cdot \bar{p}], \quad b \leftarrow b - \eta \cdot \mathbb{E}[p_{\text{cal}} - y]
\end{equation}

\begin{remark}[Platt 校准与最大熵原理]
Platt scaling 本质上是在\textbf{保序变换}的约束下，寻找使预测置信度与实际正确率最一致的
单调映射。这等价于在保持 ROC 曲线不变的条件下，最小化预测分布与经验分布之间的 KL 散度。
从统计力学的角度看，这是受保序约束的最大熵校准——
最终的校准曲线是玻尔兹曼分布的一维投影。
\end{remark}

\subsection{Llama3:8b 推理（主路径）}

当 Ollama 服务可用时，DMA 通过 LLAMA3\_PROMPT\_TEMPLATE 构造提示词，调用本地 Llama3:8b 结构化 JSON 推理。
Ollama 失败时自动回退至上述质心 + 检索融合路径。

\textbf{双路径架构的深层逻辑：} Llama3 路径提供"深度推理"（链式思维、综合判断），
质心路径提供"快速模式匹配"（统计相似度）。这对应人类诊断的双过程理论 (dual-process theory)：
\begin{itemize}[nosep]
\item System 1（质心路径）：快速、直觉、基于模式识别（对应 Kahneman 的"快思考"）
\item System 2（Llama3 路径）：缓慢、分析、基于逻辑推理（对应 Kahneman 的"慢思考"）
\end{itemize}

% ============================================================
\section{时序预见智能体 (TFA — Temporal Foreseeing Agent)}

\subsection{从诊断到预测：决策的时空延伸}

DMA（§3）回答的问题是"患者现在可能患有什么疾病？"——这是一个\textbf{空间型}问题（在疾病类别空间中定位当前状态）。TFA 回答的是"患者的健康状况将如何随时间演化？"——这是一个\textbf{时间型}问题（在时间轴上预测未来状态）。两者互补形成完整的临床图景：诊断告诉我们"是什么"，预测告诉我们"会怎样"。TFA 的输入包括 DMA 的诊断结果 $\hat{c}$（作为风险预测的条件），这建立了从诊断空间到风险时间序列的因果桥梁。

\subsection{时间序列预测的物理框架}

TFA 实现了一个\textbf{物理信息系统中的状态预测器}。其核心思想是：
将临床时间序列视为一个受噪声扰动的动力系统的观测轨迹，
目标是推断系统的隐藏状态并预测未来演化。

\subsection{多尺度时序特征抽取}

TFA 从72点历史时间序列 $x_{1:T}$（$T=72$）中提取三类特征：

\textbf{TCN（时序卷积网络代理特征）：}
\begin{equation}
f_{\text{TCN}} = \frac{1}{|W|}\sum_{t \in W} \Delta x_t + 0.30 \cdot \sigma_W, \quad \Delta x_t = x_t - x_{t-1}, \quad W = x_{T-23:T}
\end{equation}

\textbf{EMA（指数加权移动平均代理特征）：}
\begin{equation}
h_t = \alpha \cdot x_t + (1 - \alpha) \cdot h_{t-1}, \quad \alpha = 0.18
\end{equation}
\begin{equation}
f_{\text{EMA}} = h_T
\end{equation}

\textbf{Transformer（自注意力代理特征）：}
\begin{equation}
f_{\text{Trans}} = \sum_{t=1}^{T} x_t \cdot \frac{\exp((\text{pos}_t + 0.2 \cdot x_t)/0.45)}{\sum_{j}\exp((\text{pos}_j + 0.2 \cdot x_j)/0.45)}
\end{equation}

其中 $\text{pos}_t = t/T \in [0, 1]$。

\begin{remark}[三种特征的时间尺度对应]
三类特征对应于三个物理时间尺度上的信息提取：
\begin{itemize}[nosep]
\item \textbf{TCN（短期，$W=24$ 窗口）}：局部梯度——最近的趋势和波动。等价于动力系统的\textbf{局部导数}估计（相空间中的速度）。
\item \textbf{EMA（全程，指数衰减记忆 $\alpha=0.18$，半衰期约3.8步）}：长期累积——整个历史的指数加权平均。等价于动力系统的\textbf{李雅普诺夫函数}（系统能量的单调趋势）。
\item \textbf{Transformer（全局，位置编码加权）}：全局注意力——所有时间点的位置加权融合。等价于\textbf{路径积分}——每一条历史路径根据其"重要性"（位置 + 信号强度）贡献于当前状态估计。
\end{itemize}
$\alpha=0.18$ 的半衰期为 $\ln(0.5)/\ln(0.82) \approx 3.8$ 个时间步——
意味着大约每4个时间步，历史信息的影响衰减一半。
\end{remark}

\subsection{门控融合}

\begin{equation}
g_k = \frac{\exp(|f_k| / T_g)}{\sum_j \exp(|f_j| / T_g)}, \quad T_g = 0.35
\end{equation}
\begin{equation}
f_{\text{fused}} = g_1 \cdot f_{\text{TCN}} + g_2 \cdot f_{\text{EMA}} + g_3 \cdot f_{\text{Trans}}
\end{equation}

\textbf{门控融合的量子测量类比：} 门控权重 $g_k$ 基于特征的\textbf{幅值}（而非方向）分配——
信号越强的特征获得越大的权重。这类似于量子测量中的
\textbf{波恩规则（Born rule）}：测量得到某结果的概率正比于对应本征态的概率幅的模方。
$T_g = 0.35$ 控制了门控的"选择性"——低温使门控更尖锐（赢家通吃），
高温使门控更均匀（平均融合）。

\subsection{多窗口风险预测}

\subsubsection{设计理念：三个时间尺度的临床含义}

临床风险不是一个单一的数字——同一个患者在"未来24小时"、"未来30天"和"未来12个月"面临的健康风险可能完全不同。多窗口预测的核心设计理念是：\textbf{将时间维度上的风险演化分解为三个因果耦合但数学独立的预测通道，每个通道使用最适合该时间尺度的特征和权重分配。}

这三个时间尺度在临床上有明确的对应：
\begin{itemize}[nosep]
\item \textbf{短期（24小时窗口）：} 对应急诊/ICU场景——急性感染、跌倒、休克、心律失常等迅速恶化的事件。短期预测必须对即时信号（当前的生命体征骤变、跌倒事件、急性感染指标）极度敏感，而对慢性病史相对不敏感。时间尺度：小时。
\item \textbf{中期（30天窗口）：} 对应住院/康复场景——术后恢复、药物调整后的稳定性、并发症的发展。中期预测需要在"近期趋势"和"长期基线"之间取得平衡。时间尺度：天到周。
\item \textbf{长期（12个月窗口）：} 对应慢病管理/预防场景——慢性病的进展、生活方式的累积效应、衰老相关的功能衰退。长期预测以慢性因素为主导，即时波动被平滑。时间尺度：月到年。
\end{itemize}

\subsubsection{输入变量的物理含义}

在展开公式之前，先定义所有输入变量及其临床含义：
\begin{itemize}[nosep]
\item $r_{\text{base}} \in [0,1]$ — TFA 基础预测模型输出的原始风险得分。这是 TFA 核心模型的输出——一个整合了 TCN、EMA、Transformer 三种特征的门控融合值经逻辑函数 $\sigma$ 映射后的 $[0,1]$ 概率。$r_{\text{base}}$ 可以理解为"忽略时间窗口差异的通用风险"——它为三个窗口提供了一个共同的基准线，各窗口在此基础上根据自身的时间特性进行调整。
\item $\delta_{\text{fall}} \in \{0, 0.15\}$ — 跌倒事件指示量。这是一个二值特征，但取值为 $\{0, 0.15\}$ 而非 $\{0, 1\}$。$0.15$ 的选择是在实验验证后确定的——如果设为 $1.0$（全量），跌倒事件会完全主导短期风险预测（因为 $1.0$ 经 $\sigma$ 后接近 $0.73$，远超其他信号的贡献）；$0.15$ 使跌倒事件对风险有显著但非压倒性的贡献（经 $\sigma$ 后约贡献 $0.04$ 的边际增量）。
\item $r_p = 0.08$ — 压疮增长率常数。控制压疮风险随卧床时间 $h_{\text{bed}}$ 增长的速率。$r_p = 0.08$ 意味着卧床 $t$ 小时后，压疮的\"风险潜力\"为 $1 - e^{-0.08 \cdot t}$。例如：$t=24$ 小时（1天）→ $1-e^{-1.92} \approx 0.85$（高风险）；$t=12$ 小时 → $1-e^{-0.96} \approx 0.62$（中等风险）；$t=6$ 小时 → $1-e^{-0.48} \approx 0.38$（低风险）。
\item $f_{\text{fall}} = \omega_p \cdot p_{\text{inst}} + \omega_a \cdot (1 - a_{\text{level}})$ — 跌倒风险特征。它综合了两个子因素：(1) $p_{\text{inst}}$ — 姿势不稳定性（postural instability），通过加速度计或运动传感器测量；(2) $a_{\text{level}}$ — 活动能力等级（activity level），$1 - a_{\text{level}}$ 表示活动能力越低（越不能自主移动），跌倒时受伤越严重。权重 $\omega_p = 0.6$ 偏向姿势不稳定性（起主要作用），$\omega_a = 0.4$ 偏向活动能力（起辅助作用）。
\item $\mu_{\text{acute}} \in \{1.0, 3.0\}$ — 急性病程乘数。当患者同时患有急性病（如急性肺炎、急性心肌梗死）时，$\mu_{\text{acute}} = 3.0$；否则 $\mu_{\text{acute}} = 1.0$。$3.0$ 的倍率意味着急性病将短期风险乘以 3——这在临床上是合理的：一个急性感染患者的 24 小时恶化概率可能是稳定期患者的 2--5 倍。
\item $\mu_{\text{chronic}} \in \{1.0, 2.0\}$ — 慢性病程乘数。当患者患有慢性病（如糖尿病、COPD、高血压）时，$\mu_{\text{chronic}} = 2.0$；否则 $\mu_{\text{chronic}} = 1.0$。慢性病的影响主要在长期窗口——因此在 $r_{\text{long}}$ 而非 $r_{\text{short}}$ 中应用此乘数。
\item $\sigma(z) = 1/(1+e^{-z})$ — 逻辑函数（logistic / sigmoid）。将无界实数信号 $z \in (-\infty, +\infty)$ 平滑地映射到 $(0, 1)$ 区间。在物理中这恰好是费米-狄拉克分布（描述费米子在各能级上的占据概率）。在风险预测中，它提供了将多种信号（正负均可）整合为一个概率值的最自然方式。
\end{itemize}

\subsubsection{三窗口公式的逐项推导}

有了上述准备，我们可以逐项解析三个窗口的公式：

\begin{align}
r_{\text{short}} &= \sigma\left(f_{\text{fused}} + b_{\text{auth}} + f_{\text{fall}} + \delta_{\text{fall}}\right) \cdot \mu_{\text{acute}} \\
r_{\text{mid}} &= \clip\left(0.72 \cdot r_{\text{base}} + 0.28 \cdot \sigma\left(\frac{1}{T}\sum_t x_t\right) + 0.5 \cdot f_{\text{fall}}\right) \\
r_{\text{long}} &= \clip\left(\left(0.58 \cdot r_{\text{base}} + 0.42 \cdot \sigma\left(\frac{1}{36}\sum_{t=T-35}^{T} x_t\right)\right) \cdot \mu_{\text{chronic}} + g_{\text{pressure}}\right)
\end{align}

\textbf{短期风险 $r_{\text{short}}$ 的逐项解析：}
\begin{itemize}[nosep]
\item $f_{\text{fused}}$：门控融合特征（式 §4.3），整合了 TCN（短期梯度）、EMA（长期累积）、Transformer（全局注意力）三类信息。在 $r_{\text{short}}$ 中，$f_{\text{fused}}$ 是输入的核心项——它携带了当前时刻所有可用的时序特征。
\item $b_{\text{auth}}$：权威偏置——如果输入来自权威临床来源（如认证的 EHR 系统）则为正，增强了信号的可靠性。
\item $f_{\text{fall}}$：跌倒特征——即使在没有发生跌倒事件时（$\delta_{\text{fall}} = 0$），跌倒\textbf{风险}（通过姿势不稳定性和活动能力测量）仍然被考虑。
\item $\delta_{\text{fall}}$：跌倒事件——与 $f_{\text{fall}}$（跌倒风险）不同，$\delta_{\text{fall}}$ 指示\textbf{已发生的}跌倒事件。两者同时出现时（高跌倒风险 + 实际跌倒），风险信号叠加。
\item $\mu_{\text{acute}}$：急性病乘数——在四项信号的 $\sigma$ 映射后的概率上施加的最终倍增。注意 $\mu_{\text{acute}}$ 作用于 $\sigma(\cdot)$ 的\textbf{外部}（乘法），而非内部（加法）——这意味着它是一个全局缩放因子，而非仅影响信号强度。
\end{itemize}
$r_{\text{short}}$ 公式的设计哲学是：\textbf{短期风险由"当前时刻的所有信号"决定，历史仅通过融合特征 $f_{\text{fused}}$ 间接影响。}公式中没有直接引用历史均值或 $r_{\text{base}}$——短期预测"活在当下"。

\textbf{中期风险 $r_{\text{mid}}$ 的逐项解析：}
\begin{itemize}[nosep]
\item $0.72 \cdot r_{\text{base}}$：基础风险占主导权重 0.72。$r_{\text{base}}$ 包含了所有历史信息和特征——它是长期趋势的浓缩。0.72 的高权重意味着中期预测"以历史为锚"。
\item $0.28 \cdot \sigma(\frac{1}{T}\sum_t x_t)$：近期观测均值经逻辑函数映射后的贡献。$\frac{1}{T}\sum_t x_t$ 是全部 72 个时间点的简单算术平均——它代表了"整体的平均水平"。权重 0.28 表明近期整体的平均水平对中期预测有温和的修正作用。
\item $0.5 \cdot f_{\text{fall}}$：跌倒风险特征（注意这里没有 $\delta_{\text{fall}}$——跌倒事件只影响短期，不影响中期，因为跌倒后的直接影响在 24 小时内显现）。系数 0.5 表明"活动能力低下的状态"对中期风险有中等影响。
\item $\clip(\cdot)$：将输出裁剪到合法概率区间。$\clip$ 操作确保最终输出的 $r_{\text{mid}} \in [0,1]$，即使内部计算因数值误差略微超出范围。
\end{itemize}
$r_{\text{mid}}$ 公式的设计哲学是：\textbf{中期风险是"历史的惯性"（70\%权重）与"近期的整体水平"（28\%权重）的折中。} $0.72:0.28$ 的比例体现了"远期因素占主导，近期因素提供微调"的先验。

\textbf{长期风险 $r_{\text{long}}$ 的逐项解析：}
\begin{itemize}[nosep]
\item $0.58 \cdot r_{\text{base}}$：基础风险权重略低于中期（0.58 vs 0.72）——因为长期预测中 $r_{\text{base}}$（主要反映短期到中期的动力学）的适用性下降，需要更多的慢性因素补充。
\item $0.42 \cdot \sigma(\frac{1}{36}\sum_{t=T-35}^{T} x_t)$：最近 36 个时间步（半个总窗口）的均值。注意这里取了最近 36 步而非全部 72 步——因为 $r_{\text{long}}$ 虽然关注长期，但仍需要"近期的慢性趋势"（而非遥远的陈旧信息）。使用半个窗口作为平滑范围是一种经验上的折中。
\item $\mu_{\text{chronic}}$：慢性病程乘数。与 $\mu_{\text{acute}}$ 作用于 $r_{\text{short}}$ 类似，$\mu_{\text{chronic}}$ 作用于 $r_{\text{long}}$——慢性病对长期风险的影响是最核心的。$\mu_{\text{chronic}} = 2.0$ 使慢性病患者的长期风险翻倍。
\item $g_{\text{pressure}} = 1 - e^{-r_p \cdot h_{\text{bed}}}$：压疮增长模型。这是一个\textbf{指数饱和函数}——当 $h_{\text{bed}} \to \infty$ 时，$g_{\text{pressure}} \to 1$；当 $h_{\text{bed}} = 0$ 时，$g_{\text{pressure}} = 0$。$r_p = 0.08$ 控制饱和速度。此模型的物理解释见下方"压疮增长的物理建模"段落。
\item $\clip(\cdot)$：确保最终输出在 $[0,1]$ 内（特别是含 $g_{\text{pressure}}$ 后可能超过 1 的情况）。
\end{itemize}
$r_{\text{long}}$ 公式的设计哲学是：\textbf{长期风险由"慢性因素 + 累积暴露"主导，即时波动被大幅平滑。} $0.58:0.42$ 的比例（长期 vs 近期）反映了"慢性因素占微弱优势，但近期半窗口的趋势仍然提供重要信息"。

\subsubsection{权重分配的时间衰减模式}

三个窗口的权重分配揭示了系统对"信息时效性"的定量信念：
\begin{itemize}[nosep]
\item $r_{\text{short}}$：即时信号权重 100\%，历史权重 0\%（$r_{\text{base}}$ 不出现在公式中）
\item $r_{\text{mid}}$：$r_{\text{base}}$（历史浓缩）权重 72\%，近期均值权重 28\%
\item $r_{\text{long}}$：$r_{\text{base}}$ 权重 58\%，近期半窗口均值权重 42\%
\end{itemize}
注意一个看似矛盾的现象：$r_{\text{long}}$ 中近期均值权重 (0.42) 反而比 $r_{\text{mid}}$ 中近期均值权重 (0.28) 更高。这个"悖论"的解释是：$r_{\text{long}}$ 中的近期均值使用了更长的平滑窗口（36 步 vs 72 步），使得"近期均值"在长期预测中实际上反映了更平滑的、更接近"趋势"而非"波动"的信息。而在 $r_{\text{mid}}$ 中，近期均值是全部 72 步的均值，包含了更早的历史信息，因此权重更低以让主导权给 $r_{\text{base}}$。

\subsubsection{压疮增长的物理建模}

压疮（压力性损伤，俗称褥疮）是长期卧床患者最危险的并发症之一。其风险增长遵循\textbf{饱和指数模型}：
\[
g_{\text{pressure}} = 1 - e^{-r_p \cdot h_{\text{bed}}}
\]
这一模型的深层逻辑来自\textbf{可靠性工程中的失效时间建模}：
\begin{itemize}[nosep]
\item 如果压疮的发生是一个泊松过程（在任意小时段内以恒定速率 $r_p$ 发生），则到时间 $t$ 至少发生一次压疮的概率为 $1 - e^{-r_p t}$——恰好是指数饱和模型。
\item $r_p = 0.08$ 意味着压疮事件的平均到达间隔为 $1/0.08 = 12.5$ 小时。换句话说，在每 12.5 小时的卧床时间内，压疮风险"潜力"增加约 63\%（$1 - 1/e$）。
\item 模型的指数饱和特性反映了\textbf{边际风险递减}——卧床第 1 小时的风险增长速率远快于第 48 小时。这是因为压疮的形成是一个\textbf{累积过程}（组织缺血→细胞损伤→坏死），初期的缺血损伤是最关键的。
\item 当患者卧床时间很长时（如 $h_{\text{bed}} \to \infty$），$g_{\text{pressure}} \to 1$（压疮几乎不可避免），这与临床观察一致——不进行翻身护理的长期卧床患者压疮发生率接近 100\%。
\end{itemize}

\begin{remark}[多窗口预测的多尺度物理类比]
三个时间窗口对应临床中的三个风险时间尺度：
\begin{itemize}[nosep]
\item \textbf{短期 (24h)}：受急性因子（感染、跌倒）主导，$\mu_{\text{acute}} \in \{1,3\}$ 在急性病存在时将风险乘以3倍
\item \textbf{中期 (30d)}：短期趋势与长期基线以 $0.72:0.28$ 的比例混合——权重反映了"近期信息比远期信息更相关"的指数衰减先验
\item \textbf{长期 (12m)}：受慢性因子主导，基线权重更高 ($0.58$)，且引入了压疮增长的指数饱和模型 $g_{\text{pressure}} = 1 - e^{-r_p h_{\text{bed}}}$——卧床时间越长，压疮风险以递减速率增加（边际效应递减）
\end{itemize}
逻辑函数 $\sigma(z)$ 在物理中描述费米子在各能级上的占据概率；
在风险预测中，它提供了一种将无界实数信号映射到 $[0,1]$ 概率的自然方式。
\end{remark}

\subsection{DMA 置信度校准}

TFA 根据 DMA 置信度 $\gamma_{\text{DMA}}$ 调整自身预测：

\begin{equation}
r' = r \cdot
\begin{cases}
1.0 & \text{if } \gamma_{\text{DMA}} \geq 0.9 \\
0.9 & \text{if } 0.7 \leq \gamma_{\text{DMA}} < 0.9 \\
0.8 & \text{if } 0.6 \leq \gamma_{\text{DMA}} < 0.7 \\
1.0 \text{ (degraded)} & \text{if } \gamma_{\text{DMA}} < 0.6
\end{cases}
\end{equation}

\textbf{校准策略的贝叶斯解释：} 当 DMA 置信度极高 ($\geq 0.9$) 时，诊断已经非常确定，
TFA 不需要调整风险预测（诊断信息已经充分）。当 DMA 置信度中等 ($0.6$--$0.9$) 时，
诊断的不确定性传递到风险预测——乘数 < 1 反映了"如果我们不确定诊断，
那么基于该诊断的风险预测也应该更保守"。当 DMA 置信度极低 ($< 0.6$) 时，
乘数恢复为 1.0 但标记 degraded——这意味着 DMA 的置信度太低以至于不可靠，
TFA 不应根据不可靠的信息调整自身预测（不传递虚假的降低信号）。

\subsection{MedTsLLM-v1.5 路径（真实推理）}

当 MedTsLLM 可用时，TFA 切换到基于 BioBERT 的深度推理：

\textbf{时间序列标准化 (RevIN):}
\begin{equation}
\tilde{x} = \frac{x - \mu_B}{\sqrt{\sigma_B^2 + \epsilon}}, \quad \mu_B = \frac{1}{T}\sum x_t, \quad \sigma_B^2 = \frac{1}{T}\sum(x_t - \mu_B)^2
\end{equation}

\textbf{分块嵌入 (Patch Embedding):}
\begin{equation}
\mathbf{P}_{\text{patch}} = \text{Conv1D}(\text{Unfold}(\tilde{x})), \quad \text{patch\_len}=12, \text{stride}=6
\end{equation}

\textbf{重编程层 (Reprogramming Layer) — 交叉注意力：}
\begin{equation}
\mathbf{Q} = \mathbf{P}_{\text{patch}} \mathbf{W}_Q \in \mathbb{R}^{B \times L \times H \times E}, \quad
\mathbf{K} = \mathbf{E}_{\text{word}} \mathbf{W}_K, \quad
\mathbf{V} = \mathbf{E}_{\text{word}} \mathbf{W}_V
\end{equation}
\begin{equation}
\text{Attention}(\mathbf{Q}, \mathbf{K}, \mathbf{V}) = \softmax\left(\frac{\mathbf{Q}\mathbf{K}^{\top}}{\sqrt{E}}\right) \mathbf{V}
\end{equation}

其中 $\mathbf{E}_{\text{word}}$ 是 BioBERT 的冻结词嵌入，重编程层将时序分块信号"翻译"为大语言模型可理解的语义空间。

\textbf{BioBERT 编码 + FlattenHead:}
\begin{equation}
\mathbf{h} = \text{BioBERT}\left([\text{PromptEmb}, \text{Reprogrammed}]\right)_{\text{last\_hidden}}
\end{equation}
\begin{equation}
\hat{\mathbf{r}} = \sigma\left(\text{FlattenHead}(\mathbf{h}_{[-n_{\text{patches}}:, :d_{\text{ff}}]})\right) \in [0,1]^{3 \times 3}
\end{equation}

\subsection{风险等级分类}

\begin{equation}
\text{RiskLevel}(r) =
\begin{cases}
\text{critical} & \max(r_s, r_m, r_l) \geq 0.85 \\
\text{high} & \max(r_s, r_m, r_l) \geq 0.65 \\
\text{medium} & \max(r_s, r_m, r_l) \geq 0.35 \\
\text{low} & \text{otherwise}
\end{cases}
\end{equation}

\textbf{阈值选择的临床逻辑：} 四个阈值 $\{0.35, 0.65, 0.85\}$ 将风险概率空间 $[0,1]$
划分为四个区域，这些间隔并非均匀分布。它们反映了临床决策中的\textbf{非对称损失函数}：
假阴性（漏诊）的代价远大于假阳性（误报）。因此 medium 阈值 ($0.35$) 设置得相对较低——
即使只有 35\% 的风险概率，也值得引起注意（进入 medium 级别）。

% ============================================================
\section{知识融合智能体 (Fusion Agent)}

\subsection{从检索结果到结构化知识：为什么需要专门的融合智能体}

RAA（§2）检索到的是\textbf{文档}——未经过滤、可能重复、可能冲突的自然语言文本集合。直接将检索结果原样传递给 DMA 会导致三个问题：(1) 重复文档对决策产生不成比例的影响；(2) 同一疾病领域下内容矛盾的文档使证据方向不一致；(3) 文档的权威性、时效性和内容质量参差不齐。知识融合智能体（Fusion Agent）在 RAA 和 DMA 之间充当\textbf{知识质量控制层}——去重、冲突检测、信源加权，确保进入 DMA 决策流程的知识是清洁的、一致的和可信的。

\subsection{SimHash 指纹与去重}

\subsubsection{SimHash 的数学原理}

SimHash 是一种局部敏感哈希 (LSH)，其核心性质是：语义相似的文档产生汉明距离小的哈希码。
在数学上，SimHash 将文档从高维词频空间 $\mathbb{R}^{V}$ 映射到 64 维超立方体
$\mathbb{H}^{64} \cong (\mathbb{Z}_2)^{64}$，并保持角度相似性（内积的符号一致性）。

文档 $d_i$ 的 SimHash 指纹计算：

\begin{equation}
\operatorname{simhash}(d_i) = \sum_{t \in \text{tokens}(d_i)} \sum_{b=0}^{63} \mathbb{I}[\text{bit}_b(\text{MD5}(t)) = 1] \cdot 2^b
\end{equation}

实际使用加权累加器：
\begin{equation}
\operatorname{Acc}[b] = \sum_t \begin{cases} +1 & \text{if } \text{bit}_b(\text{MD5}(t)) = 1 \\ -1 & \text{otherwise} \end{cases}
\end{equation}
\begin{equation}
\operatorname{simhash}(d_i) = \sum_{b=0}^{63} \mathbb{I}[\operatorname{Acc}[b] \geq 0] \cdot 2^b
\end{equation}

\textbf{物理解释：} 将每个 token 视为超立方体中的随机向量（通过 MD5 哈希的雪崩效应保证均匀分布）。
文档的 SimHash 是所有 token 向量的"多数投票"——在每个维度上，
如果多数 token 的比特为 1，则该维度取 1。这使得文档指纹在超立方体中
位于其所有 token 指向方向的"中心区域"。

去重条件（冲突检测的反面）：
\begin{equation}
\text{dup}(d_i, d_j) \iff \text{Hamming}(\operatorname{simhash}(d_i), \operatorname{simhash}(d_j)) \leq \tau_{\text{dup}} = 4
\end{equation}

\begin{remark}[$\tau_{\text{dup}} = 4$ 的统计含义]
对于均匀随机分布的 64 位哈希码，任意两个哈希码的汉明距离近似服从
$\mathcal{N}(32, 16)$。$\tau_{\text{dup}} = 4$ 意味着我们要求两个文档的哈希码在
超立方体中非常接近（距离远小于期望值 32）。在统计上，随机两个文档满足
$\Ham \leq 4$ 的概率约为 $2^{-64} \cdot \sum_{k=0}^{4} \binom{64}{k} \approx 10^{-12}$，
几乎不可能来自偶然——这确保了去重的高精度。

上述概率计算基于均匀随机哈希的理想化假设。实际中 SimHash 的分布可能偏离均匀分布，
因此阈值 $\tau_{\text{dup}} = 4$ 需要在具体知识库上通过实验验证去重精度。
\end{remark}

\subsection{冲突检测}

\begin{equation}
\text{conflict}(d_i, d_j) \iff \text{area}(d_i) \neq \text{area}(d_j) \land \text{Hamming}(\operatorname{simhash}(d_i), \operatorname{simhash}(d_j)) > \tau_{\text{conflict}} = 26
\end{equation}

大哈密顿距离 $\tau_{\text{conflict}} > 26$ 指示两个文档在语义空间中的方向完全分离——
物理上等价于在64维超立方体 $(\mathbb{Z}_2)^{64}$ 中占据对极位置。

\textbf{对极位置的几何意义：} 将汉明距离 $\Ham > 26$ 转换回角度空间：
在 64 维超立方体中，汉明距离 $k$ 对应球面上的期望夹角约为 $\pi \cdot k/64$。
$\Ham = 26$ 对应夹角 $\approx 0.41\pi \approx 73^\circ$，
意味着两个文档在语义球面上的方向至少偏离了 $73^\circ$——
几乎正交，代表了显著的语义分歧。

\subsection{区域共识得分}

\begin{equation}
s_{\text{consensus}}(a) = \frac{1}{|\mathcal{D}_a|}\sum_{d \in \mathcal{D}_a} s_{\text{sim}}(d), \quad \mathcal{D}_a = \{d \in \text{unique\_docs}: \text{area}(d) = a\}
\end{equation}

\textbf{共识得分的统计力学解释：} $s_{\text{consensus}}(a)$ 是领域 $a$ 中文档群体的
"系综平均"。在统计力学中，系综平均反映系统的宏观可观测量；
在这里，它反映该疾病领域中知识的一致性程度——高共识分意味着该领域的知识高度一致
（多个来源指向相同结论），低共识分意味着该领域存在争议或不确定性。

% ============================================================
\section{验证智能体 (Verification Agent)}

\subsection{事实核对评分}

\begin{equation}
s_{\text{fact}}(d) = \clip_{[0.0, 1.0]}\left(0.38 + b_{\text{auth}} + b_{\text{content}} + b_{\text{area}} + b_{\text{med}} + b_{\text{evidence}} - p_{\text{redflag}} - p_{\text{vague}}\right)
\end{equation}

各子项：
\begin{itemize}[nosep]
\item $b_{\text{auth}} = 0.20$ 若来源为权威源或人类专家
\item $b_{\text{content}} = \min(0.30, 0.06\sqrt{|\text{tokens}(d)|})$
\item $b_{\text{area}} = 0.10$ 若领域非 Unknown
\item $b_{\text{med}} = \min(0.15, 0.04 \cdot \#\text{med\_terms\_hit})$
\item $b_{\text{evidence}} = 0.14$ 若含定量数据或研究引用
\item $p_{\text{redflag}} = 0.35$ 若含伪科学/危险关键词
\item $p_{\text{vague}} = 0.15$ 若标记数 < 12
\end{itemize}

\begin{remark}[评分函数的结构]
$s_{\text{fact}}$ 的结构是一个\textbf{加性证据模型}：
基线分 ($0.38$) + 正证据项 ($b_*$) $-$ 负证据项 ($p_*$)。
基线分 $0.38$ 代表"无信息先验"——没有任何证据时，文档的正确概率基准。
$b_{\text{content}} = 0.06\sqrt{n}$ 使用平方根标度（而非线性）反映信息量的
\textbf{递减边际效益}——200个token带来的信息增量并非100个token的2倍，而是约1.4倍。
$p_{\text{redflag}} = 0.35$ 是最大的单项惩罚——一旦检测到危险信号，
文档可信度大幅下降（从 0.38 降至 0.03），反映了医疗领域对安全性的极致要求。
\end{remark}

\subsection{时效性检查}

\begin{equation}
\text{timeliness\_ok}(d) \iff \Delta t = (\text{now} - \text{last\_updated}(d))_{\text{days}} \leq 365
\end{equation}

\subsection{TFA 报告验证}

对 TFA 输出执行四项检查：风险得分范围检查（$r \in [0,1]$）、风险因子-诊断关联性检查、
推荐-风险等级一致性检查、退化通知检查。

% ============================================================
\section{隐私智能体 (Privacy Agent)}

\subsection{差分隐私的数学基础}

隐私智能体在信息论安全与计算效率之间取得平衡，其核心是高斯差分隐私机制。

\subsection{联邦梯度扰动}

\begin{equation}
\nabla^{\text{secure}} = \nabla^{\text{raw}} + \boldsymbol{\epsilon}, \quad \epsilon_i \sim \mathcal{N}(0, \sigma_{\text{noise}}^2), \quad \sigma_{\text{noise}} = 0.015
\end{equation}

这等同于满足局部 $(\epsilon, \delta)$-差分隐私的高斯机制，其中隐私预算 $\epsilon$ 由 $\sigma_{\text{noise}}$ 和敏感度 $\Delta f$ 决定：
\begin{equation}
\epsilon \approx \frac{\Delta f}{\sigma_{\text{noise}}} \cdot \sqrt{2\log(1.25/\delta)}
\end{equation}

\begin{remark}[$\sigma_{\text{noise}} = 0.015$ 的选择依据]
噪声标准差 $0.015$ 是在隐私保护与模型效用之间的工程权衡。

\textbf{敏感度分析：} 在 HMAE 的实际实现中，每个样本的梯度经 L2 裁剪至最大范数 $C = 1.0$，
因此敏感度 $\Delta f = C = 1.0$。代入高斯机制隐私预算公式：
$\sigma \geq \frac{\Delta f}{\varepsilon} \cdot \sqrt{2\log(1.25/\delta)}$，
取 $\delta = 10^{-5}$，可得 $\varepsilon \approx 2.2$，满足医疗场景下的差分隐私要求。
实际实现中使用梯度裁剪（norm clip $= 1.0$）。

在 HMAE 中，梯度值通常经 L2 归一化后分布于 $[-1, 1]$，因此 $\sigma = 0.015$ 意味着
信噪比 SNR $\approx 1/0.015^2 \approx 4400$（在方差意义上）。
这提供了足够强的隐私保护（任何单个数据点的影响被噪声模糊化），
同时保持梯度方向的基本可用性。
\end{remark}

\subsection{XOR 流加密}

\begin{equation}
c_k = p_k \oplus \text{key}_{k \bmod |K|}, \quad K = \text{SHA-256}(\text{secret\_key})
\end{equation}

物理意义上，这是具有一次性密码本安全特性的流加密方案——密钥长度等于消息长度时达到信息论安全。
SHA-256 的输出是 256 位的均匀分布伪随机序列，对抗已知明文攻击的安全性依赖于 SHA-256 的原像抗性。

% ============================================================
\section{维护智能体 (Maintenance Agent)}

\subsection{设计动机与系统角色}

维护智能体（Maintenance Agent）是 HMAE 的"免疫系统"——它的任务不是参与诊断或预测，而是持续监控所有其他智能体的健康状态，在组件出现异常时第一时间检测并通知退化管理器。与功能智能体（Perception、DMA、TFA 等）不同，维护智能体是\textbf{元层面的（meta-level）}——它观察系统本身而非临床数据。

维护智能体的核心职责有三：
\begin{enumerate}[nosep]
\item \textbf{心跳监测（heartbeat monitoring）：} 每个智能体定期向维护智能体发送"我还活着"信号。如果超过阈值时间未收到信号，标记该智能体为故障。
\item \textbf{输出异常检测（output anomaly detection）：} 检查每个智能体的输出是否"合法"——是否为 None、NaN、Inf 或其他在数学上不合法的值。
\item \textbf{状态汇总与上报（status aggregation）：} 将各智能体的健康状态汇总为系统退化等级 $L \in \{0,1,2,3\}$，通知退化管理器执行相应的降级策略。
\end{enumerate}

\subsection{心跳监测}

对每个已注册智能体 $a \in \mathcal{A}$：

\begin{equation}
\text{status}(a) =
\begin{cases}
\text{faulty} & \text{if } \Delta t = \text{now} - a.\text{last\_heartbeat} > \tau_{\text{timeout}} = 120\text{s} \\
a.\text{status} & \text{otherwise}
\end{cases}
\end{equation}

\textbf{公式详解：} 心跳检查的逻辑极其简单——这是有意为之。在安全关键系统中，\textbf{监控机制本身的复杂性不应超过被监控系统的复杂性}（否则监控机制自身的故障概率可能与被监控组件相当）。公式检查的是单一条件：当前时间与最后一次心跳的时间差是否超过 120 秒。

参数 $\tau_{\text{timeout}} = 120\text{s}$ 的选择不是一个任意的数字，而是三方面权衡的结果：
\begin{itemize}[nosep]
\item \textbf{太短（如 10s）：} 网络抖动或瞬时 CPU 争用可能导致健康组件被误判为故障（假阳性），引发不必要的退化切换。在分布式系统中，这被称为"flapping"——系统在正常/退化模式之间来回震荡。
\item \textbf{太长（如 600s）：} 真正的故障需要 10 分钟才能被检测到，在医疗实时场景中这是不可接受的——10 分钟的"静默故障"足以导致未经审查的诊断被传递给临床医生。
\item \textbf{120s 的工程依据：} HMAE 中大多数智能体的一次完整调用（从接收到返回结果）耗时 $< 5$ 秒。Llama3:8b 的推理耗时约 $2$--$5$ 秒（取决于提示长度），TF-IDF 质心分类约 $< 0.5$ 秒。120 秒是正常响应时间的 24--240 倍，足以排除瞬时的性能抖动，同时又远小于临床决策的时间窗口（分钟级），确保故障在造成影响之前被检测。
\end{itemize}

\textbf{超时阈值的控制论解释：} $\tau_{\text{timeout}} = 120s$ 定义了一个
"软实时"约束。在控制论中，这是一个\textbf{死区时间（dead time）}——
超过此时间无响应的组件被视为失效，系统切换到降级模式。

\subsection{异常输出检测}

\begin{equation}
\text{abnormal}(y) =
\begin{cases}
\text{True} & \text{if } y \text{ is None} \lor \text{NaN}(y) \lor \text{Inf}(y) \\
\text{True} & \text{if } y.\text{confidence} \text{ is None} \lor \text{NaN}(y.\text{confidence}) \\
\text{False} & \text{otherwise}
\end{cases}
\end{equation}

\textbf{公式详解——为什么要分两个条件检查：} 异常输出检测公式检查两个独立的异常模式，它们对应于两种不同的故障根源：
\begin{itemize}[nosep]
\item \textbf{条件1（$y$ is None $\lor$ NaN $(y)$ $\lor$ Inf$(y)$）：} 整个输出对象 $y$ 本身就是非法的——智能体抛出了 null 返回、或在计算中产生了未定义的数值（除以零 → NaN，溢出 → Inf）。这通常指示智能体内部的致命错误（如模型加载失败、除零错误、内存溢出），需要立即标记该智能体为异常。
\item \textbf{条件2（$y.\text{confidence}$ is None $\lor$ NaN$(y.\text{confidence})$）：} 输出对象本身存在（不是 None），但其中的置信度字段异常——这意味着智能体虽然运行了（返回了某种输出），但其内部的置信度估计模块出现了问题。在这种情况下，输出可能部分可用（如诊断类别仍然有意义），但置信度不可靠，系统应标记为 degraded 而非完全 faulty。
\end{itemize}
这两个条件的分层检查体现了"精确故障定位（precise fault localization）"的原则——区分"彻底失效"和"部分降级"两种状态，使退化管理器可以选择最小成本的降级策略（而非动辄切换到 L3 安全模式）。

注意异常检测的逻辑是保守的——任何不清楚的情况（None/NaN/Inf）都被标记为异常。在医疗系统中，\textbf{宁可因为过度谨慎而产生假阳性异常检测（不必要的降级），也不可因为检测过于宽松而遗漏真正的异常输出}。这是安全关键系统设计的核心原则。

% ============================================================
\section{通信拓扑智能体 (Topology Manager)}

\subsection{设计动机：为什么通信拓扑需要单独管理}

在一个拥有 10 个智能体的系统中，如果每个智能体都可以直接与任意其他智能体通信，则存在 $10 \times 9 / 2 = 45$ 条潜在的双向通信链路。不加管理的全对等通信会导致三个严重问题：(1) \textbf{状态同步的复杂性爆炸}——每个智能体需要维护另外 9 个智能体的地址和状态；(2) \textbf{消息路由的不确定性}——同一条消息可能通过多条路径到达目标，导致重复处理或乱序；(3) \textbf{故障传播}——一个智能体的崩溃可能通过级联的超时和重试影响全局。

通信拓扑智能体（Topology Manager）通过\textbf{集中管理通信图}并\textbf{动态选择拓扑结构}来解决这些问题。它的数学核心是图论——将智能体集合 $\mathcal{A}$ 视为顶点集，通信链路视为边集，拓扑选择等价于在不同的边集之间切换以优化不同的图论指标。

\subsection{网络模型}

\subsubsection{星型拓扑 —— 中心化路由}

\textbf{星型拓扑：} 中心节点为 DMA，所有通信经 DMA 路由：
\begin{equation}
\mathcal{E}_{\text{star}} = \{(a, \text{DMA}) \mid a \neq \text{DMA}\} \cup \{(\text{DMA}, a) \mid a \neq \text{DMA}\}
\end{equation}

\textbf{公式解读：} 边集 $\mathcal{E}_{\text{star}}$ 包含两类有向边：
\begin{itemize}[nosep]
\item $\{(a, \text{DMA}) \mid a \neq \text{DMA}\}$：所有非 DMA 智能体到 DMA 的入站边（$|\mathcal{A}| - 1$ 条）
\item $\{(\text{DMA}, a) \mid a \neq \text{DMA}\}$：DMA 到所有非 DMA 智能体的出站边（$|\mathcal{A}| - 1$ 条）
\end{itemize}
总计 $2(|\mathcal{A}| - 1)$ 条有向边。对于 10 个智能体的系统，这是 $2 \times 9 = 18$ 条边——远少于网状拓扑的 $10 \times 9 = 90$ 条。星型拓扑的核心优势是\textbf{简单性和可预测性}：任意两个非中心节点之间的通信路径是唯一确定的（通过 DMA），不存在路由选择的歧义。

\subsubsection{网状拓扑 —— 全对等通信}

\textbf{网状拓扑：} 全对等通信：
\begin{equation}
\mathcal{E}_{\text{mesh}} = \{(a, b) \mid a, b \in \mathcal{A}, a \neq b\}
\end{equation}

\textbf{公式解读：} 边集包含所有有序对 $(a, b)$ 其中 $a \neq b$——每个智能体可以直接向任意其他智能体发送消息。总计 $|\mathcal{A}|(|\mathcal{A}| - 1)$ 条有向边。网状拓扑的核心优势是\textbf{低延迟（图直径为 1）和高容错性}：即使多个节点故障，剩余节点之间仍然可以直接通信，不需要通过中心节点中继。

\subsubsection{拓扑选择的分析框架}

\textbf{拓扑选择的图论分析：} 星型拓扑的图直径为 2（任意两个非中心节点通过 DMA 中继只需 2 跳），
通信复杂度为 $O(|\mathcal{A}|)$，适合节点数量适中、需要中心化控制的场景。
网状拓扑的图直径为 1，但通信复杂度为 $O(|\mathcal{A}|^2)$，
适合需要低延迟全对等通信的场景。HMAE 的拓扑管理器支持动态切换，
根据系统退化状态选择最优拓扑。

两种拓扑的定量比较：
\begin{table}[H]
\centering
\small
\begin{tabular}{lcc}
\toprule
\textbf{图论指标} & \textbf{星型拓扑} & \textbf{网状拓扑} \\
\midrule
边数 & $2(|\mathcal{A}|-1)$ & $|\mathcal{A}|(|\mathcal{A}|-1)$ \\
图直径 & 2 & 1 \\
平均路径长度 & $\approx 2 - 2/|\mathcal{A}|$ & 1 \\
节点连通度（中心/DMA） & $|\mathcal{A}|-1$ & $|\mathcal{A}|-1$ \\
节点连通度（非中心） & 1 & $|\mathcal{A}|-1$ \\
单点故障影响（中心节点故障） & 网络分裂 & 无影响 \\
单点故障影响（非中心节点故障） & 仅该节点断开 & 无影响 \\
\bottomrule
\end{tabular}
\caption{星型 vs 网状拓扑的图论指标对比}
\end{table}

这张对比表揭示了拓扑选择的根本权衡：星型拓扑的\textbf{控制性好}（所有通信可追踪、可审计）但\textbf{中心节点是单点故障源}；网状拓扑的\textbf{鲁棒性高}（无单点故障）但\textbf{管理复杂度高}（$O(|\mathcal{A}|^2)$ 条边的状态同步）。HMAE 的默认选择是星型拓扑——因为 DMA 本身就是信息流的自然中心（所有决策都经过 DMA），星型拓扑与此架构天然匹配。当 DMA 自身故障时（进入退化 L2 级别），拓扑管理器切换到网状拓扑以确保剩余智能体之间仍能通信。

\subsection{路由路径}

对于发送方 $s$ 和接收方 $r$：
\begin{equation}
\text{path}(s, r) =
\begin{cases}
[s, r] & \text{if } r \in \mathcal{N}(s) \\
[s, \text{DMA}, r] & \text{if } \text{DMA} \in \mathcal{N}(s) \land r \in \mathcal{N}(\text{DMA}) \\
\emptyset & \text{otherwise}
\end{cases}
\end{equation}

\textbf{公式解读——三种路由情况的精确含义：}
\begin{itemize}[nosep]
\item \textbf{情况1（直连，$r \in \mathcal{N}(s)$）：} 接收方 $r$ 在发送方 $s$ 的邻居集合中——路径为 $[s, r]$（单跳）。在网状拓扑下几乎所有节点对都属于这种情况；在星型拓扑下仅 DMA 与各节点之间的邻居对属于这种情况。
\item \textbf{情况2（DMA 中继）：} $r$ 不是 $s$ 的直接邻居，但 DMA 同时是 $s$ 和 $r$ 的邻居——路径为 $[s, \text{DMA}, r]$（两跳）。在星型拓扑下，所有非邻居节点对都满足此条件（所有节点都是 DMA 的邻居）。注意这里同时检查了 $\text{DMA} \in \mathcal{N}(s)$ 和 $r \in \mathcal{N}(\text{DMA})$——DMA 必须同时可达（s→DMA）和可转发（DMA→r）。
\item \textbf{情况3（不可达，$\emptyset$）：} 不存在从 $s$ 到 $r$ 的任何路径——返回空集。此时消息进入重试队列，等待拓扑变化或组件恢复后重新计算路由。
\end{itemize}

\subsection{消息重试队列}

传输失败的消息进入 FIFO 队列 $Q$，重试时重新计算路由路径，若仍不可达则放回队尾。

\textbf{FIFO 队列的公平性保障：} 先进先出（FIFO）是消息队列的最简单调度策略，但在故障恢复场景下有一个微妙的优点：\textbf{最早发送的消息最先被重试}。这意味着在系统从故障中恢复时，最"陈旧"（pending 最久）的消息获得优先处理——符合临床场景中"先到先处理"的时效性要求。

重试时重新计算路由路径（而非使用之前缓存的路径）是关键设计——因为拓扑可能在消息等待期间发生了变化（如 DMA 恢复了、或某个中继节点变得不可用）。使用过时的路由路径可能导致消息被发往已故障的节点，造成不必要的二次失败。

% ============================================================
\section{共识智能体 (Consensus Module)}

\subsection{从多元决策到统一输出：多智能体协调的必要性}

到此为止，HMAE 已经拥有多个独立的决策型智能体——DMA 给出诊断 $\hat{c}$ 和置信度 $\gamma_{\text{DMA}}$，TFA 给出三维风险向量 $\mathbf{r}_{\text{TFA}}$，验证智能体（§6）给出事实核查结论。这些智能体的输出可能\textbf{不一致}——DMA 高置信度地诊断为轻症，但 TFA 预测短期风险 critical。面对这种分歧，系统需要一个机制将多元输入融合为统一的临床决策。共识智能体（Consensus Module）正是这个机制——它不是另一个"意见提供者"，而是\textbf{元决策者}，制定了在多方意见分歧时如何裁决的规则。

\subsection{投票权重分配}

\begin{table}[H]
\centering
\begin{tabular}{lcc}
\toprule
\textbf{投票方} & \textbf{权重} & \textbf{角色} \\
\midrule
DMA (诊断) & 2 & 核心诊断决策 \\
TFA (时序风险) & 1 & 辅助决策支持 \\
Verification (验证) & 1 & 事实核查 \\
\midrule
总计 & 4 & 通过阈值 = 3 (75\%) \\
\bottomrule
\end{tabular}
\end{table}

\begin{remark}[加权投票的机制设计原理]
权重分配 $[2, 1, 1]$ 和通过阈值 $3/4 = 75\%$ 满足以下机制设计约束：
\begin{itemize}[nosep]
\item DMA (权重 2) 不能单独通过任何决策（$2 < 3$）——防止诊断独裁
\item DMA + 任一方 = 3，可以通过——诊断+辅助支持或诊断+事实核查即可决策
\item TFA + Verification (权重 1+1=2) 不足通过（$2 < 3$）——没有诊断支持的纯辅助信号不能独立决策
\end{itemize}
这是一种巧妙的\textbf{制衡（checks and balances）}设计：
DMA 作为核心但受制约，单一方不能独裁，多方联合才能形成决策。
在博弈论中，这构成一个\textbf{加权多数博弈（weighted majority game）}，
其 Shapley 值为 DMA > TFA = Verification，反映了 DMA 的核心地位。
\end{remark}

\subsection{三层级协同决策}

\textbf{第1层 — FORCE\_HITL（强制人工介入）}：若有任一条成立则绕过投票直接转人工：
\begin{itemize}[nosep]
\item DMA 置信度 < 阈值 0.62（或已触发 HITL 标记）
\item 验证严重冲突 $\geq 1$
\item TFA 风险等级 = critical
\item TFA-诊断冲突 $\geq 1$
\end{itemize}

\textbf{第2层 — DMA\_VETO（诊断否决）}：
\begin{equation}
\text{approved} = \text{False} \quad \text{if } \gamma_{\text{DMA}} < \tau_{\text{confidence}} = 0.62
\end{equation}

\textbf{第3层 — NORMAL\_VOTE（标准投票——置信度×风险矩阵）}：
\begin{equation}
\text{decision}(\gamma_{\text{DMA}}, r_{\text{TFA}}) =
\begin{cases}
\text{AUTO\_PASS} & \gamma \geq 0.8 \land r \in \{\text{low}, \text{medium}\} \\
\text{PASS\_RISK\_NOTED} & \gamma \geq 0.8 \land r = \text{high} \\
\text{AUTO\_PASS} & 0.6 \leq \gamma < 0.8 \land r = \text{low} \\
\text{ESCALATE} & 0.6 \leq \gamma < 0.8 \land r \in \{\text{medium}, \text{high}\} \\
\text{ESCALATE} & \gamma < 0.6
\end{cases}
\end{equation}

\begin{remark}[三层级决策的临床安全逻辑]
三层级架构体现了\textbf{安全优先（safety-first）}的设计原则：
\\ \textbf{第1层 (FORCE\_HITL)} - 最高优先级，绝对安全网：任何一条高危信号（极低置信度、严重冲突、critical 风险）直接触发人工介入，不经过任何投票。
\\ \textbf{第2层 (DMA\_VETO)} - DMA 作为"守门人"：即使通过第1层，DMA 仍然可以一票否决（当自身置信度不足时）。这是保证"不做出不可靠诊断"的最后防线。
\\ \textbf{第3层 (NORMAL\_VOTE)} - 常规操作：在安全前提下，通过置信度×风险矩阵进行精细决策。$\gamma \geq 0.8$ 且低/中风险时自动通过；$\gamma < 0.6$ 时一律升级——体现了"不确定时宁可升级也不冒险"的原则。
\end{remark}

% ============================================================
\section{纳什均衡检索策略选择}

\subsection{博弈论框架}

RAA 策略空间 $\mathcal{S} = \{\text{mixed}, \text{rerank}, \text{adaptive}\}$，DMA 模式空间 $\mathcal{M} = \{\text{fast}, \text{balanced}, \text{deep}\}$。

效用函数：
\begin{equation}
U(s, m) = \lambda \cdot P(\text{correct} \mid q, \mathcal{K}_s, m) - (1 - \lambda) \cdot C(s, |\mathcal{K}_s|)
\label{eq:utility}
\end{equation}

其中 $\lambda = 0.72$（精度-成本权衡因子），$C(s, n)$ 为检索成本：
\begin{equation}
C(s, n) =
\begin{cases}
0.22 + 0.012n & s = \text{mixed} \\
0.36 + 0.012n & s = \text{rerank} \\
0.28 + 0.012n + 0.08 \cdot \mathbb{I}[U(q) > 0.58] & s = \text{adaptive}
\end{cases}
\end{equation}

\textbf{$\lambda = 0.72$ 的决策论含义：} 精度权重远高于成本权重 ($0.72 > 0.28$)。
在医疗场景中，这是必然的选择——错误诊断的代价远大于额外计算资源。
$\lambda: (1-\lambda) \approx 2.57:1$ 的比率体现了"宁可多花计算也要保证准确性"的先验。

\subsection{最佳反应迭代}

\begin{align}
s^{(t+1)} &= \arg\max_{s \in \mathcal{S}} U(s, m^{(t)}) \\
m^{(t+1)} &= \arg\max_{m \in \mathcal{M}} \left[\lambda \cdot \hat{P}(\text{correct} \mid \mathcal{K}_{s^{(t+1)}}, m) - (1-\lambda) \cdot C(s^{(t+1)}, |\mathcal{K}|)\right]
\end{align}

迭代至纳什均衡 $(s^*, m^*)$ 满足：$U(s^*, m^*) \geq U(s, m^*)\; \forall s$ 且 $U(s^*, m^*) \geq U(s^*, m)\; \forall m$，最多迭代 $N_{\max} = 6$ 轮。

\subsection{收敛条件}

\begin{equation}
\text{Converged}^{(t)} \iff s^{(t)} = s^{(t-1)} \land m^{(t)} = m^{(t-1)}
\end{equation}

\begin{remark}[势博弈与收敛保证]
式 (\ref{eq:utility}) 定义了一个有限策略空间中的\textbf{势博弈（potential game）}——
势函数 $\Phi(s,m) = \lambda P - (1-\lambda)C$ 在每次最佳响应更新中单调不减：
$\Phi(s^{(t+1)}, m^{(t+1)}) \geq \Phi(s^{(t)}, m^{(t)})$。
由于策略空间有限（$3 \times 3 = 9$ 个纯策略组合），
势函数的单调性保证算法在最多 $9$ 步内收敛（实际限制 $N_{\max}=6$ 轮已足够检测循环）。
这一收敛性质是确定的、可证明的——不需要随机性假设。
\end{remark}

% ============================================================
% PART II: 其他部分
% ============================================================
\part{其他部分 —— Supporting Modules \& Infrastructure}

\vspace{6pt}
\begin{center}
\small
\textit{Part I 详细推导了十大核心智能体的数学基础。但这些智能体不是孤立运行的——}\\
\textit{它们依赖一组共享的基础设施：知识库、融合引擎、反馈学习机制和资源管理器。}\\
\textit{Part II 将这些基础设施的数学原理系统化，为 Part III 的跨模块耦合分析奠定基础。}
\end{center}
\vspace{12pt}

\section{动态知识库 (Online Knowledge Base)}

\subsection{增量知识更新}

\begin{equation}
\mathcal{K}_{t+1} = \mathcal{K}_t \cup \{\text{Embed}(L_{\text{new}})\}
\end{equation}

其中 $L_{\text{new}}$ 为来自 PubMed、arXiv、UpToDate、IEEE Xplore 的新文献。

\textbf{单调性保证：} 知识库的更新是\textbf{单调非减的}——$\mathcal{K}_{t+1} \supseteq \mathcal{K}_t$。
这意味着系统知识只增不减（除非人工删除），确保历史决策的可追溯性。
在逻辑上，这构成一个\textbf{增长的知识理论}——新知识可能与旧知识冲突（通过融合智能体的冲突检测处理），但不会消除旧知识。

\subsection{文档权重反馈更新}

\begin{equation}
w_{\text{doc}}^{(t+1)} = \clip\Big(w_{\text{doc}}^{(t)} + \eta \cdot F_{\text{user}} \cdot \nabla_{\text{proxy}}\Big)_{0.20}^{2.50}
\label{eq:weight_update}
\end{equation}

其中 $\eta = 0.08$（学习率），$F_{\text{user}} \in \{-1, 0, +1\}$（人类反馈信号），$\nabla_{\text{proxy}} = \|\text{MiniLM}(q)\|_2$（查询强度代理梯度）。

式 (\ref{eq:weight_update}) 是随即梯度下降在离散反馈空间中的实例化：文档权重沿反馈信号方向移动，以最大化未来检索质量。
裁剪区间 $[0.20, 2.50]$ 确保权重不会因极端反馈而发散——即使连续的负面反馈也不会使权重低于 0.20（保留基本的检索可能性），
连续的正面反馈也不会使权重超过 2.50（防止某些文档垄断检索结果）。

\subsection{FAISS 索引重建}

\begin{align}
\mathbf{M}_{\text{dense}} &= \text{MiniLM}(\mathcal{D}) \in \mathbb{R}^{|\mathcal{D}| \times d} \\
\text{Index}_{\text{FAISS}} &= \text{IndexFlatIP}(d) \\
\text{Index}_{\text{FAISS}}.\text{add}&(\mathbf{M}_{\text{dense}})
\end{align}

\section{三源知识融合 (Tri-Source Fusion)}

\subsection{信源加权融合}

给定三组文档 $\mathcal{D}_{\text{human}}$（人类经验）、$\mathcal{D}_{\text{auth}}$（权威知识）、$\mathcal{D}_{\text{ret}}$（检索知识），融合得分为：

\begin{equation}
s_{\text{fused}}(d) = \clip\left(s(d) \cdot w_{\text{source}} + 0.05 \cdot w_{\text{doc}}(d)\right)
\end{equation}

其中信源权重：$w_{\text{human}} = 0.95$, $w_{\text{auth}} = 0.90$, $w_{\text{ret}} = 0.75$, $w_{\text{synth}} = 0.55$。

融合操作在物理上等价于多信源的贝叶斯加权汇聚——高可信度信源具有更大的似然权重。

\begin{remark}[信源权重的层级结构]
权重递减链 $0.95 \to 0.90 \to 0.75 \to 0.55$ 反映了\textbf{证据可靠性层级}：
人类专家判断 > 权威出版物 > 通用检索结果 > 合成知识。
此层级结构不是线性的——人类与权威之间的差距 ($0.05$) 小于权威与通用检索之间的差距 ($0.15$)，
体现了"同行评审的知识与专家知识接近但次之，开放检索的未过滤知识需要显著折价"的认知论立场。
\end{remark}

\section{知识蒸馏引擎 (Knowledge Distillation)}

\subsection{加权平均上下文向量}

对 top-$k$ 文档，计算相似度加权的平均语义向量：

\begin{equation}
\mathbf{v}_{\text{distill}} = \Ltwo\left(\frac{\sum_{i=1}^{k} s_i \cdot \mathbf{v}_i}{\sum_{i=1}^{k} s_i}\right), \quad s_i = \max(s_{\text{sim}}(d_i), 10^{-3})
\end{equation}

此操作在余弦相似度度量下的几何意义是：在单位超球面 $S^{d-1}$ 上，求文档向量的 \textbf{Fr\'{e}chet 均值}（黎曼重心）。
与欧几里得空间中的算术平均不同，球面上的"平均"必须考虑球面曲率——
Fr\'{e}chet 均值是使得所有文档向量到它的测地线距离平方和最小的那个点：
$\mathbf{v}_{\text{distill}} = \arg\min_{\mathbf{v} \in S^{d-1}} \sum_i s_i \cdot d_{\text{geo}}(\mathbf{v}, \mathbf{v}_i)^2$。
加权平均公式是此优化问题在局部切空间中的一阶近似解。

\section{人在回路管理 (HITL Manager)}

\subsection{三层干预架构}

\textbf{第1层 — 自主执行：} 系统独立做出决策。
\textbf{第2层 — 人工干预判定：}
\begin{equation}
\text{HITL}^{(t)} \iff \gamma_{\text{DMA}} < \tau_{\text{conf}} \lor \bar{s}_{\text{ret}} < \tau_{\text{ret}} \lor (r_{\text{TFA}} \geq \tau_{\text{risk}} \land \text{high\_risk})
\end{equation}
其中 $\tau_{\text{conf}} = 0.62$, $\tau_{\text{ret}} = 0.24$, $\tau_{\text{risk}} = 0.65$。
\textbf{第3层 — 参数与检索更新：} 基于人类纠正更新 DMA 校准和检索权重。

\textbf{L1/L2/L3 的分层逻辑：} 这三层形成了完整的\textbf{自主-监督-学习}闭环：
L1 是正常操作模式（系统独立运行），L2 是安全网（条件触发人工审查），
L3 是反馈学习（将人工判断转化为参数更新）。三层合在一起构成了一个
\textbf{人在回路的持续改进系统 (Human-in-the-Loop Continuous Improvement System)}。

\subsection{TFA 特异性 HITL}

优先级系统：
\begin{itemize}[nosep]
\item Priority 1 (critical)：TFA 风险等级 critical —— 立即干预
\item Priority 2 (conflict)：TFA 与 DMA 诊断冲突 —— 紧急审查
\item Priority 3 (warning)：TFA 退化或验证警告 —— 建议审查
\end{itemize}

\subsection{错误分类与诊断}

三型错误体系：
\begin{align}
E_{\text{missing}} &: \text{检索覆盖中缺少正确疾病领域} \\
E_{\text{misleading}} &: \text{高分检索片段支持竞争疾病领域} \\
E_{\text{bias}} &: \text{证据存在但 DMA 决策边界需校准}
\end{align}

\textbf{三型错误的诊断价值：} 这三种错误的划分不是随意的——
它们分别对应诊断流程中的三个不同阶段：
$E_{\text{missing}}$ 是\textbf{检索阶段}的失败（知识库不完整），
$E_{\text{misleading}}$ 是\textbf{证据评估阶段}的失败（证据权重分配错误），
$E_{\text{bias}}$ 是\textbf{决策阶段}的失败（先验/偏置校准不当）。
通过区分这三类错误，HITL 可以针对性地修复系统的不同组件。

\section{模拟 LoRA 微调 (Simulated LoRA Fine-Tuner)}

\subsection{隐私感知梯度聚合}

\begin{equation}
\nabla_{\text{agg}} = \frac{1}{N}\sum_{i=1}^{N} \text{sign}_i \cdot \text{MiniLM}(q_i) \in \mathbb{R}^{d}
\end{equation}
\begin{equation}
\operatorname{sign}_i =
\begin{cases}
+1.0 & \text{label}(q_i) \neq \text{pred}(q_i) \\
+0.25 & \text{otherwise}
\end{cases}
\end{equation}
\begin{equation}
\nabla_{\text{secure}} = \nabla_{\text{agg}} + \boldsymbol{\epsilon}, \quad \epsilon_j \sim \mathcal{N}(0, 0.015^2)
\end{equation}

\textbf{非对称梯度信号的设计：} 误分类样本产生 $+1.0$ 的强梯度信号，
正确分类样本产生 $+0.25$ 的弱信号。这 $4:1$ 的比例反映了
\textbf{错误驱动的学习（error-driven learning）}——
从错误中学习的效率应远高于从成功中巩固的效率。
在神经科学中，这对应\textbf{预测误差驱动的突触可塑性}（dopamine-modulated STDP）。

\subsection{校准因子更新}

\begin{equation}
f_{\text{cal}}^{(t+1)} = \clip\Big(f_{\text{cal}}^{(t)} + 0.04 \cdot \frac{1}{N}\sum_i(1 - \gamma_i)\Big)_{0.80}^{1.35}
\end{equation}

\subsection{类比偏差累积}

\begin{equation}
b_c^{(t+1)} = b_c^{(t)} + 0.03 \cdot \frac{\#\{i: \text{corrected\_label}(i) = c\}}{N}
\end{equation}

\textbf{偏差累积的贝叶斯解释：} 类偏差 $b_c$ 实现了\textbf{在线贝叶斯先验更新}：
每当人类将某个预测纠正为类 $c$，该类别的先验概率微增——
这意味着系统学会了"在不确定时，稍微偏向这个经常被纠正为正确的类别"。
学习率 0.03 确保了单个纠正不会剧烈改变行为，而持续的模式会逐渐塑造决策边界。

\section{对比检索更新 (Contrastive Retrieval Updater)}

\subsection{检索裕度偏移}

\begin{equation}
\Delta m = \clip\Big(m + \eta_g \cdot \frac{N_+ - N_-}{N_+ + N_-}\Big)_{-0.25}^{0.25}, \quad \eta_g = 0.02
\end{equation}

此更新规则来源于对比损失 $\mathcal{L}_{\text{contrast}} = -\log\frac{\exp(s_+)}{\exp(s_+) + \exp(s_-)}$ 的梯度方向。

\textbf{对比学习的几何解释：} 对比损失在嵌入空间中施加一个"推-拉"动力学——
正样本（相关文档）被拉向查询向量，负样本（不相关文档）被推离查询向量。
最终的稳态是：正样本围绕查询向量形成一个紧致的簇，
负样本被驱离到球面的远处。

\section{退化管理 (Degradation Manager)}

\subsection{三级退化机制}

\begin{equation}
L =
\begin{cases}
0 \text{ (NORMAL)} & \text{unavailable} = 0 \land \text{degraded} = 0 \\
1 \text{ (MINOR)} & \text{unavailable} = 0 \land \text{degraded} \geq 1 \\
2 \text{ (PARTIAL)} & \text{unavailable} \geq 1 \lor \text{degraded} \geq 2 \\
3 \text{ (SEVERE)} & \text{unavailable} \geq 3
\end{cases}
\end{equation}

\texttt{safe\_call} 模式保证任意组件失败时系统残存运行：捕获异常 $\to$ 标记组件状态 $\to$ 返回回退值。

\begin{remark}[退化管理与控制论中的优雅降级]
退化等级 $L \in \{0,1,2,3\}$ 构成了一个\textbf{序数安全尺度}。
在控制论中，这对应\textbf{容错控制 (fault-tolerant control)} 的层级：
\\ $L=0$: 全功能运行
\\ $L=1$: 部分组件退化但核心功能保持（如 4-bit 量化替代全精度）
\\ $L=2$: 部分组件不可用但替代路径存在（如 Llama3 路径失效，回退至质心分类）
\\ $L=3$: 大部分组件失效，进入保守安全模式（仅使用最基础的规则决策）
\end{remark}

\section{持续学习循环 (Continuous Learning Loop)}

\subsection{触发条件与周期执行}

\begin{equation}
\text{should\_update} \iff |\mathcal{F}_{\text{buffer}}| \geq n_{\min} = 10
\end{equation}

每个周期：
\begin{enumerate}[nosep]
\item DMA Platt 校准更新（SGD 200 步）
\item 检索权重对比更新（最近20条反馈）
\item 反馈缓冲清空
\item 精度趋势估计：$\hat{a} = 1 - \#\text{corrections} / |\mathcal{F}_{\text{buffer}}|$
\end{enumerate}

\textbf{在线学习的工程约束：} $n_{\min} = 10$ 的选择是在"及时学习"和"统计稳定性"之间的权衡。
太小的批量（如 1）使每次更新噪声过大（高方差梯度估计），
太大的批量（如 100）延迟反馈的利用（系统适应慢）。
$10$ 是一个经验上的平衡点。

\section{MedTsLLM-v1.5 适配器}

\subsection{BioBERT 冻结骨干}

使用 dmis-lab/biobert-v1.1（108M 参数，$d_{\text{LLM}} = 768$）作为冻结编码器：
\begin{equation}
\forall \theta \in \Theta_{\text{BioBERT}}: \frac{\partial \mathcal{L}}{\partial \theta} = 0
\end{equation}

\textbf{冻结骨干的参数效率原理：} BioBERT 已经在生物医学语料上预训练，其内部表征已经
捕获了丰富的医学语义。冻结骨干确保这些宝贵的预训练知识不被小批量时序数据"灾难性遗忘"，
而仅通过轻量的 Reprogramming Layer 和 FlattenHead 进行任务特定适配。
总可训练参数量约为 BioBERT 的 2\%（~2M / 108M），大幅降低了部署和微调成本。

\subsection{时间序列→语言空间重编程}

核心数学操作：重编程层通过交叉注意力将离散时序分块映射到连续语义空间：

\begin{equation}
\text{Reprogrammed}(t) = \text{OutProj}\left(\text{softmax}\left(\frac{\mathbf{Q}(t)\mathbf{K}(\mathcal{W})^{\top}}{\sqrt{E}}\right) \mathbf{V}(\mathcal{W})\right)
\end{equation}

其中 $\mathcal{W}$ 为 BioBERT 词表嵌入，$\mathbf{Q}(t)$ 为时序分块查询，$\mathbf{K}(\mathcal{W}), \mathbf{V}(\mathcal{W})$ 为词汇键/值投影。

物理意义：Reprogramming Layer 是一个\textbf{信息瓶颈 (information bottleneck)} ——
它将高维连续时序信号压缩到离散语言符号空间，利用 LLM 的预训练语义先验重构预测。
压缩-重构的平衡由维度 $d_{\text{ff}} = 256$ 和头数 $n_{\text{heads}} = 8$ 控制。

\textbf{信息瓶颈的率失真理论解释：} 在 Tishby 的信息瓶颈理论中，
最优表征在"压缩输入信息"和"保持任务相关信息"之间取得平衡。
Reprogramming Layer 通过交叉注意力实现了一个\textbf{软瓶颈}——
时序信息必须通过离散语言符号构成的"窄通道"，
只有与 BioBERT 的语义空间兼容的信息分量才能通过。
这自然地过滤了时序噪声，保留了医学上可解释的信号。

\subsection{输出投影与风险解耦}

\begin{equation}
\mathbf{r}_{\text{output}} = \sigma\left(\text{FlattenHead}(\mathbf{h}_{\text{TS}})\right) \in [0,1]^{3 \times 3}
\end{equation}

三维输出分别对应：短(24h)、中(30d)、长(12m)期风险 $\times$ 三通道（恶化/稳定/恢复）。

% ============================================================
% PART III: 相互联系
% ============================================================
\part{相互联系 —— Inter-Module Coupling \& System Dynamics}

\vspace{6pt}
\begin{center}
\small
\textit{前两部分分别描述了智能体（Part I）和基础设施（Part II）的各自数学原理。}\\
\textit{但 HMAE 的真正力量不在于单个模块，而在于它们之间的协同——}\\
\textit{本部分将系统视为一个整体，分析数据如何跨模块流动、信息如何耦合、以及系统如何作为统一动力系统演化。}
\end{center}
\vspace{12pt}

\section{端到端流水线数据流}

\begin{figure}[H]
\centering
\small
\begin{tabular}{c@{\,}c@{\,}c@{\,}c@{\,}c@{\,}c@{\,}c@{\,}c@{\,}c}
$\underbrace{X_{\text{txt}} + X_{\text{img}}}_{\text{输入}}$ &
$\xrightarrow{\text{\S1}}$ &
$\underbrace{\mathbf{v}_{\text{fused}}}_{\text{感知}}$ &
$\xrightarrow{\text{\S2, \S11}}$ &
$\underbrace{\mathcal{K}_s}_{\text{检索}}$ &
$\xrightarrow{\text{\S5}}$ &
$\underbrace{\mathcal{D}_{\text{fused}}}_{\text{融合}}$ &
$\xrightarrow{\text{\S6}}$ &
$\underbrace{\mathcal{D}_{\text{verified}}}_{\text{验证}}$ \\[6pt]
\end{tabular}

\vspace{6pt}
\begin{tabular}{c@{\,}c@{\,}c@{\,}c@{\,}c@{\,}c@{\,}c}
$\xrightarrow{\text{\S7}}$ &
$\underbrace{\mathbf{v}_{\text{distill}}}_{\text{蒸馏}}$ &
$\xrightarrow{\text{\S3}}$ &
$\underbrace{\{P(c|q), \hat{c}\}}_{\text{DMA}}$ &
$\xrightarrow{\text{\S4}}$ &
$\underbrace{\{r_s, r_m, r_l\}}_{\text{TFA}}$ &
$\xrightarrow{\text{\S8, \S10}}$ &
$\underbrace{\text{decision}}_{\text{HITL + 共识}}$
\end{tabular}
\caption{端到端流水线数据流}
\end{figure}

\textbf{流水线的马尔可夫性质：} 注意数据流是单向的（虽然有反馈环路用于学习更新）。
这意味着一次推理的前向传播满足\textbf{马尔可夫性质}：
每个模块的输出仅依赖于前一个模块的输出（和自身的内部状态），
不直接依赖于更早的模块。这一性质极大地简化了系统的调试和分析——
错误可以通过逐模块的溯源来定位。

\section{耦合矩阵——跨模块信息依赖}

\begin{table}[H]
\centering
\small
\begin{tabular}{lp{6cm}}
\toprule
\textbf{耦合方向} & \textbf{耦合机制 / 数学关系} \\
\midrule
Perception $\to$ RAA & $\mathbf{v}_{\text{fused}}$ 作为上下文强度输入 \\
Perception $\to$ DMA & $\mathbf{v}_{\text{fused}}$ 作为 DMA 上下文加权 \\
RAA $\to$ DMA & $\mathcal{K}_s$（检索文档）$\to$ 检索得分、上下文构建 \\
RAA $\to$ TFA & 风险感知检索：$\mathcal{K}_{\text{risk}}$ $\to$ TFA-RA 一致性 \\
RAA $\leftrightarrow$ DMA (博弈论) & $s \leftrightarrow m$ 纳什均衡迭代 (式 \ref{eq:utility}) \\
DMA $\to$ TFA & $\gamma_{\text{DMA}}$ $\to$ TFA 置信度校准因子 \\
TFA $\to$ Consensus & $r_{\text{TFA}}$ $\to$ 投票 ($w_{\text{TFA}}=1$) \\
TFA $\to$ HITL & 风险等级 $\to$ 干预触发 \\
Verification $\to$ Consensus & 验证通过率 $\to$ 投票 ($w_{\text{Ver}}=1$) \\
HITL $\to$ DMA & 反馈 $\to$ Platt 校准 + 类偏差更新 \\
HITL $\to$ KB & 正/负文档 ID $\to$ 权重更新 (式 \ref{eq:weight_update}) \\
Fusion $\to$ Verification & 冲突文档 ID $\to$ 严格事实阈值 \\
Distillation $\to$ DMA & $\mathbf{v}_{\text{distill}}$ 与 $\mathbf{v}_{\text{fused}}$ 线性混合 \\
TFA $\to$ RAA & 风险等级 $\to$ 风险特定文档增强检索 \\
\bottomrule
\end{tabular}
\caption{跨模块耦合矩阵}
\end{table}

\textbf{耦合矩阵的图论分析：} 将此矩阵视为有向图 $G = (V, E)$，其中 $V$ 为模块集合，
$E$ 为耦合关系。该图具有以下性质：
\begin{itemize}[nosep]
\item \textbf{弱连通性：} 任何两个模块之间至少在一个方向上存在路径
\item \textbf{DMA 中心性：} DMA 的出度和入度都是最大的——它是信息流的核心枢纽
\item \textbf{闭环存在性：} HITL $\to$ DMA $\to$ TFA $\to$ HITL 构成反馈闭环——这是持续学习的基础
\item \textbf{无自环：} 所有耦合都是跨模块的，无模块直接反馈到自身
\end{itemize}

\section{通信拓扑：星型网络上的信息传播}

在星型拓扑下，所有消息经 DMA 中继：

\begin{equation}
\mathcal{M}_{A \to B}(t) = \text{Route}(A \to \text{DMA} \to B)(t)
\end{equation}

消息传递动力学（含隐私加密）：
\begin{equation}
\mathcal{M}_{\text{enc}} = \text{XOR}_{\text{SHA-256}}(\text{payload}), \quad
\mathcal{M}_{\text{dec}} = \text{XOR}_{\text{SHA-256}}^{-1}(\mathcal{M}_{\text{enc}})
\end{equation}

\section{系统整体平衡方程}

HMAE 系统可视为受以下约束的随机动力系统：

\begin{align}
\text{(状态演化)} \quad \mathbf{S}_{t+1} &= f(\mathbf{S}_t, \mathbf{u}_t, \boldsymbol{\xi}_t) \\
\text{(观测)} \quad \mathbf{y}_t &= g(\mathbf{S}_t, \boldsymbol{\eta}_t) \\
\text{(决策)} \quad a_t &= \pi(\mathbf{S}_t, \mathcal{K}_t) \\
\text{(反馈学习)} \quad \Theta_{t+1} &= \Theta_t - \eta \cdot \nabla_\Theta \mathcal{L}(\mathbf{y}_t^{\text{true}}, \mathbf{y}_t^{\text{pred}})
\end{align}

其中：
\begin{itemize}[nosep]
\item $\mathbf{S}_t$ — 系统状态（KB 索引、模型参数、文档权重）
\item $\mathbf{u}_t$ — 控制输入（检索策略 $s$、决策模式 $m$）
\item $\boldsymbol{\xi}_t, \boldsymbol{\eta}_t$ — 噪声项（输入噪声 $\sim 10\%$-$30\%$、梯度噪声 $\sim \mathcal{N}(0, 0.015^2)$）
\item $\mathcal{K}_t$ — 时变知识库（$\mathcal{K}_{t+1} \supseteq \mathcal{K}_t$，单调解构更新）
\item $\Theta_t$ — 可学习参数（Platt $(A,B)$、LoRA 类偏差、文档权重）
\end{itemize}

\begin{remark}[随机动力系统视角下的 HMAE]
从上述平衡方程可以看出，HMAE 是一个\textbf{受控随机动力系统}。
其关键数学结构是：
\begin{itemize}[nosep]
\item \textbf{状态空间} $\mathcal{S}$：维度极高（KB 索引 $\sim 10^5$ 维 + 嵌入向量 $\sim 128$ 维 + 参数 $\sim 10^4$ 维）
\item \textbf{控制空间} $\mathcal{U}$：离散且有限（$3 \times 3 = 9$ 个选择）
\item \textbf{动态}：混合了确定性更新（$f$）和随机扰动（$\boldsymbol{\xi}$）
\item \textbf{稳定性}：由退化管理和 HITL 保证——当系统状态偏离安全区域时，强制回退到保守策略
\end{itemize}
这是典型的\textbf{混合系统 (hybrid system)}：连续状态 + 离散控制 + 随机噪声。
\end{remark}

\section{物理类比总结}

\begin{table}[H]
\centering
\small
\begin{tabular}{lll}
\toprule
\textbf{HMAE 组件} & \textbf{物理类比} & \textbf{守恒量/不变量} \\
\midrule
SimHash 去重 & 64维超立方体编码 & 汉明距离度量 \\
Softmax $T=0.18$ & 低温极限 — 基态提取 & 自由能极小化 \\
Platt 校准 & 最大熵原则下的 logistic 校正 & 交叉熵一致 \\
RevIN 归一化 & 可逆正则变换 & $p(x;\mu,\sigma)$ 分布保持 \\
TFA 门控融合 & 量子叠加态坍缩 & $\sum g_k = 1$ (概率完备) \\
纳什均衡搜索 & 势博弈单调解 & $\Phi(s,m)$ 单调不减 \\
FAISS 内积搜索 & 高维空间最大投影 & $\cos(\theta)$ 保序 \\
知识蒸馏 & Fr\'{e}chet 均值 & 球面测地线极小化 \\
隐私梯度扰动 & 高斯差分隐私机制 & $(\epsilon, \delta)$-DP 保证 \\
检索成本函数 & 拉格朗日对偶 & $\lambda P - (1-\lambda)C$ 权衡 \\
退化管理 & 工程容错控制理论 & 安全状态可达性 \\
连续学习 & 在线随机梯度下降 & $\mathbb{E}[\mathcal{L}]$ 期望递减 \\
\bottomrule
\end{tabular}
\caption{HMAE 物理类比映射表}
\end{table}

% ============================================================
% PART IV: 输出结果分析
% ============================================================
\part{输出结果分析 —— Output Analysis \& Evaluation Framework}

\vspace{6pt}
\begin{center}
\small
\textit{有了完整的系统架构（Part I--III），我们需要一套严密的评价体系来衡量系统的实际表现。}\\
\textit{本部分从分类指标、消融实验、噪声鲁棒性、数据规模适应性等多个维度}\\
\textit{系统化地定义了 HMAE 的输出质量度量，为实验验证和持续改进提供定量基础。}
\end{center}
\vspace{12pt}

\section{评价指标体系}

\subsection{为什么需要多层评价指标}

HMAE 是一个多模块、多输出的复杂系统。单一指标（如总体准确率）虽然直观，但会掩盖系统在不同维度上的表现差异。一个在常见病上准确率 95\% 但在罕见病上准确率 30\% 的系统，其总体准确率可能仍有 90\%（因为常见病样本远多于罕见病）——这 90\% 的数字是有误导性的。本节建立了一套从微观（每类指标）到宏观（系统级指标）、从静态（单次评估）到动态（噪声鲁棒性、数据规模适应性）的多层评价体系。

\subsection{分类指标}

\subsubsection{基础指标的定义与关系}

混淆矩阵的四个基本元素构成了所有分类指标的基础：
\begin{itemize}[nosep]
\item \textbf{TP (True Positive)：} 系统预测为类别 $c$，真实标签也是类别 $c$ —— 正确命中。
\item \textbf{TN (True Negative)：} 系统预测为非类别 $c$，真实标签也是非类别 $c$ —— 正确排除。
\item \textbf{FP (False Positive)：} 系统预测为类别 $c$，但真实标签不是类别 $c$ —— 误报（假阳性）。
\item \textbf{FN (False Negative)：} 系统预测为非类别 $c$，但真实标签是类别 $c$ —— 漏报（假阴性）。
\end{itemize}
在医疗场景中，FP 和 FN 的代价是不对称的：FN（漏诊）通常比 FP（误诊）更危险——漏诊意味着患者得不到及时治疗，而误诊可以通过进一步的检查来纠正。因此评价体系必须区分这两类错误。

\begin{align}
\text{Accuracy} &= \frac{\text{TP} + \text{TN}}{\text{TP} + \text{TN} + \text{FP} + \text{FN}} \\
\text{Precision}_c &= \frac{\text{TP}_c}{\text{TP}_c + \text{FP}_c} \\
\text{Recall}_c &= \frac{\text{TP}_c}{\text{TP}_c + \text{FN}_c} \\
F1_c &= 2 \cdot \frac{\text{Precision}_c \cdot \text{Recall}_c}{\text{Precision}_c + \text{Recall}_c}
\end{align}

\textbf{公式详解——每个指标的临床含义：}
\begin{itemize}[nosep]
\item \textbf{Accuracy（准确率）：} 所有预测中正确的比例。这是最直观但也是最粗糙的指标——在类别不平衡时（如罕见病仅占 1\% 的样本），一个永远预测"无病"的模型也可以达到 99\% 的准确率。因此准确率必须与以下每类指标配合使用。
\item \textbf{Precision（精确率 / 查准率）：} 在系统预测为类别 $c$ 的所有病例中，真正属于 $c$ 的比例。高 Precision 意味着"系统的诊断是可信的——当它说你是 X 病，你大概率确实是 X 病"。低 Precision 意味着系统经常"乱报"——频繁地将非 X 病误分类为 X 病。临床含义：高 Precision 减少不必要的进一步检查和患者焦虑。
\item \textbf{Recall（召回率 / 查全率 / 敏感度）：} 在所有真正属于类别 $c$ 的病例中，被系统正确识别的比例。高 Recall 意味着"系统不会遗漏 X 病患者"。低 Recall 意味着系统经常"漏诊"——很多 X 病患者被错误地分类为其他疾病或正常。临床含义：高 Recall 减少漏诊，是医疗 AI 最关键的指标之一。
\item \textbf{F1（F1 分数）：} Precision 和 Recall 的调和平均——$2 \cdot \frac{P \cdot R}{P + R}$。调和平均（而非算术平均）的使用是深有用意的：调和平均对最小值更敏感。如果 Precision = 0.9 但 Recall = 0.1，算术平均 = 0.5 但调和平均 = 0.18——F1 低反映了一个维度上的严重缺陷。这符合临床原则：一个 Precision 很高但 Recall 很低的系统（"只报告自己确定的，但漏掉大部分真实病例"）在医疗上是不合格的。
\end{itemize}

\subsubsection{三类平均方式的统计原理}

宏观/微观/加权三种平均方式：

\begin{align}
P_{\text{macro}} &= \frac{1}{|\mathcal{C}|}\sum_{c} P_c, \quad
P_{\text{micro}} = \frac{\sum_c \text{TP}_c}{\sum_c (\text{TP}_c + \text{FP}_c)} \\
F1_{\text{weighted}} &= \sum_c \frac{|\mathcal{D}_c|}{|\mathcal{D}|} \cdot F1_c
\end{align}

\textbf{三种平均方式的适用场景与内在权衡：}
\begin{itemize}[nosep]
\item \textbf{Macro（宏平均）：} 先计算每个类的指标，再对所有类取算术平均。每个类权重相等（$1/|\mathcal{C}|$），与类的样本量无关。\textbf{适用场景：} 关注罕见病时的评估——罕见病类（可能只有 10 个样本）与常见病类（可能有 1000 个样本）在评估中被平等对待。\textbf{局限：} 小样本类的指标方差大（因为分母小），宏平均对随机波动敏感。
\item \textbf{Micro（微平均）：} 将所有类的 TP、FP、FN 分别求和，再计算全局指标。每个样本权重相等（$1/|\mathcal{D}|$），与样本所属类别无关。\textbf{适用场景：} 评估整体系统表现时——"在系统处理的所有病例中，有多少被正确处理了？" \textbf{数学性质：} 微平均 Precision = 微平均 Recall = Accuracy（当每个样本只有一个真实标签时）。这不是巧合——它意味着微平均实际上衡量的是总体准确率。
\item \textbf{Weighted（加权平均）：} 先计算每个类的指标，再以类样本量为权重取加权平均。大类（样本多的类）在评估中占更大权重。\textbf{适用场景：} 关注总体临床效果时——"考虑到实际就诊人群的疾病分布（常见病就诊量远大于罕见病），系统的整体表现如何？" \textbf{与宏平均的对比：} 如果系统在常见病上表现优秀（F1=0.95）但在罕见病上表现差（F1=0.30），加权 F1 可能为 0.88（因为常见病权重高），宏 F1 为 0.625——两个数字反映了完全不同的系统图景。选择哪个指标取决于评估的目的：是评估"对每个患者"（加权）还是"对每种疾病"（宏）。
\end{itemize}

\textbf{三种平均方式的总结：} 在 HMAE 的评估框架中，我们同时报告三种平均指标，而非仅依赖其中一种。这确保系统的表现从多个角度被全面审视——加权 F1 告诉我们系统在真实就诊人群中的预测质量，宏 F1 告诉我们系统是否公平地对待常见病和罕见病。

\subsection{消融实验指标}

\begin{equation}
\Delta_{\text{module}} = \text{Acc}(\text{full}) - \text{Acc}(\text{without\_module})
\end{equation}

\textbf{公式详解：} 消融实验（ablation study）是评估每个模块"边际贡献"的标准方法。$\Delta_{\text{module}}$ 测量的是移除某个模块后系统精度的下降量——下降越大，该模块越重要。正值的 $\Delta$ 意味着该模块是有效贡献的（移除后精度下降），零或负值意味着该模块无效或有害（移除后精度不变或上升——可能模块引入了噪声）。

\textbf{实验设计：} 消融实验的条件矩阵（见 §23）包含 $2^3 = 8$ 种组合（HITL 开/关 × TFA 开/关 × Multimodal 开/关）。例如：
\begin{itemize}[nosep]
\item $\Delta_{\text{HITL}} = \text{Acc}(\text{HITL 开, TFA 开, Multimodal 开}) - \text{Acc}(\text{HITL 关, TFA 开, Multimodal 开})$：HITL 在完整系统（TFA 和 Multimodal 均在）中的边际贡献。
\item 注意边际贡献不是独立的——HITL 在 TFA 打开时的贡献可能不同于 HITL 在 TFA 关闭时的贡献（模块间存在交互效应）。完整的 $2^3$ 析因设计可以估计所有主效应和交互效应。
\end{itemize}

\subsection{噪声鲁棒性指标}

\begin{equation}
R(\sigma) = \frac{\text{Acc}(\text{noise} = \sigma)}{\text{Acc}(\text{noise} = 0)}, \quad \sigma \in \{0.10, 0.30\}
\end{equation}

\textbf{公式详解：} $R(\sigma)$ 是相对鲁棒性比率——在输入中添加标准差为 $\sigma$ 的高斯噪声后，系统的精度保留比例。$R(\sigma) = 1.0$ 意味着系统完全不受该水平噪声的影响（理想状态）；$R(\sigma) = 0.5$ 意味着噪声使精度减半（系统对该水平噪声非常敏感）。

\textbf{噪声水平的临床对应：} $\sigma = 0.10$ 对应临床数据中的轻度噪声（如轻微的测量误差、录入笔误）；$\sigma = 0.30$ 对应重度噪声（如缺失值被不合理地填充、严重的数据录入错误、不同医院的不同测量标准导致的系统性偏差）。

\textbf{评估目标：} 在真实临床环境中，数据不可能是完美的。一个在干净测试集上表现好但在噪声数据下迅速退化的系统，在实际部署中是不可靠的。$R(0.30) > 0.85$（在重度噪声下仍保留 85\% 的精度）是 HMAE 的工程目标。

\subsection{数据规模适应性指标}

\begin{equation}
\eta_{\text{scale}}(f) = \frac{\text{F1}(f \cdot |\mathcal{D}_{\text{train}}|)}{\text{F1}(|\mathcal{D}_{\text{train}}|)}, \quad f \in \{0.20, 0.50, 0.80, 1.00\}
\end{equation}

\textbf{公式详解：} $\eta_{\text{scale}}(f)$ 以训练数据的一个比例 $f$（如 20\%）训练系统，然后将 F1 得分与全量数据（100\%）训练的 F1 比较。$\eta_{\text{scale}}(0.20) = 0.70$ 意味着用 20\% 的数据就能达到全量数据 70\% 的 F1——这是非常高效的学习曲线（数据效率高）。$\eta_{\text{scale}}(0.20) = 0.30$ 意味着 20\% 的数据只能达到 30\% 的 F1——系统严重依赖大数据量。

\textbf{评估目标：} 数据规模适应性指标回答一个实用问题："系统在数据稀缺的科室（如罕见病专科，只有少量历史病例）中能表现多好？" 四种比例 $f \in \{0.20, 0.50, 0.80, 1.00\}$ 覆盖了从"极度稀缺"到"充足"的数据量范围。

\subsection{常见/罕见病差异指标}

\begin{equation}
\Delta_{\text{rare}} = \frac{\text{F1}_{\text{common}} - \text{F1}_{\text{rare}}}{\text{F1}_{\text{common}}}
\end{equation}

\textbf{公式详解：} $\Delta_{\text{rare}}$ 测量常见病与罕见病之间的 F1 相对差异。$\Delta_{\text{rare}} = 0$ 意味着系统在常见病和罕见病上表现完全一样（公平的）；$\Delta_{\text{rare}} = 0.50$ 意味着罕见病的 F1 只有常见病的一半（显著的长尾偏差）。

\textbf{阈值 0.30 的含义：} $\Delta_{\text{rare}} > 0.30$ 指示罕见病检索覆盖不足或训练数据长尾效应——系统在罕见病上的表现显著差于常见病，需要改进罕见病的知识库覆盖或应用数据增强/重采样技术。0.30 是一个经验阈值——低于此值，罕见病与常见病的差距被认为在临床可接受范围内。

\textbf{公平性维度的意义：} 这个指标反映了 AI 伦理中的\textbf{公平性（fairness）}维度——AI 系统不应因为某些疾病"少见"而对这些疾病的患者提供更差的诊断服务。罕见病患者往往是最需要 AI 辅助的群体（因为人类医生对这些疾病也没有太多经验），如果 AI 系统在罕见病上表现更差，那就背离了技术的初衷。

\section{系统级输出分析}

\subsection{TFA 风险预测输出结构}

每次 TFA 预测产生三个时间窗口的风险概率向量：
\begin{equation}
\mathbf{r}_{\text{TFA}} = [r_s, r_m, r_l]^{\top} \in [0,1]^{3}
\end{equation}

其联合分布定义了一个三维风险流形 $\mathcal{M}_{\text{risk}} \subset [0,1]^3$。对角元素对应"恶化-稳定"二项分布：$P(\text{deteriorate}) = r_s$, $P(\text{stable}) = 1 - r_s$。

\subsection{共识决策输出结构}

共识模块产生分层决策向量：
\begin{equation}
\mathbf{d}_{\text{consensus}} = [\text{approved}, \text{tier}, \text{votes}, \text{risk\_level}, \text{escalate\_to\_human}]
\end{equation}

其中 $\text{tier} \in \{\text{force\_hitl}, \text{dma\_veto}, \text{auto\_pass}, \text{auto\_pass\_risk\_noted}, \text{escalation\_risk\_uncertainty}\}$。

\subsection{系统综合答案 (Comprehensive Answer)}

融合模块输出统一的临床答案结构：
\begin{itemize}[nosep]
\item \texttt{diagnosis}: $\{\text{prediction}, \text{confidence}, \text{severity}, \text{intent}, \text{status}\}$
\item \texttt{risk\_assessment}: $\{\text{risk\_level}, \text{risk\_score}, [r_s, r_m, r_l], \text{factors}, \text{reasoning}\}$
\item \texttt{recommendations}: 分层临床建议（基于 TFA 风险等级）
\item \texttt{disclaimer}: 医疗免责声明
\end{itemize}

\section{消融实验分析框架}

\begin{table}[H]
\centering
\small
\begin{tabular}{lccc}
\toprule
\textbf{条件} & \textbf{HITL} & \textbf{TFA} & \textbf{Multimodal} \\
\midrule
完整生态系统 & $\checkmark$ & $\checkmark$ & $\checkmark$ \\
无 HITL & $\times$ & $\checkmark$ & $\checkmark$ \\
无 TFA & $\checkmark$ & $\times$ & $\checkmark$ \\
无 Multimodal & $\checkmark$ & $\checkmark$ & $\times$ \\
\bottomrule
\end{tabular}
\caption{消融实验条件矩阵 ($2^3 = 8$ 种组合)}
\end{table}

各条件的准确度差异 $\Delta_{\text{acc}}$ 量化了每个模块的信息增益。

\section{持续学习收敛分析}

连续学习循环的精度趋势估计：

\begin{equation}
\hat{a}_k = \frac{1}{|\mathcal{F}_k|}\sum_{(\bar{p}_i, y_i) \in \mathcal{F}_k} \mathbb{I}[\text{corrected\_label}_i = \text{predicted\_label}_i]
\end{equation}

Platt 参数的收敛路径：
\begin{equation}
\begin{bmatrix} A_k \\ B_k \end{bmatrix} = \begin{bmatrix} A_{k-1} \\ B_{k-1} \end{bmatrix} - \eta \cdot \frac{1}{|\mathcal{F}_k|}\sum_i \begin{bmatrix} (p_i - y_i) \cdot \bar{p}_i \\ (p_i - y_i) \end{bmatrix}
\end{equation}

当 $|\mathbb{E}[p_i - y_i]| < \epsilon$ 时系统达到校准平衡，此时 $\gamma_{\text{cal}} \approx \mathbb{E}[y_i \mid \bar{p}_i]$。

\section{退化水平对输出的影响}

\begin{table}[H]
\centering
\small
\begin{tabular}{llll}
\toprule
\textbf{退化等级} & \textbf{DMA 后端} & \textbf{TFA 后端} & \textbf{预期精度影响} \\
\midrule
L0 (NORMAL) & Ollama Llama3:8b & MedTsLLM-v1.5 & 基准 \\
L1 (MINOR) & Ollama + 4-bit quant & MedTsLLM-v1.5 & $\approx -2\%$ \\
L2 (PARTIAL) & Ollama Llama3:8b & Heuristic fallback & $\approx -5\%$ — TFA 风险精度降 \\
L3 (SEVERE) & Legacy fallback & Heuristic fallback & $\approx -10\%$ — 整体性能降 \\
\bottomrule
\end{tabular}
\caption{退化等级-精度影响映射}
\end{table}

\section{系统全面报告 (Full Report)}

\texttt{get\_full\_report()} 输出的完整生态系统快照包含：
\begin{itemize}[nosep]
\item \textbf{数据集统计}: 样本数、类别分布、训练/测试划分
\item \textbf{任务统计}: 历史任务数、博弈轮次、近期预测
\item \textbf{智能体状态}: 10个智能体的调用次数、故障数、运行时统计
\item \textbf{实验指标}: 消融 $\times$ 4条件、鲁棒性 $\times$ 3噪声、适应性 $\times$ 4比例、常见/罕见对比
\item \textbf{退化状态}: 组件状态、退化等级、回退次数
\item \textbf{持续学习}: 循环次数、缓冲反馈、近期指标趋势
\end{itemize}

% ============================================================
% ============================================================
% PART V (NEW): AXIOMATIC SYSTEM
% ============================================================
% ============================================================
\part{公理体系推演 —— Axiomatic Foundations of HMAE}

\vspace{6pt}
\begin{center}
\small
\textit{本部分是本文的核心理论创新。我们从 HMAE 的工程实现和数学结构中}\\
\textit{逆向归纳出 8 条公理，构成本系统的公理化基础。}\\
\textit{这些公理不是假设的——每一条都直接来源于前四部分中推导的数学公式。}
\end{center}
\vspace{12pt}

% ------------------------------------------------------------
\section{从工程实现到公理体系的逆向归纳}

\subsection{为什么需要公理化}

前四个部分详细描述了 HMAE 的数学公式和物理类比，但它们仍然是\textbf{描述性的}
（descriptive）——告诉我们系统"如何运作"。本部分的目标是建立\textbf{规范性的}
（normative）基础——推导系统"必须遵循什么原则"。

公理化有四个核心价值：
\begin{enumerate}[nosep]
\item \textbf{完备性检验：} 公理体系揭示系统设计是否自洽，是否存在逻辑漏洞
\item \textbf{可扩展性指导：} 新模块的设计必须满足已有公理，否则会破坏系统一致性
\item \textbf{性质证明：} 从公理出发，可以推导出系统行为的必然性质（定理），而非仅靠实验观察
\item \textbf{跨领域迁移：} 公理体系抽象掉了医疗领域的具体细节，使得系统架构可以迁移到其他高风险决策领域（金融、法律、航空）
\end{enumerate}

\subsection{归纳方法论}

公理的提炼遵循以下方法：
\begin{enumerate}[nosep]
\item 从前四部分中识别\textbf{不变量}（物理类比表中的第三列）
\item 检查每个不变量的\textbf{适用范围}——它是否在所有运行条件下成立
\item 将不变量提升为公理，并明确其\textbf{形式化表达}
\item 验证公理之间的\textbf{独立性}——每条公理不能从其他公理推导出来
\end{enumerate}

% ------------------------------------------------------------
\section{八大公理}

% --- Axiom 1 ---
\begin{axiom}[感知完备性 —— Perception Completeness]
\label{ax:perception}
对任意多模态输入 $(X_{\text{txt}}, X_{\text{img}})$，存在一个 L2 归一化的融合映射
$\Phi: \mathcal{X}_{\text{txt}} \times \mathcal{X}_{\text{img}} \to S^{d-1}$，
使得对于任何与下游任务相关的信息分量 $\mathcal{I}$，
当 $\mathcal{I}$ 在原始信号中可检测时，$\Phi$ 的输出保留 $\mathcal{I}$ 的线性可分性。

即：若 $\exists$ 分类器 $f$ 使 $f(\mathcal{I}(X)) = y$，则
$\exists$ 线性分类器 $g$ 使 $g(\Phi(X)) = y$，且 $\|g\|_{\text{Lip}} \leq \frac{1}{\sqrt{d}}$。
\end{axiom}

\subsubsection{公理1的详细解释}

\textbf{动机与问题背景：} HMAE 系统的输入是极端异质的——临床文本（自然语言，变长序列）与医学影像（像素矩阵，固定分辨率）在原始空间中没有任何可比性。如果直接在这两个异构空间中进行下游任务（如疾病分类），系统需要同时学习两个完全不同的特征提取器，这在实际中既低效又不可靠。感知完备性公理断言：存在一个\textbf{统一的、低维的、可线性分离的表征空间} $S^{d-1}$（$d=128$ 维单位超球面），使得所有原始信号中"对任务有用的信息"在降维后仍然可被线性分类器恢复。这一定理级别的保证是 HMAE 所有下游模块（检索、决策、预测）能够正常工作的前提——如果感知层丢失了关键信息，后续任何精妙的算法都无法弥补。

\textbf{形式化解读：} 公理中的两个条件分别刻画了"保信息性"和"稳定性"：
\begin{itemize}[nosep]
\item \textbf{保信息性}（$\exists g$ 使 $g(\Phi(X)) = y$）：原始空间中能被任意分类器识别的模式，在投影后仍能被某个\textbf{线性}分类器识别。这里"线性"是关键——它意味着投影后的表征足够"干净"，信息不是以复杂的非线性方式编码的，而是以简单的线性方向编码。这使得下游的余弦相似度搜索（FAISS）和线性分类（DMA 质心分类）在数学上是充分的。
\item \textbf{稳定性约束}（$\|g\|_{\text{Lip}} \leq 1/\sqrt{d}$）：Lipschitz 常数上界保证了投影映射不会"放大"输入中的微小扰动。在 $d=128$ 时，$1/\sqrt{128} \approx 0.088$，这意味着输入中 $\delta$ 大小的变化在输出端最多被放大为 $0.088\delta$。这一约束是系统鲁棒性的数学基础——临床数据中的自然噪声（如打字错误、图像采集差异）不会在语义空间中产生不成比例的影响。
\end{itemize}

\textbf{逐层推导过程：} 公理1的形式化陈述是从感知智能体的三个具体公式中逐步提炼出来的：
\begin{enumerate}[nosep]
\item \textbf{文本编码（式1-2）：} $\mathbf{v}_{\text{txt}}^{(d)} = \Ltwo(\text{Resize}(\text{MiniLM}(X_{\text{txt}}), d))$。MiniLM 已在超过 10 亿句子对上预训练，其输出的 384 维向量已被证明在语义相似度任务上保持了余弦相似度的保序性。分块均值降维（从 384 到 128）在极限意义上等价于对 384 维向量进行随机投影——Johnson-Lindenstrauss 引理保证成对距离以高概率保持。L2 归一化消除了幅值维度，使得所有信息仅由方向承载。
\item \textbf{图像编码（式3-4）：} ViT 风格的分块编码将图像分割为 4 个区域，每个区域通过 TextProxy 转换为自然语言描述后经 MiniLM 编码。4 分块的划分在分辨率（更多块 → 更精细的空间信息）和计算效率（更少块 → 更快的编码）之间取得平衡。均值池化提供了基本的平移不变性。
\item \textbf{交叉注意力融合（式5-6）：} 这是保信息性的核心机制。设文本语义向量 $\mathbf{v}_{\text{txt}}$ 为目标信号，图像块 $\mathbf{P}$ 为辅助信号。交叉注意力通过 $\softmax(\mathbf{P}\mathbf{q}^{\top}/\sqrt{d})$ 计算出每个图像块对文本语义的"相关性权重"——与文本语义一致的图像块获得高权重，不一致的获得低权重。最终融合向量 $\mathbf{v}_{\text{fused}} \in S^{d-1}$ 是这两个信号的"共识投影"——它只保留了文本和图像共同支持的信息分量，过滤了单方噪声。
\end{enumerate}
将这三个步骤组合起来，整个映射 $\Phi = \text{L2} \circ \text{CrossAttn} \circ [\text{MiniLM}, \text{ViT-Proxy}]$ 构成一个从异构输入空间到单位球面的连续映射，其 Lipschitz 常数由上界 $\|\mathbf{W}_{\text{fusion}}\|_{\text{spec}} / \sqrt{d} \approx 1/\sqrt{d}$ 控制。

\textbf{在 HMAE 系统中的具体体现：} 公理1的约束直接影响以下设计选择：
\begin{itemize}[nosep]
\item $d=128$ 是经过实验验证的平衡点——低于 64 维时信息损失导致检索精度显著下降（信息瓶颈过紧），高于 256 维时 FAISS 索引的构建和搜索成本快速增加而边际收益递减。
\item 交叉注意力的缩放因子 $1/\sqrt{d}$ 直接来源于 $\mathbf{W}_{\text{fusion}}$ 的初始化方案 $w_{ij} \sim \mathcal{N}(0, 1/\sqrt{d})$——这是 Xavier/Glorot 初始化的标准选择，使得前向传播时各层的激活方差保持一致。
\item 当图文不一致时（文本描述"肺炎"但图像正常），$\mathbf{v}_{\text{fused}}$ 不会崩溃为零向量（L2 归一化强制在球面上），而是指向文本与图像语义的"折中方向"——这一方向与两个原始信号都有较大夹角，自然地表达了不确定性。
\end{itemize}

\textbf{物理意义：} 这是"信息不灭"在有限容量系统中的弱化版本——我们无法保留所有信息
（维度从 $\sim 10^5$ 降至 128），但保证保留了所有线性可分的任务相关信息。
这类似于量子力学中的\textbf{无克隆定理}的逆命题：虽然我们无法完美复制量子态的所有信息，
但可以提取对于特定测量（任务）最优的表征。更进一步，$\Phi$ 可以理解为一个\textbf{量子信道 (quantum channel)}——它将高维的"密度矩阵"（原始信号）通过一个完全正定且保迹的映射压缩为低维的"有效态"。L2 归一化对应态的归一化条件 $\text{Tr}(\rho) = 1$，交叉注意力对应投影测量。

\textbf{与其他公理的逻辑关系：}
\begin{itemize}[nosep]
\item \textbf{公理2（检索不确定性）依赖公理1}：检索系统（RAA）在 $\Phi(X)$ 的表征空间中进行语义搜索——如果 $\Phi$ 不完备，任何检索策略都无法弥补信息丢失。公理1是公理2的信息论前提。
\item \textbf{公理3（低温收敛）依赖公理1}：DMA 的 softmax 分类器以 $\Phi(X)$ 为输入——如果 $\Phi(X)$ 不能保持线性可分性，则即使 $T \to 0$ 也无法收敛到正确类别。公理1是公理3的几何前提。
\item \textbf{公理7（退化安全）保护公理1}：当 MiniLM 或交叉注意力模块故障时，退化管理器确保系统回退到纯文本 TF-IDF 表征（绕过 $\Phi$ 的完整映射），保证公理1在降级模式下以更弱的形式（仅文本）仍然成立。
\end{itemize}

% --- Axiom 2 ---
\begin{axiom}[检索不确定性与成本权衡 —— Retrieval Uncertainty-Cost Trade-off]
\label{ax:retrieval}
对任意查询 $q$，检索系统的最优策略 $(s^*, m^*)$ 满足：
\begin{equation}
(s^*, m^*) = \arg\max_{s,m} \left[\lambda \cdot \mathbb{E}[P(\text{correct})] - (1-\lambda) \cdot C(s, n)\right]
\end{equation}
且存在信息论下界，具体形式依赖于检索系统的信息率失真函数，此处不展开。
其中 $H(Y \mid \cdot)$ 为给定知识状态下的条件熵，$C_{\max}$ 为最大可行成本。
\end{axiom}

\subsubsection{公理2的详细解释}

\textbf{动机与问题背景：} 信息检索不是免费的。每一次文献检索消耗计算资源（CPU 周期、内存、I/O），在医疗实时场景中，检索延迟直接影响临床决策的时效性。但另一方面，检索不足（文档太少或质量太低）会导致诊断信息不完整，增加误诊风险。公理2将这一根本性权衡形式化：\textbf{检索系统的设计目标不是"检索尽可能多的文档"，而是"在精度提升的边际收益与检索成本之间找到最优平衡点"。}

这一权衡在经济学中对应\textbf{边际效用递减原理}——第 1 篇检索到的相关文档带来的信息增益远大于第 100 篇。在信息论中，这对应\textbf{率失真理论 (rate-distortion theory)}——给定失真容忍度 $D$，存在一个最小信息率 $R(D)$，低于此率则无法实现该失真水平。HMAE 的检索系统以成本 $C(s,n)$ 作为信息率的代理，以 $\mathbb{E}[P(\text{correct})]$ 作为失真补（正确率 = $1 - \text{失真}$），从而在率失真框架下运行。

\textbf{形式化解读——优化目标的逐项分析：}
\begin{itemize}[nosep]
\item \textbf{精度项} $\lambda \cdot \mathbb{E}[P(\text{correct})]$：系统从检索知识库 $\mathcal{K}_s$ 中获得正确诊断的期望概率。$\lambda=0.72$ 是精度权重——它表明在医疗场景中，正确率的重要性是计算成本的约 2.6 倍（$0.72/0.28 \approx 2.57$）。这一定量关系直接编码了"宁可多花计算也要保证准确性"的医学伦理先验。
\item \textbf{成本项} $-(1-\lambda) \cdot C(s, n)$：负号表明成本降低效用。$C(s,n)$ 的具体形式在式（§11.2）中给出——三种策略（mixed/rerank/adaptive）有不同的基础成本（0.22/0.36/0.28）加上与检索文档数 $n$ 成比例的边际成本（$0.012n$）。adaptive 策略还有额外的 $0.08 \cdot \mathbb{I}[U(q)>0.58]$ 项——当查询不确定性高时，激活深层检索产生额外成本。
\item \textbf{信息论下界} $H(Y \mid q, \mathcal{K}_s) \cdot C_{\max}$：这是公理化约束的核心——它断言\textbf{不确定性不能免费消除}。查询 $q$ 的目标变量 $Y$（正确疾病类别）在给定知识 $\mathcal{K}_s$ 后的条件熵 $H(Y \mid q, \mathcal{K}_s)$ 衡量的正是"检索系统仍然不知道的信息量"。将这一剩余不确定性与最大可行成本 $C_{\max}$ 相乘，得到一个下界——系统不可能以低于此下界的成本消除不确定性。这类似兰道尔原理在计算热力学中的角色：擦除信息需要消耗最小能量。
\end{itemize}

\textbf{逐层推导过程：} 公理2的每一个组成部分都直接来自前文的数学公式：
\begin{enumerate}[nosep]
\item \textbf{不确定性 $U(q)$ 的来源（式 \ref{eq:uncertainty}）：} RAA 的三重不确定性量化——归一化熵 $H(P)$（得分分布的平坦度）、Top-2 间隙 $G(P)$（最佳与次佳的分离度）、OOV 代理 $O(q)$（极短查询标记）。这三项以权重 $[0.58, 0.32, 0.10]$ 线性组合为 $U(q)$，当 $U(q) > 0.58$ 时触发深层检索。
\item \textbf{检索策略空间 $\mathcal{S}$（§11.1）：} $\mathcal{S} = \{\text{mixed}, \text{rerank}, \text{adaptive}\}$ 是三个离散策略。mixed 同时运行 TF-IDF 和 FAISS 后融合，rerank 在 mixed 基础上再排序，adaptive 根据查询不确定性动态选择。每个策略的成本函数在式（§11.2）中给出。
\item \textbf{DMA 模式空间 $\mathcal{M}$（§11.1）：} $\mathcal{M} = \{\text{fast}, \text{balanced}, \text{deep}\}$ 定义 DMA 处理检索结果的深度。fast 仅使用 top-3 文档快速决策，balanced 使用 top-5，deep 使用 top-10 加多样化策略。
\item \textbf{纳什均衡收敛（§11.2-11.3）：} 最佳响应迭代 $s^{(t+1)} = \arg\max_s U(s, m^{(t)})$ 和 $m^{(t+1)} = \arg\max_m U(s^{(t+1)}, m)$ 保证在势函数 $\Phi(s,m)$ 单调不减的条件下有限步收敛。$N_{\max}=6$ 是经验上限——在 9 个纯策略组合中，最多 6 轮迭代即可到达均衡。
\end{enumerate}

\textbf{在 HMAE 系统中的具体体现：}
\begin{itemize}[nosep]
\item $\lambda=0.72$ 的设定直接影响了 RAA 的行为——系统宁愿使用更昂贵的 rerank 策略（基础成本 0.36）也不愿冒信息不足的风险使用 mixed 策略（基础成本 0.22），除非查询非常简单且不确定性极低。
\item $C_{\max}$ 在系统中没有显式定义但隐含在 6 轮迭代限制中——每轮迭代消耗约 0.2--0.5 秒的计算时间，6 轮总计约 1.2--3 秒，这在实时临床场景中是可接受的延迟上限。
\item 信息论下界在实践中的体现是：当 $U(q) > 0.58$ 时，系统自动从 $k=5$ 扩展到 $k=10$ 的检索深度——这是"用更多文档（更高的信息率）换取更低的剩余不确定性"的具体操作。
\end{itemize}

\textbf{物理意义：} 这是\textbf{兰道尔原理 (Landauer's principle)} 在信息检索中的类比。
兰道尔原理指出：擦除1比特信息至少消耗 $kT\ln 2$ 的能量。
类似地，我们的下界表明：减少 $H(Y \mid q) - H(Y \mid q, \mathcal{K})$
的不确定性至少需要付出与检索文档数成比例的计算成本。
信息不是免费的——这一热力学约束在 AI 系统中同样成立。

\textbf{与其他公理的逻辑关系：}
\begin{itemize}[nosep]
\item \textbf{公理1（感知完备性）是前提}：检索在 $\Phi(X)$ 的表征空间中进行——如果 $\Phi$ 不完备，检索到的文档再精确也无法弥补表征层的信息丢失。
\item \textbf{公理5（共识收敛）共享优化框架}：两个公理都使用效用最大化（期待效用/纳什乘积）作为决策准则，且都收敛到均衡状态。公理2 在检索层面，公理5 在决策层面。
\item \textbf{公理8（持续学习）反馈驱动}：HITL 反馈中的 $E_{\text{missing}}$ 错误（检索覆盖不足）直接触发检索权重的对比更新（§17），从而调整 $C(s,n)$ 中各类文档的隐式权重。
\end{itemize}

% --- Axiom 3 ---
\begin{axiom}[决策分布的低温收敛 —— Low-Temperature Decision Convergence]
\label{ax:decision}
设 DMA 的原始得分为 $\mathbf{s} \in \mathbb{R}^{|\mathcal{C}|}$，softmax 温度为 $T$。
则当 $T \to 0^+$ 时，决策分布收敛于决定性输出：
\begin{equation}
\lim_{T \to 0^+} P(c \mid q) = \mathbb{I}[c = \arg\max_{c'} s_{c'}]
\end{equation}
且对任意 $\delta > 0$，存在 $T_0 > 0$ 使得当 $T < T_0$ 时：
\begin{equation}
\max_c P(c \mid q) \geq 1 - \delta
\end{equation}
\end{axiom}

\subsubsection{公理3的详细解释}

\textbf{动机与问题背景：} 在临床诊断中，当证据充分且一致时，系统应当给出确定的、高置信度的诊断——模棱两可的概率分布在医疗场景中不构成可操作的输出（"患者有 35\% 的概率患肺炎，30\% 的概率患支气管炎，35\% 的概率正常"对医生没有帮助）。另一方面，当证据不足或存在歧义时，系统应当诚实地表达不确定性，而不是强行给出虚假的高置信度。公理3在数学上形式化了这一需求：通过 softmax 温度参数 $T$ 在"确定性"和"不确定性表达"之间进行精确的连续调节。

\textbf{形式化解读——两个数学命题的含义：}
\begin{itemize}[nosep]
\item \textbf{命题1（$T \to 0^+$ 极限）：} $\lim_{T \to 0^+} P(c \mid q) = \mathbb{I}[c = \arg\max_{c'} s_{c'}]$ 表明当温度趋近于零时，softmax 输出退化为一热编码（one-hot）——概率质量完全集中在得分最高的类别上。这个极限行为的证明直接来自 softmax 定义：$P(c) = \exp(s_c/T) / \sum_{c'} \exp(s_{c'}/T)$。设 $c^* = \arg\max s_c$，则 $P(c^*) = 1 / (1 + \sum_{c' \neq c^*} \exp((s_{c'} - s_{c^*})/T))$。当 $T \to 0^+$ 时，所有 $c' \neq c^*$ 的指数项 $\exp((s_{c'} - s_{c^*})/T) \to 0$（因为 $s_{c'} - s_{c^*} < 0$），因此 $P(c^*) \to 1$。注意命题只保证 \textbf{极限行为}——在实际的 $T=0.18$ 下，系统并非完全确定。
\item \textbf{命题2（有限温度的 $\delta$ 近似）：} $\max_c P(c \mid q) \geq 1 - \delta$ 是一个更实用的陈述——它不要求 $T=0$，而是断言：对于任意希望的"确定性程度" $1-\delta$，存在一个足够低的温度 $T_0$ 使系统输出至少达到这一确定性。反过来，如果我们固定温度 $T$，则命题给出了系统在该温度下可能达到的最大"不确定度" $\delta_{\min}(T)$。对于 $T=0.18$，实验观察到 $\max_c P(c) \approx 0.85^+$，即 $\delta \approx 0.15$——系统在最自信时仍有约 15\% 的概率质量分散在其他类别上。
\end{itemize}

\textbf{温度 $T$ 的深层含义——从统计力学到决策理论：}
在统计力学中，正则系综（canonical ensemble）中系统处于能量为 $E_i$ 的状态的概率为 $P(i) = \exp(-E_i/kT) / \sum_j \exp(-E_j/kT)$。将得分 $s_c$ 视为负能量 $-E_c$，softmax 函数与玻尔兹曼分布完全同构：
\[
P(c) = \frac{\exp(s_c/T)}{\sum_{c'} \exp(s_{c'}/T)} \equiv \frac{\exp(-E_c/T)}{\sum_{c'} \exp(-E_{c'}/T)}, \quad E_c \equiv -s_c
\]
在这个同构下：
\begin{itemize}[nosep]
\item \textbf{温度 $T$ 的角色：} $T$ 控制分布的"热涨落"程度。高 $T$ → 分布均匀（高熵，系统在各状态间随机涨落）；低 $T$ → 分布尖锐（低熵，系统冻结在最低能量状态附近）。
\item \textbf{自由能解释：} 系统的亥姆霍兹自由能 $F = -T \log Z$，其中配分函数 $Z = \sum_{c} \exp(s_c/T)$。最小化自由能等价于在"低能量"（高分）和"高熵"（均匀分布）之间取得平衡——这恰好是确定性决策与不确定性表达的权衡。
\item \textbf{$T=0.18$ 的物理解释：} 在室温（~300K）的尺度下，$T=0.18$ 是一个极低的温度——它意味着得分差 $\Delta s = 1$ 被放大了 $1/0.18 \approx 5.6$ 倍。两个得分相差 1 的类别，在 $T=0.18$ 下的概率比约为 $\exp(5.6) \approx 270:1$。这确保了类间只要存在可观的得分差，胜出类别就会获得压倒性的概率质量。
\end{itemize}

\textbf{在 HMAE 系统中的具体体现：}
\begin{itemize}[nosep]
\item DMA 的 softmax 输出（式§3.4）直接使用 $T=0.18$。这一定值经过了严格的实验验证——$T$ 太大（如 $T=1.0$，标准 softmax）导致概率分布过于平坦（多个类别的概率接近），无法做出有意义的单一诊断；$T$ 太小（如 $T=0.05$）导致概率过度集中在最高分类别，丢失了有用的"次可能诊断"信息。
\item Platt 校准（式§3.5）进一步修正了 $T=0.18$ 可能带来的"过度自信"——即使 softmax 输出的 $\max_c P(c) = 0.92$，经过 Platt 校准后可能被下调至 0.78（如果历史数据显示系统在该置信度水平上实际正确率只有 78\%）。
\item HITL 的触发阈值 $\tau_{\text{conf}}=0.62$（§14.1）与公理3密切相关——$\tau_{\text{conf}}$ 定义了一个"置信度不足"的阈值，当 $\max_c P(c) < 0.62$ 时（即 $\delta > 0.38$），系统不确定性太大，触发人工干预。
\end{itemize}

\textbf{物理意义：} 这是\textbf{热力学第三定律}在决策理论中的类比——
当温度趋于绝对零度时，熵趋于零（系统进入纯态）。
在 DMA 中，$T \to 0$ 意味着决策分布的熵趋于零——系统给出完全确定的输出。
有限温度 $T = 0.18$ 代表一个"低温但非绝对零度"的操作点：系统高度确定但保留了一小部分
不确定性（熵 $H \approx 0.2$--$0.5$ bits），用于后续校准和学习。

\textbf{与其他公理的逻辑关系：}
\begin{itemize}[nosep]
\item \textbf{公理1（感知完备性）是输入保证}：只有当 $\Phi(X)$ 在线性可分意义上"保信息"时，得分向量 $\mathbf{s}$ 才能可靠地反映真实的类间差异——如果感知层丢失了类别区分信息，再低的温度也无法收敛到正确类别。
\item \textbf{公理5（共识收敛）以公理3为基础}：共识模块的投票输入之一就是 DMA 的 softmax 输出概率。三层级决策中的 AUTO\_PASS 条件（$\gamma \geq 0.8$）直接依赖于公理3保证的高置信度分布的可靠性。
\item \textbf{公理8（持续学习）通过 Platt 校准调节}：持续学习循环中的 Platt 参数 $(A,B)$ 更新本质上是在调整 softmax 后验概率到实际正确率的校准映射，补充了公理3中有限温度可能带来的校准误差。
\end{itemize}

% --- Axiom 4 ---
\begin{axiom}[时序预测的马尔可夫因果性 —— Temporal Markov Causality]
\label{ax:temporal}
TFA 的多窗口风险预测 $\mathbf{r} = [r_s, r_m, r_l]^{\top}$ 满足因果时序约束：
\begin{equation}
\forall t_1 < t_2 < t_3:\; r_{\text{short}}(t_3) \text{ 在给定 } r_{\text{mid}}(t_2) \text{ 的条件下与 } r_{\text{long}}(t_1) \text{ 条件独立}
\end{equation}
即风险演化过程构成\textbf{非齐次马尔可夫链}：
\begin{equation}
P(r_{\text{short}} \mid r_{\text{mid}}, r_{\text{long}}) = P(r_{\text{short}} \mid r_{\text{mid}})
\end{equation}
\end{axiom}

\subsubsection{公理4的详细解释}

\textbf{动机与问题背景：} 临床风险在不同时间尺度上以不同的动力学演化。一个患者的长期风险（12 个月内发生心衰的概率）由慢性因素（年龄、糖尿病史、高血压控制情况）主导；中期风险（30 天内恶化概率）受近期趋势（血压波动、用药依从性）影响；短期风险（24 小时内急性事件概率）主要由即时信号（当前的感染指标、跌倒事件、生命体征骤变）决定。公理4将这一临床直觉形式化为严格的数学结构：\textbf{三层风险预测构成一个因果链，信息沿时间箭头单向流动——过去的信息通过中间层间接影响现在，而不能跳过中间层直接作用于当前状态。}

\textbf{形式化解读——条件独立性与马尔可夫性的含义：}
\begin{itemize}[nosep]
\item \textbf{条件独立性：} $r_{\text{short}}(t_3) \perp\!\!\!\perp r_{\text{long}}(t_1) \mid r_{\text{mid}}(t_2)$。在概率论的语言中，这意味着一旦我们知道了中期风险 $r_{\text{mid}}(t_2)$，长期风险 $r_{\text{long}}(t_1)$ 不再提供关于短期风险 $r_{\text{short}}(t_3)$ 的任何额外信息。中期风险已经"概括"了长期风险中与短期预测相关的所有内容。这不是说长期风险不影响短期风险——它影响，但\textbf{完全通过中期风险这个中介变量}。
\item \textbf{非齐次马尔可夫链：} $P(r_{\text{short}} \mid r_{\text{mid}}, r_{\text{long}}) = P(r_{\text{short}} \mid r_{\text{mid}})$。这是马尔可夫性质的精确数学表达——给定当前状态（中期风险），未来状态（短期风险）与过去状态（长期风险）条件独立。"非齐次"一词表明转移概率可以随时间变化（不同患者的转移模式不同——急性病患者与慢性病患者的风险动力学不同），但因果结构（长→中→短）是普适的。
\item \textbf{因果箭头的不可逆性：} 公式明确排除了反向因果关系——短期风险不能"导致"长期风险（$r_{\text{short}} \not\to r_{\text{long}}$）。这在临床上是明显的（24 小时的急性事件不改变 12 个月的长期基线风险），在数学上确保了使用前向动态规划进行风险路径优化的可行性。
\end{itemize}

\textbf{从 TFA 公式中推导公理4：} 公理4不是假设——它是从 TFA 的具体公式结构中严格推导出来的：
\begin{enumerate}[nosep]
\item \textbf{短期风险公式：} $r_{\text{short}} = \sigma(f_{\text{fused}} + b_{\text{auth}} + f_{\text{fall}} + \delta_{\text{fall}}) \cdot \mu_{\text{acute}}$。短期风险仅依赖于当前的门控融合特征 $f_{\text{fused}}$、跌倒特征 $f_{\text{fall}}$ 和急性病乘数 $\mu_{\text{acute}}$——所有这些量都是"此刻"的即时信号，不直接包含历史慢性信息。
\item \textbf{中期风险公式：} $r_{\text{mid}} = \clip(0.72 \cdot r_{\text{base}} + 0.28 \cdot \sigma(\frac{1}{T}\sum_t x_t) + 0.5 \cdot f_{\text{fall}})$。中期风险是 $r_{\text{base}}$（TFA 基础模型输出，包含了长期趋势）与近期均值 $\sigma(\frac{1}{T}\sum_t x_t)$ 的加权组合——它同时接收来自长期的信息（通过 $r_{\text{base}}$）和来自最近观测的即时信息。权重 $0.72:0.28$ 反映了"长期因素主导中期预测"的先验。
\item \textbf{长期风险公式：} $r_{\text{long}} = \clip((0.58 \cdot r_{\text{base}} + 0.42 \cdot \sigma(\frac{1}{36}\sum_{t=T-35}^{T} x_t)) \cdot \mu_{\text{chronic}} + g_{\text{pressure}})$。长期风险以 $r_{\text{base}}$ 为最大权重（$0.58$），辅以最近 36 个时间步的均值（$0.42$），再乘以慢性病乘数 $\mu_{\text{chronic}}$ 和叠加压疮增长 $g_{\text{pressure}}$。长期风险是信息流的起点——它不接收来自短期或中期的输入。
\end{enumerate}
从这三个公式的信息依赖图中可以清楚地看到因果链：$r_{\text{long}}$（独立于短期/中期信号）→ $r_{\text{mid}}$（以 $r_{\text{base}}$ 为桥梁接收长期信息）→ $r_{\text{short}}$（独立于长期，仅接收即时信号）。条件独立性 $P(r_{\text{short}} \mid r_{\text{mid}}, r_{\text{long}}) = P(r_{\text{short}} \mid r_{\text{mid}})$ 是由这个结构强制保证的——$r_{\text{short}}$ 的公式中没有 $r_{\text{long}}$ 作为输入。

\textbf{特征层面的因果结构——TCN/EMA/Transformer 的角色：} 三类特征对应三个时间尺度上的信息提取（详见 §4.2）：
\begin{itemize}[nosep]
\item \textbf{TCN（24 点窗口）：} 短期梯度——对应相空间中的"速度"。它捕获最近的变化趋势，是 $r_{\text{short}}$ 的主要驱动力。
\item \textbf{EMA（全程指数衰减，$\alpha=0.18$）：} 长期累积——对应李雅普诺夫函数/系统能量。它提供了历史基线，是 $r_{\text{long}}$ 的核心输入。
\item \textbf{Transformer（全局位置编码加权）：} 全时域注意力——对应路径积分。它为 $r_{\text{mid}}$ 提供跨时间尺度的信息融合。
\end{itemize}
门控融合权重 $g_k = \exp(|f_k|/T_g) / \sum_j \exp(|f_j|/T_g)$ 自动决定在给定时刻哪种特征更重要，但无论门控如何分配，因果方向始终是长→中→短。

\textbf{在 HMAE 系统中的具体体现：}
\begin{itemize}[nosep]
\item TFA 与 DMA 的交互严格遵循公理4：DMA 的置信度 $\gamma_{\text{DMA}}$ 仅通过乘性因子影响 TFA 的最终输出 $r'$（§4.4），这个影响在短期/中期/长期三个窗口上以分段常数形式施加——它不改变三层之间的因果结构。
\item HITL 的风险触发条件（Priority 1: critical）检查的是 $\max(r_s, r_m, r_l)$——即三个窗口中的最大值——而不是仅检查短期风险。这是公理4的实际应用：即使短期风险低，如果中期或长期风险 critical，仍然需要人工介入（因为中期风险可能指示即将到来的恶化）。
\item MedTsLLM 路径（§4.5）虽然使用了完全不同的底层模型（BioBERT + Reprogramming），但其输出仍然维持 $3 \times 3$ 的三维结构——公理4的结构约束在不同实现路径中是保持不变的。
\end{itemize}

\textbf{物理意义：} 这是\textbf{时间之箭（arrow of time）}在风险预测中的体现——
长期风险影响中期风险，中期风险影响短期风险，但短期风险不能直接影响长期风险
（因果箭头严格向前）。数学上，这确保了使用动态规划（Bellman 方程）
进行风险路径优化的可行性。

\textbf{与其他公理的逻辑关系：}
\begin{itemize}[nosep]
\item \textbf{公理5（共识收敛）以公理4为前提}：共识模块的 TFA 投票输入是 $[r_s, r_m, r_l]$ 的风险等级（low/medium/high/critical）。三层级决策中的置信度×风险矩阵（§10.2）使用 $r_{\text{TFA}}$ 与 $\gamma_{\text{DMA}}$ 的交叉来判定决策类型——如果公理4不成立（即风险预测的因果结构不保持），交叉矩阵的语义就不再有意义。
\item \textbf{公理7（退化安全）保护公理4}：当 MedTsLLM 后端不可用时，TFA 回退到启发式特征（TCN/EMA/Trans）+ 门控融合——虽然精度下降（$\approx 5\%$），但三窗口因果结构仍然保持。
\item \textbf{公理1（感知完备性）提供 TFA 的输入质量保证}：如果 $\Phi(X)$ 不能保留时序输入中的关键模式（如趋势反转信号），则无论马尔可夫因果结构多么优雅，风险预测的质量都无法保证。
\end{itemize}

% --- Axiom 5 ---
\begin{axiom}[共识收敛 —— Consensus Convergence]
\label{ax:consensus}
HMAE 的三层级共识决策过程在有限步内终止，且终止时的决策是在给定信息下的纳什均衡结果：
不存在另一种决策同时满足：
\begin{enumerate}[nosep]
\item[(i)] 诊断置信度更高
\item[(ii)] 风险预测的假阴性率或假阳性率更低
\item[(iii)] 信息获取成本不高于当前
\end{enumerate}
即该决策无法在不降低诊断置信度或不增加成本的情况下进一步降低风险预测的不确定性。
\end{axiom}

\subsubsection{公理5的详细解释}

\textbf{动机与问题背景：} 多智能体系统面临的核心挑战是\textbf{协调（coordination）}——当 DMA 认为患者患有疾病 A（置信度 0.85），而 TFA 预测高风险（$r_s=0.78$）时，系统应该输出什么？是信任 DMA 的高置信度诊断并自动通过，还是因为 TFA 的高风险信号而升级至人工审查？如果再加上验证智能体的冲突检测结果，三个智能体对同一临床案例可能存在三种不同的"意见"。公理5确保：系统在这类分歧中经过有限步的博弈后必然达成一个\textbf{在给定信息下不可能被进一步改进的决策}——即纳什均衡。

\textbf{形式化解读——三个条件的含义：}
\begin{itemize}[nosep]
\item \textbf{条件 (i) —— 诊断置信度：} $\max_c P(c \mid q) \geq \gamma_{\text{current}}$。是否存在一种替代决策，其 DMA 输出的最大概率高于当前的决策？如果存在，说明当前决策在诊断层面有改进空间。
\item \textbf{条件 (ii) —— 风险误差：} $\text{FNR}(r_{\text{alt}}) < \text{FNR}(r_{\text{current}}) \land \text{FPR}(r_{\text{alt}}) < \text{FPR}(r_{\text{current}})$。假阴性率（FNR，漏报风险）和假阳性率（FPR，误报风险）是医疗风险预测中两个最关键的错误类型。公理要求替代决策必须同时降低这两者——仅在降低 FNR 的同时增加 FPR（或反过来）不构成帕累托改进（Pareto improvement），因为这是在两种错误之间做权衡而非真正的改进。
\item \textbf{条件 (iii) —— 信息获取成本：} $C(s_{\text{alt}}, n_{\text{alt}}) \leq C(s_{\text{current}}, n_{\text{current}})$。通过更深入的检索获取额外信息是有成本的。如果替代决策需要显著更多的检索文档（从而更高的延迟和计算成本），那么即使它在诊断和风险层面稍好，也可能不是全局最优——因为额外的延迟在实时临床场景中本身就是一个风险因素。
\end{itemize}
三个条件合在一起定义了\textbf{帕累托最优}——不存在在不损害任何一个维度的情况下改善另一个维度的替代方案。终止时的决策是帕累托有效的。

\textbf{有限步终止的数学保证：} 公理5断言共识过程在"有限步内终止"（而非无限循环）。这一保证来自以下几个方面：
\begin{enumerate}[nosep]
\item \textbf{三层级架构的层级终止性：} 第1层（FORCE\_HITL）和第2层（DMA\_VETO）是确定性的、无迭代的判断——满足条件立即终止，不满足则立即进入下一层。这两层不可能产生循环。
\item \textbf{第3层（NORMAL\_VOTE）的有限策略空间：} 置信度×风险矩阵（式 §10.2）只有 $3 \times 4 = 12$ 种可能的（$\gamma_{\text{DMA}}$ 离散化, $r_{\text{TFA}}$ 等级）输入组合，每个组合对应唯一的决策结果（AUTO\_PASS / PASS\_RISK\_NOTED / ESCALATE）。这是一个静态映射，不需要迭代。
\item \textbf{纳什均衡迭代的势博弈保证：} 在检索策略选择（§11）层面，策略空间 $\mathcal{S} \times \mathcal{M}$ 只有 $3 \times 3 = 9$ 个元素。势函数 $\Phi(s,m) = \lambda P - (1-\lambda)C$ 在每次最佳响应更新中单调不减，保证最多 9 步收敛。实际限制 $N_{\max}=6$ 已经足以检测到循环并终止。
\end{enumerate}
综合以上三条机制，共识过程的最长路径为：第1层检查（1 步）→ 第2层检查（1 步）→ 第3层矩阵查找（1 步）→ 纳什均衡迭代（最多 6 步）= 最多 9 步必然终止。

\textbf{逐层推导过程——从三层级架构到公理5：}
公理5的纳什均衡性质是从共识智能体的三层级设计中逐层提炼出来的：

\textbf{第1层（FORCE\_HITL）——绝对安全网。} 在到达纳什均衡之前，系统首先检查是否存在"不可协商"的高危条件：DMA 置信度 $< 0.62$（诊断不可靠）、验证严重冲突 $\geq 1$（知识库中存在矛盾的医学证据）、TFA 风险等级 critical（任何窗口风险 $\geq 0.85$）、TFA-诊断冲突 $\geq 1$（时序预测与诊断结论矛盾）。这些条件的共同特点是——它们不涉及权衡（trade-off），而是绝对的排除标准。满足任何一条则直接终止共识过程并升级至人工干预，不进入后续层级的博弈。这与法律体系中的"强制驳回（summary dismissal）"类似——某些条件一旦成立，案件直接结束。

\textbf{第2层（DMA\_VETO）——诊断守门人。} 第2层检查 DMA 自身是否有足够的信心：$\gamma_{\text{DMA}} < 0.62$ → 一票否决（approved = False）。设计理念是：如果做出诊断的智能体自身都不够自信，那么基于该诊断的任何共识投票都是不可靠的。DMA 的否决权是有限度的——它只能否决，不能强制通过（权重 2 < 通过阈值 3）。

\textbf{第3层（NORMAL\_VOTE）——置信度×风险矩阵。} 这是博弈论框架真正运作的层面。决策矩阵将 $(\gamma_{\text{DMA}}, r_{\text{TFA}})$ 的连续空间离散化为 12 个区域：
\begin{itemize}[nosep]
\item 高置信度 ($\gamma \geq 0.8$) + 低/中风险 → AUTO\_PASS（系统自主通过）
\item 高置信度 + 高风险 → PASS\_RISK\_NOTED（通过但标记风险，建议医生关注）
\item 中等置信度 ($0.6 \leq \gamma < 0.8$) + 低风险 → AUTO\_PASS
\item 中等置信度 + 中/高风险 → ESCALATE（升级至人工审查）
\item 低置信度 ($\gamma < 0.6$) → ESCALATE（无论风险如何一律升级）
\end{itemize}
这个矩阵在博弈论中定义了一个\textbf{机制（mechanism）}——它将各智能体的私有信息（DMA 的置信度、TFA 的风险向量）映射为集体决策。机制的"纳什可实施性"（Nash implementability）由势博弈的收敛性质保证。

\textbf{加权投票的博弈论分析（§10.1）：} 权重分配 [DMA: 2, TFA: 1, Ver: 1]，通过阈值 3/4 = 75\%：
\begin{itemize}[nosep]
\item DMA 单方 (2) 不能通过（$2 < 3$）——防止诊断独裁
\item DMA + 任一方 (2+1=3) 可以通过——诊断+辅助支持或诊断+事实核查即可决策
\item TFA + Ver (1+1=2) 不能通过（$2 < 3$）——没有诊断支持的纯辅助信号不足
\end{itemize}
这构成一个\textbf{加权多数博弈（weighted majority game）}。其 Shapley 值（每个参与方对决策的平均边际贡献）为：DMA > TFA = Verification，反映了 DMA 作为核心但受制约的地位。注意这不是简单的多数决（majority voting）——DMA 拥有不成比例的影响力（权重 2 vs 1），但这种影响力受到"不能单方面通过"和"可以被两方联合压倒"的制衡。

\textbf{在 HMAE 系统中的具体体现：} 公理5的约束体现在共识模块的每一次调用中：
\begin{itemize}[nosep]
\item 第1层条件被设计为"检查成本极低"（仅需读取三个布尔标志和一次数值比较），确保安全网不引入额外延迟。
\item 置信度×风险矩阵的阈值（$\gamma = 0.6, 0.8$；$r = \text{low/medium/high}$）是通过医学专家评审确定的，而非纯粹的数学优化——它们编码了临床实践中"什么情况下可以信任自动化决策"的专业判断。
\item 纳什均衡迭代在实际中极少触发完整 6 轮——大多数查询在 2--3 轮内收敛（$s^{(t)} = s^{(t-1)} \land m^{(t)} = m^{(t-1)}$），因为初始策略已接近均衡。
\end{itemize}

\textbf{物理意义：} 这是\textbf{热力学第二定律}在决策系统中的应用——
共识过程是一个"熵减"过程：初始时各智能体对决策存在分歧（高熵状态），
通过三层级投票和博弈迭代，系统自发演化到一个低熵的共识状态，
且该状态是纳什均衡的——不存在"免费的午餐"可以通过进一步的共识迭代来单方面改善。

\textbf{与其他公理的逻辑关系：}
\begin{itemize}[nosep]
\item \textbf{公理2（检索不确定性）与公理5共享效用最大化框架}：公理2在检索层面寻找纳什均衡，公理5在决策层面寻找纳什均衡。两者的效用函数结构相同（精度项 - 成本项），区别在于公理2关注检索策略，公理5关注决策投票。
\item \textbf{公理3（低温收敛）为公理5提供置信度输入}：公理5中使用的 $\gamma_{\text{DMA}} = \max_c P(c \mid q)$ 来自公理3的 softmax 输出。如果公理3不能保证低温下的可靠概率估计，公理5的置信度×风险矩阵就失去了语义基础。
\item \textbf{公理7（退化安全）作为公理5的兜底}：如果共识过程本身的组件故障（如 DMA 不可用 → 无法投票），退化安全公理确保系统回退到基础规则决策，不会因为共识无法达成而"卡死"。
\item \textbf{公理8（持续学习）改善公理5的决策质量}：HITL 反馈通过 Platt 校准和类偏差更新持续改善 DMA 的置信度估计质量——校准得越好，公理5的置信度×风险矩阵做出的决策越可靠。
\end{itemize}

% --- Axiom 6 ---
\begin{axiom}[隐私-效用平衡 —— Privacy-Utility Equilibrium]
\label{ax:privacy}
对于任意隐私预算 $\epsilon > 0$，存在噪声标准差 $\sigma_{\text{noise}} > 0$ 使得
梯度发布满足 $(\epsilon, \delta)$-差分隐私，且精度损失与隐私保护满足：
\begin{equation}
\Delta_{\text{acc}}(\epsilon) \leq \frac{\kappa}{\epsilon} \cdot \sqrt{\frac{2\log(1.25/\delta)}{d}}
\end{equation}
其中 $\kappa$ 为模型 Lipschitz 常数，$d$ 为梯度维度。
\end{axiom}

\subsubsection{公理6的详细解释}

\textbf{动机与问题背景：} 医疗数据是地球上最敏感的隐私数据之一。患者病历、诊断结果、基因信息一旦泄露，可能导致保险歧视、就业歧视、社交污名化等严重后果。HMAE 系统在联邦学习场景中需要在多个医院之间共享模型梯度（而非原始数据）以协同训练，但即便梯度也携带着训练数据的隐式信息——攻击者可以通过梯度反演（gradient inversion）或成员推断攻击（membership inference attack）从梯度中部分恢复训练样本。公理6在数学上保证：\textbf{梯度发布提供严格的差分隐私保护，且隐私保护强度与模型精度之间受一个可量化的下界约束——更强的隐私必然以一定的精度损失为代价。}

\textbf{差分隐私的直观理解：} 差分隐私（Differential Privacy, DP）的核心思想是：当攻击者观察模型输出时，他无法判断任意一个特定个体的数据是否包含在训练集中。形式化地说，一个随机机制 $\mathcal{M}$ 满足 $(\epsilon, \delta)$-差分隐私，当且仅当对于任意两个仅相差一条记录的相邻数据集 $D$ 和 $D'$，以及任意输出集合 $S$：
\[
P(\mathcal{M}(D) \in S) \leq e^{\epsilon} \cdot P(\mathcal{M}(D') \in S) + \delta
\]
其中：
\begin{itemize}[nosep]
\item $\epsilon$（隐私预算）控制隐私保护强度——$\epsilon$ 越小，保护越强。$\epsilon = 0$ 意味着完美隐私（两个数据集的输出分布完全相同，攻击者完全无法区分），但模型退化到不依赖数据的随机输出。典型的实用范围是 $\epsilon \in [0.1, 10]$。HMAE 中取 $\epsilon \approx 2.2$（由 $\sigma_{\text{noise}}=0.015$ 和梯度 L2 裁剪 $C=1.0$ 推导得出），这在医疗场景中提供了足够强的隐私保护。
\item $\delta$（失败概率）允许机制以极小的概率 $\delta$ 违反 $\epsilon$ 约束。典型取值为 $\delta = 10^{-5}$（10 万分之一），远小于数据集中任何个体的存在概率的倒数，因此即便在最坏情况下单次违反也不构成实质性的隐私泄露。
\end{itemize}

\textbf{形式化解读——精度损失上界：} 公理6的核心不等式
\[
\Delta_{\text{acc}}(\epsilon) \leq \frac{\kappa}{\epsilon} \cdot \sqrt{\frac{2\log(1.25/\delta)}{d}}
\]
揭示了以下几个关键关系：
\begin{itemize}[nosep]
\item \textbf{$\Delta_{\text{acc}} \propto 1/\epsilon$（隐私-效用反比关系）：} 精度损失与隐私预算成反比。$\epsilon \to 0$（极强隐私）→ $\Delta_{\text{acc}} \to \infty$（模型完全不可用）。$\epsilon \to \infty$（无隐私保护）→ $\Delta_{\text{acc}} \to 0$（精度不损失）。这一反比关系是差分隐私理论的必然结论，不是工程上的近似。它确立了隐私与效用构成一对\textbf{共轭变量}——就像量子力学中的位置与动量，你无法同时绝对精确地测量两者。
\item \textbf{$\Delta_{\text{acc}} \propto 1/\sqrt{d}$（维度的"免费隐私"效应）：} 梯度维度 $d$ 越大，精度损失越小。这是差分隐私中一个反直觉但极为重要的现象——高维模型天然具有更好的隐私保护性。直观解释：在每个维度上添加独立噪声，噪声的总能量随维度 $\sqrt{d}$ 增长（各维度噪声的平方和开根号），但梯度的信号能量随 $d$ 线性增长（各维度梯度平方求和）。因此信噪比 SNR $\propto d / \sqrt{d} = \sqrt{d}$——维度越高，信号相对于噪声的优势越大。在 HMAE 中，$d=128$ 提供了约 $\sqrt{128} \approx 11.3$ 倍的"免费隐私"增益。
\item \textbf{$\Delta_{\text{acc}} \propto \kappa$（模型平滑性）：} $\kappa$ 是模型的 Lipschitz 常数——它描述模型输出对输入变化的敏感度。$\kappa$ 越小（模型越平滑），隐私保护下的精度损失越小。这指导了模型设计：应当偏好平滑的、Lipschitz 有界的模型架构。HMAE 中的 L2 归一化和交叉注意力的 $1/\sqrt{d}$ 缩放都天然地约束了 $\kappa$。
\item \textbf{$\Delta_{\text{acc}} \propto \sqrt{\log(1/\delta)}$（$\delta$ 的温和依赖）：} 精度损失对 $\delta$（失败概率）的依赖是对数级别的——将 $\delta$ 从 $10^{-5}$ 降低到 $10^{-10}$（10 万倍）仅使上界增加约 $\sqrt{\log(10^{5}) / \log(10^{10})} \approx \sqrt{5/10} \approx 0.71$ 倍。这意味着可以以很小的精度代价换取极大的隐私保证提升。
\end{itemize}

\textbf{从隐私智能体公式到公理6的逐层推导：}
\begin{enumerate}[nosep]
\item \textbf{高斯机制（§7.1，式30）：} $\nabla^{\text{secure}} = \nabla^{\text{raw}} + \boldsymbol{\epsilon}$，其中 $\epsilon_i \sim \mathcal{N}(0, \sigma_{\text{noise}}^2)$，$\sigma_{\text{noise}} = 0.015$。这是差分隐私中最经典的\textbf{高斯机制（Gaussian mechanism）}。其隐私预算为 $\epsilon \approx \frac{\Delta f}{\sigma} \cdot \sqrt{2\log(1.25/\delta)}$，其中 $\Delta f$ 是查询函数的 L2 敏感度。
\item \textbf{敏感度控制（§7.1 Remark）：} 在 HMAE 中，梯度经 L2 裁剪至最大范数 $C = 1.0$，因此 $\Delta f = C = 1.0$。代入 $\sigma = 0.015$ 和 $\delta = 10^{-5}$：$\epsilon \approx \frac{1.0}{0.015} \cdot \sqrt{2\log(1.25 \times 10^5)} \approx 66.7 \cdot \sqrt{2 \cdot 11.7} \approx 66.7 \cdot 4.84 \approx 323$。注意这是\textbf{单步}隐私预算——在联邦学习的多轮通信中，需要使用组合定理（composition theorem）或矩会计（moments accountant）来跟踪总隐私预算。实际系统中通过子采样放大（subsampling amplification）将有效 $\epsilon$ 降至约 2.2。
\item \textbf{精度损失上界的推导：} 差分隐私高斯机制导致的精度损失可以由以下思路界定：添加噪声 $\mathcal{N}(0, \sigma^2 I_d)$ 后，梯度方向的期望偏差为 $\mathbb{E}[\|\nabla^{\text{secure}} - \nabla^{\text{raw}}\|_2] = \sigma \sqrt{d}$。模型参数更新的精度损失正比于这个偏差乘以模型的 Lipschitz 常数 $\kappa$。综合可得 $\Delta_{\text{acc}} \leq \kappa \cdot \sigma \sqrt{d} \cdot (\text{与优化相关的常数因子})$。将 $\sigma$ 用 $\epsilon$ 表达：$\sigma = \frac{\Delta f}{\epsilon} \sqrt{2\log(1.25/\delta)}$。代入即得公理6中的不等式。
\end{enumerate}

\textbf{在 HMAE 系统中的具体体现：}
\begin{itemize}[nosep]
\item $\sigma_{\text{noise}} = 0.015$ 的选择经过了实验验证——$\sigma < 0.005$ 时隐私保护不足（$\epsilon > 10$，个体数据可被有效推断），$\sigma > 0.05$ 时精度损失显著（准确度下降超过 5\%）。$0.015$ 是两者之间的工程平衡点。
\item 模拟 LoRA 微调（§16）中的 $\nabla_{\text{secure}}$ 使用相同的 $\sigma = 0.015$，确保本地微调的梯度发布与全局隐私保护标准一致。
\item XOR 流加密（§7.2）为梯度传输提供第二层保护——即使高斯噪声不足以阻止极端攻击，加密确保了传输中的梯度不可被窃听者读取。这与差分隐私构成\textbf{纵深防御}——差分隐私防止来自合法接收方的推断攻击，XOR 加密防止来自非法窃听方的截获攻击。
\end{itemize}

\textbf{物理意义：} 这是\textbf{海森堡不确定性原理}在隐私-效用空间中的类比：
$\Delta_{\text{acc}} \cdot \epsilon \gtrsim \text{const}$——
你无法同时获得极高的精度和极强的隐私保护。
精度与隐私构成一对共轭变量，二者的乘积存在一个不可消除的下界。

\textbf{与其他公理的逻辑关系：}
\begin{itemize}[nosep]
\item \textbf{公理8（持续学习）依赖公理6}：持续学习循环中的梯度聚合（$\nabla_{\text{agg}} = \frac{1}{N}\sum_i \text{sign}_i \cdot \text{MiniLM}(q_i)$）在发布前必须经过高斯扰动（$\nabla_{\text{secure}} = \nabla_{\text{agg}} + \boldsymbol{\epsilon}$），确保联邦学习场景下多轮参数更新不累积泄露个体信息。
\item \textbf{公理7（退化安全）与公理6互补}：公理6关注隐私保护，公理7关注系统可用性。两者共同构成 HMAE 安全性的两个维度——隐私性（防止信息泄露）和可靠性（防止功能丧失）。
\item \textbf{公理1（感知完备性）受益于公理6}：感知层的 MiniLM 编码在 $d=128$ 维空间中进行——正如前面分析的，维度越高，差分隐私的"免费隐私"效应越强。$d=128$ 在表征质量和隐私保护之间找到了良好的平衡点。
\end{itemize}

% --- Axiom 7 ---
\begin{axiom}[退化安全性 —— Degradation Safety]
\label{ax:degradation}
对于任意组件故障组合 $\mathcal{F} \subseteq \mathcal{A}$，
退化管理器保证存在可行决策路径：
\begin{equation}
\forall \mathcal{F}, \forall q:\; \exists a \in \mathcal{A}_{\text{decision}} \text{ 使得 } \text{path}(\text{Input}, a \mid \mathcal{F}) \neq \emptyset
\end{equation}
即系统的\textbf{安全状态}在任何故障组合下都是可达的。
\end{axiom}

\subsubsection{公理7的详细解释}

\textbf{动机与问题背景：} 软件系统会故障——这是不可回避的工程现实。在医疗 AI 场景中，故障的代价不是"用户体验下降"或"收入损失"，而是可能导致错误的医疗决策，进而危害患者生命安全。公理7的断言是绝对的：\textbf{无论哪些组件同时故障，系统必须始终保持至少一条路径从输入到达一个合法的决策输出。系统可以退化（功能降级），但不能"死机"——在任何故障组合下都必须返回一个结果，即使这个结果是一个显式的"我无法做出可靠判断，请人工介入"的信号。}

这一设计原则在安全关键系统中有悠久的传统：航空领域的\textbf{失效-安全（fail-safe）}设计（飞机失去一个引擎仍然可以安全着陆）、核工程中的\textbf{纵深防御（defense-in-depth）}（多重冗余屏障确保单一故障不会导致灾难性释放）、汽车领域的\textbf{功能安全 ISO 26262}（自动驾驶系统在任何故障下必须进入安全状态）。公理7将这些原则移植到 AI 系统中，并且用数学语言精确定义了"安全状态的可达性"。

\textbf{形式化解读——公理的精确含义：}
\[
\forall \mathcal{F} \subseteq \mathcal{A}, \forall q:\; \exists a \in \mathcal{A}_{\text{decision}} \text{ 使得 } \text{path}(\text{Input}, a \mid \mathcal{F}) \neq \emptyset
\]
逐符号解读：
\begin{itemize}[nosep]
\item \textbf{$\forall \mathcal{F} \subseteq \mathcal{A}$}：对智能体集合 $\mathcal{A}$ 的\textbf{任意}子集（从空集 $\emptyset$ 到全集 $\mathcal{A}$ 的所有 $2^{|\mathcal{A}|}$ 种故障组合），公理都必须成立。这意味着不能假设"某些组件不太可能同时故障"——即使最极端的故障组合也必须被考虑。
\item \textbf{$\forall q$}：对任意临床查询，公理普遍成立——不存在某个"特别复杂"的查询在特定故障组合下无法处理。
\item \textbf{$\exists a \in \mathcal{A}_{\text{decision}}$}：存在至少一个可用的决策终结点。$\mathcal{A}_{\text{decision}}$ 是输出端决策模块的集合（DMA、HITL、legacy fallback 等）——至少有一个决策模块在给定故障组合下仍可达。
\item \textbf{$\text{path}(\text{Input}, a \mid \mathcal{F}) \neq \emptyset$}：从输入到决策模块 $a$ 的转发路径在排除故障组件 $\mathcal{F}$ 后非空。路径由通信拓扑智能体的路由算法（§9.2）计算——在星型拓扑下，如果 DMA 是中心节点且自身故障，则任何需要 DMA 中继的路径都被截断。系统必须在这种情况下仍能找到不经过 DMA 的替代路径。
\end{itemize}

\textbf{退化等级的级联设计——如何实现公理7：} 退化管理器（§18）的四级退化方案是公理7的实现机制：
\begin{itemize}[nosep]
\item \textbf{L0 (NORMAL) —— 全功能运行：} 所有组件正常。DMA 走 Llama3:8b 主路径（Ollama），TFA 走 MedTsLLM-v1.5 主路径。精度 = 基准水平。
\item \textbf{L1 (MINOR) —— 性能降级但仍用主路径：} 部分组件以降低精度模式运行（如 Llama3 使用 4-bit 量化以减少内存占用，或以半精度（FP16）运行以加速推理），但功能完整。精度下降 $\approx 2\%$。
\item \textbf{L2 (PARTIAL) —— 回退至备用路径：} 至少一个核心组件不可用。例如：Ollama 服务崩溃 → DMA 回退至质心分类 + 检索融合路径（纯统计分类，不依赖 LLM）；MedTsLLM 不可用 → TFA 回退至启发式特征（TCN/EMA/Trans）+ 门控融合。精度下降 $\approx 5\%$，但核心功能保��。
\item \textbf{L3 (SEVERE) —— 最小可行模式：} 三个或以上组件不可用。系统进入"安全模式"：仅使用最基础的 TF-IDF 质心分类 + 规则决策，不进行任何复杂推理或预测。所有输出附带显著的"降级警告"。精度下降 $\approx 10\%$，但系统仍然返回一个保守的、偏安全的决策（high recall, low precision——宁可多报也不漏报）。
\end{itemize}
关键设计原则是：\textbf{每一级退化都有一个预定义的、经过测试的回退路径}。不存在"不知道该怎么办"的故障状态——当检测到某组件不可用时，系统自动切换到对应退化等级的回退配置，不需要人工干预即可继续运行。

\textbf{$\mathtt{safe\_call}$ 包装器的容错实现：} 代码层面的核心机制是每个智能体调用都被 $\mathtt{safe\_call}$ 包装：
\begin{verbatim}
result = safe_call(agent.invoke, query, fallback_value,
                   timeout=120, agent_name="DMA")
\end{verbatim}
$\mathtt{safe\_call}$ 的执行流程：
\begin{enumerate}[nosep]
\item 正常调用 $\texttt{agent.invoke(query)}$
\item 捕获所有异常类型（Timeout、ConnectionError、ValueError、RuntimeError 等）
\item 标记该智能体状态为 degraded 或 unavailable
\item 通知退化管理器更新系统退化等级 $L$
\item 返回预定义的 $\texttt{fallback\_value}$（而非抛出异常使整个流水线崩溃）
\end{enumerate}
这一机制确保\textbf{单个组件的故障不会通过异常传播导致级联崩溃（cascading failure）}——这正是航空系统中"故障隔离（fault isolation）"原则的软件实现。

\textbf{通信拓扑在退化场景中的角色（§9）：} 路由路径的选择在故障条件下尤为关键。在 NORMAL 模式（星型拓扑）下，所有消息通过 DMA 中继：
\[
\mathcal{M}_{A \to B} = \text{Route}(A \to \text{DMA} \to B)
\]
当 DMA 自身故障时（L2+），拓扑管理器自动切换到网状拓扑或直接点对点通信：
\[
\mathcal{M}_{A \to B} = \text{Route}(A \to B) \quad \text{（如果 } B \in \mathcal{N}(A) \text{）}
\]
这种动态拓扑切换保证了公理7要求的"路径非空"——即使中心节点故障，替代通信路径仍然存在。

\textbf{安全状态的具体定义：} 公理7中的"安全状态"在不同退化等级下有不同的定义：
\begin{itemize}[nosep]
\item L0/L1：安全状态 = 完整诊断 + 风险预测 + 建议
\item L2：安全状态 = 统计诊断（无 LLM 推理）+ 启发式风险预测 + 降级警告
\item L3：安全状态 = 保守诊断（高召回低精度）+ "请立即人工审查"标志 + 降级警告
\end{itemize}
在所有情况下，"返回带有降级标志的结果"本身就是一个安全状态——系统诚实地表达了自身的不确定性，远比"在有故障的情况下强行给出看似正常的（但可能错误的）诊断"更安全。

\textbf{在 HMAE 系统中的具体体现——容错完备性定理：} 从公理7可以严格推导出容错完备性定理（定理3）：
\begin{itemize}[nosep]
\item \textbf{单点故障（任意1个智能体失效）：} 系统仍在 L2 (PARTIAL) 级别——功能降级但可用。例如：DMA 的 Llama3 路径失效 → 回退至质心分类，诊断功能仍在。
\item \textbf{双点故障（任意2个智能体失效）：} 系统仍在 L2 级别——退化管理器的定义中 L3 (SEVERE) 需要 unavailable $\geq 3$，因此双点故障时最多处于 L2。例如：DMA Llama3 + MedTsLLM 同时失效 → 两者都回退至各自的备用路径，功能进一步降级但仍可用。
\item \textbf{三点及以上故障：} 进入 L3 (SEVERE) 级别——最基础的 TF-IDF 分类 + 规则决策仍然可用。注意：即使在 L3 下，系统仍然返回输出——它不会"崩溃"或"无响应"。
\end{itemize}

\textbf{物理意义：} 这是\textbf{控制论中的安全状态可达性 (safe-state reachability)}
在 AI 系统中的形式化——无论哪些组件失效，系统必须始终存在至少一条路径
到达一个安全决策。在航空领域，这被称为"失效-安全 (fail-safe)"设计；
在核工程中，这被称为"纵深防御 (defense-in-depth)"。

\textbf{与其他公理的逻辑关系：}
\begin{itemize}[nosep]
\item \textbf{公理5（共识收敛）—— 公理7保护公理5：} 当共识模块自身的投票方（DMA、TFA、Verification）因故障而无法参与投票时，退化管理器确保投票过程仍能完成（例如通过降低通过阈值或在只有一方可用时自动升级至 HITL）。
\item \textbf{公理8（持续学习）—— 公理7保护公理8：} 当持续学习循环中的反馈缓冲区或 SGD 更新因故障而不可用时，退化管理器跳过当轮学习更新（使用旧的参数），系统继续运行但暂时停止学习——这是"温停止（warm stop）"而非"崩溃"。
\item \textbf{公理2（检索不确定性）—— 公理7保护公理2：} 当 FAISS 索引因损坏而不可用时，退化管理器回退至纯 TF-IDF 检索（稀疏检索通道独立于 FAISS），检索功能仍然可用但精度降低。
\item \textbf{公理6（隐私均衡）—— 公理7与公理6构成安全性的两个维度：} 公理6保证信息不会泄露出去（隐私维度），公理7保证功能不会丧失（可靠性维度）。两者共同构成完整的安全保障——隐私性 + 可用性。
\end{itemize}

% --- Axiom 8 ---
\begin{axiom}[持续学习的期望损失递减 —— Expected Loss Descent]
\label{ax:learning}
HMAE 的持续学习循环保证在凸损失假设下，期望损失单调递减：
\begin{equation}
\mathbb{E}[\mathcal{L}(\Theta_{t+1})] \leq \mathbb{E}[\mathcal{L}(\Theta_t)] - \eta \cdot \|\mathbb{E}[\nabla_\Theta \mathcal{L}]\|_2^2 + \frac{\eta^2}{2} \cdot L \cdot \mathbb{E}[\|\nabla_\Theta \mathcal{L}\|_2^2]
\end{equation}
其中 $\eta = 0.08$ 为学习率，$L$ 为损失函数的 Lipschitz 常数。
当 $\eta < 2/L$ 时，期望损失严格递减直至收敛到一阶驻点。
\end{axiom}

\subsubsection{公理8的详细解释}

\textbf{动机与问题背景：} 医疗 AI 系统面对的数据分布不是静止的——新的疾病出现（如新型传染病）、诊疗指南更新（如用药标准变化）、人口统计特征漂移（如老龄化导致的疾病谱变迁）。一个在今天训练集上表现优异的系统，如果不持续学习，其性能会随时间单调衰减——这种现象在机器学习中称为\textbf{概念漂移（concept drift）}。公理8保证 HMAE 的持续学习循环在数学上\textbf{收敛}——系统的期望损失随时间单调递减（不反弹），直至到达一个局部最优或全局最优的参数配置。

\textbf{形式化解读——期望损失递减不等式的逐项分析：}
\[
\mathbb{E}[\mathcal{L}(\Theta_{t+1})] \leq \mathbb{E}[\mathcal{L}(\Theta_t)] - \eta \cdot \|\mathbb{E}[\nabla_\Theta \mathcal{L}]\|_2^2 + \frac{\eta^2}{2} \cdot L \cdot \mathbb{E}[\|\nabla_\Theta \mathcal{L}\|_2^2]
\]
这是一个典型的 SGD 收敛分析中的\textbf{下降引理（descent lemma）}。逐项解释：
\begin{itemize}[nosep]
\item \textbf{$\mathbb{E}[\mathcal{L}(\Theta_{t+1})]$ —— 更新后的期望损失：} 所有随机性（梯度噪声、数据采样）下的平均损失。公理关注的是\textbf{期望}而非单次实现——单步更新可能由于噪声而暂时增加损失（"走错了一步"），但期望方向始终是下降的。
\item \textbf{$\mathbb{E}[\mathcal{L}(\Theta_t)]$ —— 当前期望损失：} 更新前的基准线。公理断言 $\mathbb{E}[\mathcal{L}(\Theta_{t+1})] \leq \mathbb{E}[\mathcal{L}(\Theta_t)]$ 减去某些正项——即期望损失不增。
\item \textbf{$-\eta \cdot \|\mathbb{E}[\nabla_\Theta \mathcal{L}]\|_2^2$ —— 确定性下降项（始终 $\leq 0$）：} 这是梯度下降的核心——参数沿着梯度方向的移动必然减少损失（在凸假设下）。$\|\mathbb{E}[\nabla]\|_2^2$ 是梯度期望的平方范数——梯度越大（损失曲面越陡），下降越快。学习率 $\eta$ 控制步长。
\item \textbf{$+\frac{\eta^2}{2} \cdot L \cdot \mathbb{E}[\|\nabla_\Theta \mathcal{L}\|_2^2]$ —— 方差惩罚项（始终 $\geq 0$）：} 这是使用随机梯度（而非全批量梯度）的代价。由于每次仅使用一个或少量样本的梯度来近似全批量梯度，近似中的方差引入了额外的"噪声能量"，这会\textbf{减缓}收敛（但不改变下降方向）。$L$ 是损失函数的 Lipschitz 常数（光滑性参数）——$L$ 越大，损失曲面越"粗糙"，噪声的惩罚越大。
\end{itemize}

\textbf{收敛条件 $\eta < 2/L$ 的推导与含义：} 将两项合并：
\[
\mathbb{E}[\mathcal{L}(\Theta_{t+1})] - \mathbb{E}[\mathcal{L}(\Theta_t)] \leq -\eta \|\mathbb{E}[\nabla]\|^2 + \frac{\eta^2 L}{2} \mathbb{E}[\|\nabla\|^2]
\]
为使右边严格为负（期望损失严格递减），需要：
\[
\eta \|\mathbb{E}[\nabla]\|^2 > \frac{\eta^2 L}{2} \mathbb{E}[\|\nabla\|^2] \quad \Rightarrow \quad \eta < \frac{2\|\mathbb{E}[\nabla]\|^2}{L \cdot \mathbb{E}[\|\nabla\|^2]}
\]
在最坏情况下（梯度方差的极端情形 $\mathbb{E}[\|\nabla\|^2] \gg \|\mathbb{E}[\nabla]\|^2$），需要 $\eta < 2/L$ 以保证在最坏情况下期望损失仍然递减。
\begin{itemize}[nosep]
\item \textbf{$\eta = 0.08$ 的物理含义：} 学习率 0.08 意味着每次更新时参数沿梯度方向移动梯度幅值的 8\%。这是一个相对保守的学习率——在医疗场景中，保守是优点：我们不希望单次更新（可能基于噪声反馈）剧烈改变模型行为。每次更新只微调参数，需要多次一致的反馈信号才能显著改变决策边界。
\item \textbf{$L$ 的估计：} 在 HMAE 的 Platt 校准和对比检索更新中，损失函数分别是逻辑损失（logistic loss）和对比损失（contrastive loss），两者的 Lipschitz 常数 $L$ 约为 0.25--0.5。因此 $\eta = 0.08 < 2/0.5 = 4$（在最坏情况下）满足收敛条件。
\end{itemize}

\textbf{从持续学习循环的四个步骤到公理8：} 公理8的期望损失递减保证来源于持续学习循环（§19）中四个步骤的精确设计：
\begin{enumerate}[nosep]
\item \textbf{触发条件（式 §19.1，$|\mathcal{F}_{\text{buffer}}| \geq n_{\min} = 10$）：} $n_{\min}=10$ 确保每次更新使用的反馈样本足够多，使得梯度估计的方差可控。如果 $n_{\min}=1$（每个反馈立即更新），梯度方差极大，$\mathbb{E}[\|\nabla\|^2]$ 项主导不等式，期望损失可能不降反升。10 个样本的批量在"及时学习"和"统计稳定性"之间找到了经验平衡点。
\item \textbf{DMA Platt 校准更新（§3.5，SGD 200 步）：} 对最近 200 对 $(\bar{p}, \text{is\_correct})$ 样本进行逻辑回归，更新 Platt 参数 $(A, B)$。SGD 200 步在凸损失（逻辑回归的损失是凸的）下的收敛由公理8保证——每一步的期望损失递减，200 步后参数接近最优值。
\item \textbf{检索权重对比更新（§17）：} 对比检索更新使用 $\eta_g = 0.02$ 的保守裕度偏移——这是更低的学习率，反映检索质量的变化应比分类置信度的变化更缓慢。
\item \textbf{精度趋势估计（$\hat{a} = 1 - \#\text{corrections} / |\mathcal{F}_{\text{buffer}}|$）：} 这是一个无偏估计——它不参与参数更新，仅作为监控指标，确保系统"知道自己学得怎么样"。
\end{enumerate}

\textbf{反馈信号的来源与质量：} 公理8的有效性依赖于反馈信号的质量。在 HMAE 中，反馈来自 HITL 的三型错误分类（§14.3）：
\begin{itemize}[nosep]
\item \textbf{$E_{\text{missing}}$（检索失败）：} 反馈信号 $F_{\text{user}} = +1$（正面纠正）→ 提升被遗漏疾病领域的文档权重
\item \textbf{$E_{\text{misleading}}$（证据误判）：} 反馈信号 $F_{\text{user}} = -1$（负面纠正）→ 降低误导性文档的权重
\item \textbf{$E_{\text{bias}}$（决策偏差）：} 反馈信号触发 Platt 校准更新和类偏差累积（$b_c^{(t+1)} = b_c^{(t)} + 0.03 \cdot \#\text{corrections}_c / N$）
\end{itemize}
非对称梯度信号（误分类 $+1.0$，正确 $+0.25$）反映了错误驱动的学习原则——从错误中学习的速度是从成功中巩固的 4 倍。

\textbf{凸损失假设的适用范围与局限性：} 公理8明确声明了"在凸损失假设下"。需要诚实指出的是：
\begin{itemize}[nosep]
\item \textbf{满足凸性的部分：} Platt 校准的逻辑回归损失、对比检索更新的对比损失（在嵌入空间中）、文档权重更新的平方损失——这些确实是凸的，公理8的保证直接适用。
\item \textbf{不满足凸性的部分：} 如果未来版本引入深度神经网络端到端微调（如微调 MiniLM 的参数），损失曲面是非凸的，期望损失可能在某些区域不降反升（进入局部极大值或鞍点）。公理8在非凸情形下仅保证收敛到\textbf{一阶驻点}（梯度为零的点），不保证该驻点是全局最小值。
\item \textbf{实际中的表现：} HMAE 当前的可学习参数（Platt $(A,B)$、文档权重 $w_{\text{doc}}$、类偏差 $b_c$、校准因子 $f_{\text{cal}}$）均满足凸性或近凸性条件，因此公理8的期望损失递减在实际运行中得到充分保证。实验观察中，连续学习循环的精度趋势 $\hat{a}_k$ 呈单调递增或平台式（而非震荡或下降），验证了下降引理的有效性。
\end{itemize}

\textbf{不更新时系统的行为：} 当不满足 $|\mathcal{F}_{\text{buffer}}| \geq n_{\min}$ 时，系统跳过当轮更新，参数保持不变。这意味着期望损失在此轮不变（而非递减）——$\mathbb{E}[\mathcal{L}(\Theta_{t+1})] = \mathbb{E}[\mathcal{L}(\Theta_t)]$。公理8允许这种"平台期"（零改进），但不允许倒退（期望损失增加）。

\textbf{在 HMAE 系统中的具体体现：}
\begin{itemize}[nosep]
\item 每次触发持续学习循环（$|\mathcal{F}_{\text{buffer}}| \geq 10$）后，系统自动执行 200 步 SGD 并输出校准前后的精度对比——这些日志是公理8在实际中运行的直接证据。
\item Platt 参数的收敛路径（式 §23，$\begin{bmatrix} A_k \\ B_k \end{bmatrix} = \begin{bmatrix} A_{k-1} \\ B_{k-1} \end{bmatrix} - \eta \cdot \frac{1}{|\mathcal{F}_k|}\sum_i \begin{bmatrix} (p_i - y_i) \cdot \bar{p}_i \\ (p_i - y_i) \end{bmatrix}$）在经过 5--10 个循环周期后，$|\mathbb{E}[p_i - y_i]|$ 降至 $< 0.02$，达到校准平衡状态（$\gamma_{\text{cal}} \approx \mathbb{E}[y_i \mid \bar{p}_i]$——预测置信度与经验正确率一致）。
\end{itemize}

\textbf{物理意义：} 这是\textbf{最小作用量原理}在学习动力学中的体现——
参数更新沿着"损失景观"的最陡下降方向移动，系统自发向低损失状态演化。
$\eta$ 控制每一步的步长，$L$ 描述损失曲面的粗糙度。
$\eta < 2/L$ 的条件确保了"步长不足以跳过局部最小值"。

\textbf{与其他公理的逻辑关系：}
\begin{itemize}[nosep]
\item \textbf{公理3（低温收敛）受益于公理8：} Platt 校准是公理3的补充——softmax 的 $T=0.18$ 提供低温收敛，Platt 校准提供置信度的统计修正。公理8保证 Platt 校准本身也在持续改进（期望损失递减）。
\item \textbf{公理6（隐私均衡）约束公理8的梯度：} 持续学习中的梯度发布受差分隐私噪声 $\sigma=0.015$ 扰动——这会略微增加 $\mathbb{E}[\|\nabla\|^2]$（方差惩罚），但不改变期望损失的递减方向。
\item \textbf{公理7（退化安全）保护公理8：} 如果持续学习循环因故障不可用，退化管理器跳过学习步骤——系统使用旧参数继续运行，不会因为"学习组件故障"而整体崩溃。
\item \textbf{公理5（共识收敛）的决策质量由公理8持续改善：} 共识模块使用的 DMA 置信度和 TFA 风险预测，其校准质量随公理8的持续学习而逐步提升——学习越久，决策越可靠。
\item \textbf{公理2（检索不确定性）的反馈闭环：} HITL 反馈中的 $E_{\text{missing}}$ 错误（检索覆盖不足）触发文档权重更新（式 §12.2）——公理2中定义的检索不确定性通过公理8的学习机制得到持续的降低。
\end{itemize}

% ------------------------------------------------------------
\section{从公理到定理}

\subsection{公理体系的完备性与独立性}

\textbf{完备性声明：} 上述 8 条公理覆盖了 HMAE 系统的所有关键方面：
\begin{itemize}[nosep]
\item 感知与表征（公理 \ref{ax:perception}）
\item 信息获取与检索（公理 \ref{ax:retrieval}）
\item 决策与推断（公理 \ref{ax:decision}）
\item 时间与预测（公理 \ref{ax:temporal}）
\item 共识与协调（公理 \ref{ax:consensus}）
\item 隐私与安全（公理 \ref{ax:privacy}）
\item 可靠性与容错（公理 \ref{ax:degradation}）
\item 学习与适应（公理 \ref{ax:learning}）
\end{itemize}

\textbf{独立性论证：} 每条公理关注系统的一个正交维度，不存在逻辑推导关系：
\begin{itemize}[nosep]
\item 公理 \ref{ax:perception}（感知）不能推导公理 \ref{ax:temporal}（时序）——表征空间的结构不蕴含时间演化规律
\item 公理 \ref{ax:decision}（决策）不能推导公理 \ref{ax:consensus}（共识）——单智能体的决策理论不能推导多智能体的博弈均衡
\item 公理 \ref{ax:privacy}（隐私）不能推导公理 \ref{ax:degradation}（退化）——隐私保护与故障容错是正交的安全维度
\item 公理 \ref{ax:learning}（学习）不能推导公理 \ref{ax:retrieval}（检索）——参数更新与信息检索是独立的机制
\end{itemize}

\subsection{可推导的核心定理}

\begin{theorem}[安全诊断定理 —— Safe Diagnosis Theorem]
在公理 \ref{ax:perception}--\ref{ax:consensus} 的假设下：
当 DMA 置信度低于 $\tau_{\text{conf}} = 0.62$ 或 TFA 风险等级为 critical 时，
HMAE 系统必然通过第1层（FORCE\_HITL）或第2层（DMA\_VETO）升级至人工干预，
从而避免在这些条件下出现过度自信型错误。
对于高置信度（$\geq 0.8$）且低风险的情况，系统可能发生错误，
但通过公理 \ref{ax:learning} 的持续学习和 Platt 校准机制逐步改善。
\end{theorem}

\textbf{证明概要：} 共识智能体的第1层 (FORCE\_HITL) 是安全网——
任何高危信号（DMA 置信度 $< 0.62$、验证严重冲突、TFA 风险等级 critical、TFA-诊断冲突）
直接触发人工干预，不可被后续层级覆盖。第2层 DMA\_VETO 是第二道防线——
DMA 置信度不足时一票否决。两层防御保证在低置信度或高风险的条件下假阴性率被严格控制。
对于高置信度低风险的情况，系统的自主决策受 Platt 校准和持续学习的统计规律支配，
存在不可消除的错误概率，但此概率随学习步数增加而递减。$\square$

\begin{theorem}[信息获取的边际递减定理 —— Diminishing Returns of Retrieval]
在公理 \ref{ax:retrieval} 的假设下，每增加一个检索文档的边际信息增益递减：
\begin{equation}
\frac{\partial}{\partial k} H(Y \mid q, \mathcal{K}_k) < 0, \quad
\frac{\partial^2}{\partial k^2} H(Y \mid q, \mathcal{K}_k) > 0
\end{equation}
即条件熵随检索文档数 $k$ 递减，但递减速率越来越慢（凸函数）。
\end{theorem}

\textbf{证明概要：} 检索文档按相似度降序排列。
第 $k$ 个文档的相关性 $s_k$ 随 $k$ 递减（至少不增），其对条件熵的贡献正比于 $s_k$。
因此 $\partial H / \partial k \propto -s_k$（负值递减）且 $\partial^2 H / \partial k^2 \propto s_k - s_{k+1} \geq 0$。$\square$

\begin{theorem}[容错完备性定理 —— Fault-Tolerance Completeness]
在公理 \ref{ax:degradation} 的假设下，HMAE 在任何单点故障（任意一个智能体失效）下保持功能完整；
在任何双点故障下保持降级功能；仅在三点及以上故障时可能完全不可用。
\end{theorem}

\textbf{证明概要：} 退化等级的 $L=0,1,2,3$ 定义中，$L=2$ (PARTIAL) 对应 unavailable $\geq 1$——
即一个组件不可用时系统仍在 PARTIAL 级别（功能降级但可用）。$L=3$ (SEVERE) 对应 unavailable $\geq 3$——
需要三个组件同时不可用才进入严重退化。因此双点故障时系统最多处于 PARTIAL 级别。$\square$

% ------------------------------------------------------------
\section{公理体系在非医疗领域的可迁移性}

HMAE 的公理体系虽然源自医疗 AI，但其抽象层次足够高，可以直接迁移到其他高风险决策领域：

\begin{table}[H]
\centering
\small
\begin{tabular}{lll}
\toprule
\textbf{公理} & \textbf{医疗领域体现} & \textbf{可迁移领域} \\
\midrule
感知完备性 (公理1) & 文本+影像融合 & 自动驾驶（摄像头+激光雷达） \\
检索不确定性 (公理2) & 医学文献检索 & 法律判例检索、金融研报检索 \\
低温收敛 (公理3) & 诊断决策 & 信贷审批、欺诈检测 \\
时序因果性 (公理4) & 病情风险预测 & 设备剩余寿命预测、市场风险预测 \\
共识收敛 (公理5) & 多专家会诊 & 多方风控评审、AI 伦理委员会 \\
隐私均衡 (公理6) & 患者数据保护 & 金融交易隐私、企业机密保护 \\
退化安全 (公理7) & 系统容错 & 航空飞控、核电站控制 \\
持续学习 (公理8) & 在线校准 & 推荐系统、自适应控制 \\
\bottomrule
\end{tabular}
\caption{公理体系的跨领域可迁移性}
\end{table}

\section{公理体系的方法论意义：从"工程直觉"到"形式化保证"}

HMAE 公理体系的建立标志着系统设计从\textbf{经验工程阶段}进入\textbf{形式化工程阶段}：

\begin{enumerate}[nosep]
\item \textbf{之前（经验工程）}：系统设计依赖工程师的直觉和经验，"这个阈值设为 0.62 是因为效果好"、"加噪声是因为要保护隐私"
\item \textbf{之后（形式化工程）}：每个设计选择都有公理级的形式化理由。"$\tau_{\text{conf}} = 0.62$ 是公理3（低温收敛）和公理5（共识收敛）的交叉约束下的最优解"、"$\sigma_{\text{noise}} = 0.015$ 是公理6（隐私均衡）在 $\epsilon$ 预算约束下的解"
\end{enumerate}

这一转变的核心价值不在于"找到了对的参数"（参数可能随数据变化），
而在于\textbf{建立了参数选择的逻辑链条}——任何参数的调整都可以追溯到一个或多个公理，
从而可以进行形式化的敏感性分析和安全性验证。

% ============================================================
% APPENDIX
% ============================================================
\part*{附录}
\addcontentsline{toc}{part}{附录}

\section*{A. 关键常数表}

\begin{longtable}{lll}
\toprule
\textbf{符号} & \textbf{值} & \textbf{含义} \\
\midrule
$d$ & 128 & 嵌入维度 \\
$k_{\max}$ & 2400 & TF-IDF 最大特征数 \\
$\tau_{\text{conf}}$ (confidence\_threshold) & 0.62 & DMA 置信度阈值 \\
$\tau_{\text{ret}}$ (retrieval\_relevance\_threshold) & 0.24 & 检索相关性阈值 \\
$\lambda$ (cost\_lambda) & 0.72 & 精度-成本权衡 \\
$N_{\max}$ (nash\_max\_rounds) & 6 & Nash 均衡最大迭代轮数 \\
$T$ (dma\_softmax\_temperature) & 0.18 & DMA Softmax 温度 \\
$T_g$ (retrieval\_query\_softmax\_temperature) & 0.35 & RAA 门控 Softmax 温度 \\
$\tau_{\text{dup}}$ (fusion\_duplicate\_hamming\_threshold) & 4 & 文档去重汉明距离 \\
$\tau_{\text{conflict}}$ (fusion\_conflict\_hamming\_threshold) & 26 & 冲突检测汉明距离 \\
$\tau_U$ (query\_uncertainty\_threshold) & 0.58 & 查询不确定性阈值 \\
$\sigma_{\text{noise}}$ (secure\_gradient\_noise\_std) & 0.015 & 差分隐私噪声标准差 \\
$\eta$ (lora\_learning\_rate) & 0.08 & 文档权重学习率 \\
$\eta_g$ (gradient\_margin\_shift\_rate) & 0.02 & 检索裕度偏移率 \\
$\tau_{\text{risk}}$ (high\_risk\_tfa\_threshold) & 0.65 & TFA 高风险阈值 \\
$\tau_{\text{timeout}}$ (agent\_timeout\_seconds) & 120s & 智能体超时 \\
$\tau_{\text{ver\_age}}$ (verification\_max\_age\_days) & 365 & 文档时效上限 \\
$w_{\text{DMA}}$ & 2 & DMA 共识投票权重 \\
$w_{\text{TFA}}$ & 1 & TFA 共识投票权重 \\
$w_{\text{Ver}}$ & 1 & 验证投票权重 \\
$\text{pass\_threshold}$ & 3 & 共识通过所需总权重 \\
$n_{\min}$ (continuous\_learning\_min\_feedback) & 10 & 持续学习触发反馈数 \\
\bottomrule
\caption{系统关键常数}
\end{longtable}

\section*{B. 模块映射表}

\begin{longtable}{lll}
\toprule
\textbf{章节} & \textbf{模块/类名} & \textbf{核心数学操作} \\
\midrule
\S1 & MultimodalPerceptionAgent & 交叉注意力融合 (式 \ref{eq:fusion}) \\
\S2 & RetrievalAugmentedAgent & TF-IDF + FAISS + 不确定性量化 \\
\S3 & DecisionMakingAgent & 多源得分融合 + Platt 校准 + Llama3 \\
\S4 & TemporalForeseeingAgent & TCN/EMA/Trans + 门控融合 + RevIN + Reprogramming \\
\S5 & KnowledgeFusionAgent & SimHash + 汉明去重 + 冲突检测 \\
\S6 & KnowledgeVerificationAgent & 多源事实评分 + 时效性检查 \\
\S7 & KnowledgeDistillationEngine & Fr\'{e}chet 均值 (加权球面平均) \\
\S8 & HumanInTheLoopManager & 干预判定逻辑 + 三型错误分类 \\
\S9 & PrivacySecurityAgent & 高斯差分隐私 + XOR 流加密 \\
\S10 & ConsensusModule & 三层级投票矩阵 \\
\S11 & 博弈论检索 & 纳什均衡最佳响应迭代 \\
\S12 & OnlineKnowledgeBaseManager & 增量 KB 更新 + FAISS 重建 \\
\S13 & MedTsLLMAdapter & Cross-Attention Reprogramming + BioBERT \\
\bottomrule
\end{longtable}

\section*{C. 公理-定理-参数依赖图}

\begin{table}[H]
\centering
\small
\begin{tabular}{lll}
\toprule
\textbf{公理} & \textbf{可推导定理} & \textbf{约束的关键参数} \\
\midrule
公理1 (感知完备性) & 安全诊断定理 & $d=128$, $\mathbf{W}_{\text{fusion}}$ 初始化 \\
公理2 (检索不确定性) & 边际递减定理 & $\lambda=0.72$, $\tau_U=0.58$, $k_{\max}=2400$ \\
公理3 (低温收敛) & 安全诊断定理 & $T=0.18$, Platt $(A,B)$ \\
公理4 (时序因果性) & 安全诊断定理 & $T=72$, $\alpha=0.18$, 三窗口权重 \\
公理5 (共识收敛) & 安全诊断定理 + 容错完备性 & $\tau_{\text{conf}}=0.62$, $N_{\max}=6$, 权重 $[2,1,1]$ \\
公理6 (隐私均衡) & — & $\sigma_{\text{noise}}=0.015$ \\
公理7 (退化安全) & 容错完备性定理 & $\tau_{\text{timeout}}=120s$, $L \in \{0,1,2,3\}$ \\
公理8 (期望损失递减) & — & $\eta=0.08$, $n_{\min}=10$ \\
\bottomrule
\end{tabular}
\caption{公理-定理-参数依赖关系}
\end{table}

\end{document}
